% Build: make paper (figures from scripts/paper_figures.py, then tectonic -X compile main.tex from paper/). Tectonic 0.15.0 prints
% "TeX rerun seems needed, but stopping at 6 passes" for any document with a BibTeX bibliography (its bookkeeping of the .bbl,
% reproduced on a one-entry article); the aux, bbl and out files are identical between runs, so the output is complete.
% For submission switch the class options to [sigconf,review,anonymous].
% Source comments: a line beginning "% src:" names the report JSON path, table and key behind the sentence that follows it.
% $RUN6 = paper/data/run8-randomization-window (the final run: run 8, the pipeline of run 7 on a ledger with the date rule of the archive's June 2024 randomization notice, docs/method.md call 275, from one code state on commit 5e4fedf, with run 7's configuration as the ablation include-randomization-window; the macro keeps its name because scripts/paper_audit.py resolves it; runs 5, 6 and 7 are kept beside it under paper/data); the data root's reports/ hold the same files.
% The planted-ledger backtests with every rung (positive_control_rungs) sit under $RUN6/positive-control/rungs/{n2000,n11100}; the held-out gain of the positive control is docs/method.md call 230 and the rung comparison call 231.
% No marker remains: the extras rerun of round 02 (issued_row_brier_skill, guard-pairs at 500 resamples) is in the run folder, and every source line names a file under $RUN6 (a few cite code or documentation).
\documentclass[sigconf,nonacm]{acmart}
\usepackage{booktabs}
\usepackage{graphicx}
\usepackage{amsmath}
\usepackage{xcolor}
\newcommand{\pwin}{p_{\mathrm{win}}}
\newcommand{\psig}{p_{\mathrm{sig}}}
\newcommand{\todo}[1]{{\color{red}[#1]}}
\setlength{\emergencystretch}{2em}

\begin{document}

\title{Can Past Experiments Predict the Next One? A Public Benchmark for Pre-Launch Forecasts of A/B Tests}

\author{Yingjie Li}
\affiliation{\institution{Independent researcher}\city{Seattle}\country{USA}}
\email{liyingjie@journey-ai.net}

\begin{abstract}
\textbf{Problem.} Before launching a change, a team wants to know what its experiment is likely to show: what result to expect, whether the experiment is worth running and how large to make it. Its past experiments may hold evidence on this. Industrial systems now screen candidates before a live test; the evidence they report rests on private data. We found no public benchmark that shows whether forecasts from past experiments can be trusted, or when those experiments say too little.

% src: $RUN6/backtest-upworthy-20261106-2dc461f52ea4.json header.tests_scored 2949; $RUN6/backtest-asos-20261103-6bb5a15c5ff0.json header.tests_scored 68
\textbf{Method.} We simulate how a team would have forecast each test at the time, using two public archives of concluded A/B tests: 2,949 scored headline and image tests from an online publisher (Upworthy), after leaving out a period its maintainers flag for randomization problems and rows that fail our admission rules, and 68 scored experiments from a fashion retailer (ASOS). We fit models only on earlier tests and predict each new test before it starts. The retailer's archive has no dates, so its order is arbitrary and tests only the mechanics. The archives record the traffic each test received, not what was planned, so forecasts use a figure a real forecast would not have; we also report a run with it fixed. A synthetic archive with a planted signal checks whether the backtest can detect one. All results are exploratory: we reused the same data and changed the method after seeing results, so they need testing on fresh data, such as the publisher's unopened confirmatory split.

\textbf{Findings.}
\setlength{\leftmargini}{1em}
\begin{itemize}
% src: $RUN6/backtest-upworthy-20261106-2dc461f52ea4.json tables.aggregate[cell=ALL,rung=P0].logscore_vs_p1 -0.11966 (P1 above P0 by 0.120 nats); tables.aggregate[cell=ALL,rung=P1].brier_skill 0.00228 (-0.00116, 0.00612), brier_skill_planned 0.00464 (0.00077, 0.00827), ece_percentile 100; $RUN6/paper-extras-upworthy-20261106-2dc461f52ea4.json p1_minus_p0_by_origin.by_origin_clear_of_zero 7, exposure_fixed.auc_event_within_ties_exact main.P1 0.525, fixed.P1 0.500 (tables/fmt-exposure-fixed.tex: 21 of 21 strata hold one forecast), calibration_within_period.rungs.P1.by_stratum_group.ALL.share_of_ece_removed 0.70; 90 percent bootstrap intervals, Section 5; $RUN6/paper-extras-asos-20261103-6bb5a15c5ff0.json gates_by_stratum.strata.ALL.verdict "undecidable at this n"
\item \emph{Can past experiments predict new ones? Partly.} For the publisher, a prior fitted per group of tests forecasts how results will be distributed better than the base rate does (0.120 nats per comparison on the log score, in every quarter scored). It describes the spread of results, not which test will succeed: within a group and period, its small ability to rank tests comes entirely from the traffic figure. Its skill at forecasting a significant positive result, measured against the base rate, is at most small: the 90 percent interval spans zero on our main definition and lies just above zero, unadjusted, when significance is set at the precision expected from traffic. Its probabilities are miscalibrated, mostly in their overall level for each period and group. The retailer's tests are too few to decide.
% src: $RUN6/paper-extras-synthetic-20261103-bf4fd286e2b2.json positive_control_rungs.runs[n_tests=11100].by_cell ALL rung=P2 recovered_fraction_logscore 0.0919, headline_only rung=P2 recovered_fraction_logscore 0.0319 (tables/positive-control-rungs.tex); $RUN6/backtest-upworthy-20261106-2dc461f52ea4.json tables.decision_value origin=ALL budget=0.25 policy=p_win rung=text_lr diff_matched 0.0457; tables/fmt-ranking-paired.tex (P2 below the arm-count order at every budget; P2 minus least precise first spans zero or lies below it); tables/fmt-abstention-rows.tex (issued above abstained on log score; not on Brier skill)
\item \emph{How?} Weighting similar past tests matches the prior on log score and skill, and ranks tests better than the prior but no better than rules reading only a test's number of arms or traffic. The backtest is weak at detecting such a signal: a planted signal returns 0.09 of the gain of a model told the signal, and 0.03 in a group with about as many comparisons as the publisher's largest. A text classifier ranks best among the models, its top quarter of tests capturing 4.6 points more of the positive lift than a random choice matched on number of arms, but it fails that control. Issued forecasts score above withheld ones on log score, not on skill, and the two sets differ in period and group, so abstaining is not shown to help.
% src: $RUN6/ablations/backtest-upworthy-20261106-2dc461f52ea4--prior-normal.json tables.aggregate[cell=ALL,rung=P1] logscore -0.8316 against the main run's -0.7566 (0.075 nats); $RUN6/backtest-upworthy-20261106-2dc461f52ea4.json tables.aggregate[cell=ALL,rung=P0].logscore_vs_p1 -0.11966 (0.120 nats); tables/fmt-prior-settings.tex and tables/fmt-p2-settings.tex (other settings, forced weight); $RUN6/paper-extras-upworthy-20261106-2dc461f52ea4.json prior_family_scale (tables/fmt-prior-family.tex), sizing_evaluation.test_and_roll (tables/fmt-sizing.tex: prior minus base-rate at multiple 20, adjusted interval spans zero)
\item \emph{What improves prediction?} A heavy-tailed forecast shape gained most, 0.075 of the 0.120 nats, much of it from pooling tests at different click-through levels, so we do not claim that effects themselves are heavy-tailed. Other settings barely moved the score, and forcing weight onto similar past tests made it worse. In a simulation with no money involved, sizing tests from the forecast shows no measurable advantage over the base rate.
\end{itemize}

\textbf{Tools.} Forecasts come from a base rate (how often past tests in a group came out each way), a prior fitted to past tests, that prior moved toward similar past tests, and three learned models: a boosted model, a local comparison set and a text classifier.

\textbf{Release.} A benchmark on those public archives, which adds no new data: a fixed protocol, data format, baselines and reports. It comes with Moot Court, an MIT-licensed Python package that regenerates the results tables and figures from its reports, reruns those reports from the public archives, and runs on a company's private aggregates. It has not been deployed in production.
\end{abstract}

\maketitle

\section{Introduction}
\label{sec:intro}
Before launching a change, a team wants to know what its experiment is likely to show. This paper is about that forecast: the result to expect and whether it will be significant (Section~\ref{sec:ladder-results}), with choosing tests (Section~\ref{sec:decision-results}) and sizing them (Section~\ref{sec:sizing-results}) as uses of the forecast, each evaluated here. We do not measure what happens after a change is deployed. A team that runs online controlled experiments accumulates a \emph{ledger}, one record for each concluded test. Published methods mostly use such a ledger after a test, to shrink a noisy estimate, though it could also inform the forecast before one. Three 2026 industrial systems now screen candidates before a live test; the evidence they report rests on private data (Section~\ref{sec:related}). We found no public benchmark that does three things: fixes a protocol under which a forecast sees only rows from before its test, with one stated exception for traffic (Section~\ref{sec:setting}); scores forecasts so that one that looks calibrated but cannot tell tests apart is seen for what it is; and uses gates, so that a reader knows how many tests a claim needs. This paper releases such a benchmark and applies it to two public archives and a synthetic one.

The system issues a \emph{card} before a test launches: a frozen forecast of the test's result, with the probability of a significant positive result and the probability that the variant wins. The forecasts come from a \emph{ladder} of models, each step or \emph{rung} building on the one below. P0 is the base rate. P1, the prior, is a distribution of effects estimated from concluded tests, allowing for each test's noise (an empirical-Bayes prior). P2, the precedent rung, moves that prior toward similar past tests when a \emph{guard}, a check on earlier tests held back for validation, admits the move. Three learned models, among them a text classifier, run beside the ladder as challengers. Models are fit per \emph{cell}, a group of comparisons of the same kind (on the publisher's archive the kind of change, on the retailer's the metric), with partial pooling across cells. A \emph{gate} is the number of tests at which a simulated forecaster of a stated skill would clear its interval with high power: below the gate a claim at that skill is undecidable, and above it the interval still has to be read. A card can also \emph{abstain}, naming why, and then shows the base rate instead of its forecast.

We ask three questions, and the discussion answers them in order (Section~\ref{sec:discussion}).
\begin{enumerate}
\item[RQ1] Can a team predict a new experiment's result from its concluded experiments?
\item[RQ2] If so, how, and how well?
\item[RQ3] What improves the prediction?
\end{enumerate}
Every result in this paper is exploratory. We read the same data repeatedly, one split of the publisher's archive (the \emph{headline archive} below) and the whole of the retailer's, and changed the method between runs after seeing results. The headline archive's confirmatory split is unopened, and the archive asks that analyses of it be registered first, so it remains available to test these results. Separately, because the archive's maintainers report randomization problems in one period, the main analysis leaves out the tests that started then, a rule fixed before any number of this run existed (Section~\ref{sec:setting}).

% src: $RUN6/paper-extras-synthetic-20261103-bf4fd286e2b2.json positive_control_rungs.runs[n_tests=11100].planted_text_effect (a lift of one percent), by_cell ALL rung=P2 recovered_fraction_logscore 0.0919; tables/fmt-exposure-fixed.tex; tables/fmt-ranking-paired.tex; tables/fmt-control-rungs.tex; tables/event-definition-upworthy.tex
In short, a prior fitted to the ledger forecasts how the next result will be distributed better than a base rate does, but it describes the spread of results and not which test will succeed. Its skill on which test will be significant is at most small and depends on how that event is defined: the interval spans zero at the test's realized precision and lies just above zero, unadjusted, at the precision expected from traffic or with traffic fixed. Among the models a text classifier ranks tests best, but it fails the planted-signal control, so that control neither validates nor explains the result. Precedents match the prior on score and rank tests better than it does, but no better than rules reading only the number of arms or traffic. This backtest recovers little of a planted lift of one percent, so it cannot rule out a precedent signal of that size. Table~\ref{tab:glance} sets the numbers side by side.

\begin{table*}[t]\centering\footnotesize
% src: $RUN6/backtest-upworthy-20261106-2dc461f52ea4.json header.bootstrap_resamples 500; $RUN6/backtest-asos-20261103-6bb5a15c5ff0.json header.tests_in_ledger 78, header.tests_scored 68; tables/coverage-ratio-asos.tex tables/coverage-ratio-synthetic.tex ratio range 1.00 to 1.00; tables/glance.tex
\caption{Results at a glance: the main numbers of the two public archives and the synthetic archive side by side, for the prior (P1) unless the row names another model. Intervals are 90 percent bootstrap by test. Section~\ref{sec:protocol} defines each measure; the results sections read each row. Sources: Tables~\ref{tab:main-upworthy}, \ref{tab:main-asos}, \ref{tab:main-synthetic}, \ref{tab:coverage-ratio}, \ref{tab:decision-intervals-upworthy}, \ref{tab:decision-intervals-asos}, \ref{tab:decision-intervals-synthetic}, \ref{tab:event-definition} and~\ref{tab:fmt-gates-strata}.}
\label{tab:glance}
% generated by paper/build_tables.py from the report keys named above each row (front)
\begin{tabular}{lrrr}
\toprule
 & Upworthy (headline archive) & ASOS (deployment-shaped) & Synthetic (no planted text signal) \\
\midrule
% src: $RUN6/backtest-upworthy-20261103-4bfb354b87e3.json header.tests_scored; $RUN6/backtest-upworthy-20261103-4bfb354b87e3.json header.scored_rows; $RUN6/backtest-asos-20261103-6bb5a15c5ff0.json header.tests_scored; $RUN6/backtest-asos-20261103-6bb5a15c5ff0.json header.scored_rows; $RUN6/backtest-synthetic-20261103-bf4fd286e2b2.json header.tests_scored; $RUN6/backtest-synthetic-20261103-bf4fd286e2b2.json header.scored_rows
scored tests / rows & 3,797 / 13,247 & 68 / 337 & 907 / 1,814 \\
% src: $RUN6/backtest-upworthy-20261103-4bfb354b87e3.json tables.aggregate[cell=ALL,rung=P0].logscore_vs_p1; $RUN6/backtest-upworthy-20261103-4bfb354b87e3.json tables.aggregate[cell=ALL,rung=P0].logscore_vs_p1_lo; $RUN6/backtest-upworthy-20261103-4bfb354b87e3.json tables.aggregate[cell=ALL,rung=P0].logscore_vs_p1_hi; $RUN6/backtest-asos-20261103-6bb5a15c5ff0.json tables.aggregate[cell=ALL,rung=P0].logscore_vs_p1; $RUN6/backtest-asos-20261103-6bb5a15c5ff0.json tables.aggregate[cell=ALL,rung=P0].logscore_vs_p1_lo; $RUN6/backtest-asos-20261103-6bb5a15c5ff0.json tables.aggregate[cell=ALL,rung=P0].logscore_vs_p1_hi; $RUN6/backtest-synthetic-20261103-bf4fd286e2b2.json tables.aggregate[cell=ALL,rung=P0].logscore_vs_p1; $RUN6/backtest-synthetic-20261103-bf4fd286e2b2.json tables.aggregate[cell=ALL,rung=P0].logscore_vs_p1_lo; $RUN6/backtest-synthetic-20261103-bf4fd286e2b2.json tables.aggregate[cell=ALL,rung=P0].logscore_vs_p1_hi
P1 log score gain over P0, nats (90\%) & 0.149 (0.136, 0.162) & 0.368 (0.192, 0.550) & 0.010 (0.001, 0.018) \\
% src: $RUN6/backtest-upworthy-20261103-4bfb354b87e3.json tables.aggregate[cell=ALL,rung=P1].brier_skill; $RUN6/backtest-upworthy-20261103-4bfb354b87e3.json tables.aggregate[cell=ALL,rung=P1].brier_skill_lo; $RUN6/backtest-upworthy-20261103-4bfb354b87e3.json tables.aggregate[cell=ALL,rung=P1].brier_skill_hi; $RUN6/backtest-asos-20261103-6bb5a15c5ff0.json tables.aggregate[cell=ALL,rung=P1].brier_skill; $RUN6/backtest-asos-20261103-6bb5a15c5ff0.json tables.aggregate[cell=ALL,rung=P1].brier_skill_lo; $RUN6/backtest-asos-20261103-6bb5a15c5ff0.json tables.aggregate[cell=ALL,rung=P1].brier_skill_hi; $RUN6/backtest-synthetic-20261103-bf4fd286e2b2.json tables.aggregate[cell=ALL,rung=P1].brier_skill; $RUN6/backtest-synthetic-20261103-bf4fd286e2b2.json tables.aggregate[cell=ALL,rung=P1].brier_skill_lo; $RUN6/backtest-synthetic-20261103-bf4fd286e2b2.json tables.aggregate[cell=ALL,rung=P1].brier_skill_hi
P1 Brier skill against the base rate (90\%) & 0.001 (-0.002, 0.005) & -0.107 (-0.237, -0.002) & 0.002 (-0.013, 0.017) \\
% src: $RUN6/gates.json by_base_rate["0.09"].by_target["0.90"].n_400; $RUN6/backtest-upworthy-20261103-4bfb354b87e3.json tables.aggregate[cell=ALL,rung=P1].gate_base_rate; $RUN6/backtest-upworthy-20261103-4bfb354b87e3.json tables.aggregate[cell=ALL,rung=P1].decidable_by_tests; $RUN6/gates.json by_base_rate["0.14"].by_target["0.90"].n_400; $RUN6/backtest-asos-20261103-6bb5a15c5ff0.json tables.aggregate[cell=ALL,rung=P1].gate_base_rate; $RUN6/backtest-asos-20261103-6bb5a15c5ff0.json tables.aggregate[cell=ALL,rung=P1].decidable_by_tests; $RUN6/gates.json by_base_rate["0.14"].by_target["0.90"].n_400; $RUN6/backtest-synthetic-20261103-bf4fd286e2b2.json tables.aggregate[cell=ALL,rung=P1].gate_base_rate; $RUN6/backtest-synthetic-20261103-bf4fd286e2b2.json tables.aggregate[cell=ALL,rung=P1].decidable_by_tests
release gate, tests at base rate; pooled claim & 250 at 0.09; decidable & 400 at 0.14; undecidable & 400 at 0.14; decidable \\
% src: $RUN6/backtest-upworthy-20261103-4bfb354b87e3.json tables.aggregate[cell=ALL,rung=P1].ece_percentile; $RUN6/backtest-asos-20261103-6bb5a15c5ff0.json tables.aggregate[cell=ALL,rung=P1].ece_percentile; $RUN6/backtest-synthetic-20261103-bf4fd286e2b2.json tables.aggregate[cell=ALL,rung=P1].ece_percentile
P1 ECE, percentile of the clustered null & 100 & 98.8 & 95.2 \\
% src: $RUN6/backtest-upworthy-20261103-4bfb354b87e3.json tables.aggregate[cell=ALL,rung=P1].coverage_80; $RUN6/backtest-upworthy-20261103-4bfb354b87e3.json tables.aggregate[cell=ALL,rung=P1].coverage_90; $RUN6/backtest-asos-20261103-6bb5a15c5ff0.json tables.aggregate[cell=ALL,rung=P1].coverage_80; $RUN6/backtest-asos-20261103-6bb5a15c5ff0.json tables.aggregate[cell=ALL,rung=P1].coverage_90; $RUN6/backtest-synthetic-20261103-bf4fd286e2b2.json tables.aggregate[cell=ALL,rung=P1].coverage_80; $RUN6/backtest-synthetic-20261103-bf4fd286e2b2.json tables.aggregate[cell=ALL,rung=P1].coverage_90
P1 coverage of the 80 / 90 percent intervals & 0.81 / 0.89 & 0.78 / 0.86 & 0.78 / 0.89 \\
% src: $RUN6/backtest-upworthy-20261103-4bfb354b87e3.json tables.coverage_by_se_ratio[rung=P1,decile=1].coverage_90; $RUN6/backtest-upworthy-20261103-4bfb354b87e3.json tables.coverage_by_se_ratio[rung=P1,decile=10].coverage_90
P1 90 percent coverage, most to least precise decile & 1.00 to 0.66 & n/a & n/a \\
% src: $RUN6/backtest-upworthy-20261103-4bfb354b87e3.json tables.aggregate[cell=ALL,rung=P1].abstention_rate; $RUN6/backtest-asos-20261103-6bb5a15c5ff0.json tables.aggregate[cell=ALL,rung=P1].abstention_rate; $RUN6/backtest-synthetic-20261103-bf4fd286e2b2.json tables.aggregate[cell=ALL,rung=P1].abstention_rate
abstention rate & 0.141 & 0.457 & 0.182 \\
% src: $RUN6/backtest-upworthy-20261103-4bfb354b87e3.json tables.aggregate[cell=ALL,rung=P2].logscore_vs_p1; $RUN6/backtest-upworthy-20261103-4bfb354b87e3.json tables.aggregate[cell=ALL,rung=P2].logscore_vs_p1_lo; $RUN6/backtest-upworthy-20261103-4bfb354b87e3.json tables.aggregate[cell=ALL,rung=P2].logscore_vs_p1_hi; $RUN6/backtest-asos-20261103-6bb5a15c5ff0.json tables.aggregate[cell=ALL,rung=P2].logscore_vs_p1; $RUN6/backtest-synthetic-20261103-bf4fd286e2b2.json tables.aggregate[cell=ALL,rung=P2].logscore_vs_p1
P2 minus P1, paired log score (90\%) & 0.0004 (-0.0003, 0.0012) & 0.0000; guard holds $\lambda$ at zero & 0.0000; guard holds $\lambda$ at zero \\
% src: $RUN6/backtest-upworthy-20261103-4bfb354b87e3.json tables.decision_value[origin=ALL,budget=0.25,policy=p_win,rung=text_lr].diff_matched; $RUN6/backtest-upworthy-20261103-4bfb354b87e3.json tables.decision_value[origin=ALL,budget=0.25,policy=p_win,rung=text_lr].diff_matched_lo; $RUN6/backtest-upworthy-20261103-4bfb354b87e3.json tables.decision_value[origin=ALL,budget=0.25,policy=p_win,rung=text_lr].diff_matched_hi; $RUN6/backtest-asos-20261103-6bb5a15c5ff0.json tables.decision_value[origin=ALL,budget=0.25,policy=p_win,rung=P1].diff_matched; $RUN6/backtest-asos-20261103-6bb5a15c5ff0.json tables.decision_value[origin=ALL,budget=0.25,policy=p_win,rung=P1].diff_matched_lo; $RUN6/backtest-asos-20261103-6bb5a15c5ff0.json tables.decision_value[origin=ALL,budget=0.25,policy=p_win,rung=P1].diff_matched_hi; $RUN6/backtest-synthetic-20261103-bf4fd286e2b2.json tables.decision_value[origin=ALL,budget=0.25,policy=p_win,rung=P1].diff_matched; $RUN6/backtest-synthetic-20261103-bf4fd286e2b2.json tables.decision_value[origin=ALL,budget=0.25,policy=p_win,rung=P1].diff_matched_lo; $RUN6/backtest-synthetic-20261103-bf4fd286e2b2.json tables.decision_value[origin=ALL,budget=0.25,policy=p_win,rung=P1].diff_matched_hi
best decision value at a quarter budget (90\%) & text classifier 0.039 (0.022, 0.054) & all rungs 0.000 (0.000, 0.000) & prior 0.195 (0.173, 0.216) \\
% src: $RUN6/backtest-upworthy-20261103-4bfb354b87e3.json tables.decision_value[origin=ALL,budget=0.25,policy=p_win,rung=P1].diff_matched; $RUN6/backtest-upworthy-20261103-4bfb354b87e3.json tables.decision_value[origin=ALL,budget=0.25,policy=p_win,rung=P1].diff_matched_lo; $RUN6/backtest-upworthy-20261103-4bfb354b87e3.json tables.decision_value[origin=ALL,budget=0.25,policy=p_win,rung=P1].diff_matched_hi; $RUN6/backtest-asos-20261103-6bb5a15c5ff0.json tables.decision_value[origin=ALL,budget=0.25,policy=p_win,rung=P1].diff_matched; $RUN6/backtest-asos-20261103-6bb5a15c5ff0.json tables.decision_value[origin=ALL,budget=0.25,policy=p_win,rung=P1].diff_matched_lo; $RUN6/backtest-asos-20261103-6bb5a15c5ff0.json tables.decision_value[origin=ALL,budget=0.25,policy=p_win,rung=P1].diff_matched_hi; $RUN6/backtest-synthetic-20261103-bf4fd286e2b2.json tables.decision_value[origin=ALL,budget=0.25,policy=p_win,rung=P1].diff_matched; $RUN6/backtest-synthetic-20261103-bf4fd286e2b2.json tables.decision_value[origin=ALL,budget=0.25,policy=p_win,rung=P1].diff_matched_lo; $RUN6/backtest-synthetic-20261103-bf4fd286e2b2.json tables.decision_value[origin=ALL,budget=0.25,policy=p_win,rung=P1].diff_matched_hi
P1 decision value at a quarter budget (90\%) & -0.004 (-0.012, 0.005) & 0.000 (0.000, 0.000) & 0.195 (0.173, 0.216) \\
\bottomrule
\end{tabular}

\par\smallskip\parbox{0.92\textwidth}{\scriptsize\raggedright
\textbf{Reading the table.}
\begin{itemize}\setlength{\itemsep}{0pt}
\item Log scores sit on each archive's own effect scale, so read that row within a column.
\item The two Brier skill rows differ in how the significant-positive event is defined: at the test's realized standard error, the main reading, or at its planned one. On ASOS and synthetic the two events coincide.
\item On ASOS the gate is not met, so no interval there decides a claim.
\item The gate row is for a skill of 0.10. The skills measured on the publisher's archive are far smaller, so a verdict of decidable says the interval is worth reading, not that a skill of this size can be told from zero.
% src: tables/glance.tex, row P1 ECE; $RUN6/backtest-upworthy-20261106-2dc461f52ea4.json tables.aggregate[cell=ALL,rung=P1].ece_percentile 100; $RUN6/backtest-synthetic-20261103-bf4fd286e2b2.json tables.aggregate[cell=ALL,rung=P1].ece_percentile 95.2
\item A percentile above 95 rejects calibration. The synthetic archive's 95.2 is just past that line; the publisher's 100 is above every draw of the null.
\item The gates of the headline and ASOS columns are simulated at that archive's own dependence (Table~\ref{tab:fmt-gates-strata}). No such simulation ran for the synthetic archive, so its gate is the report's scaling of the nearest simulated gate to its dependence.
\item The abstention rate is of all scored rows, the burn-in origin included.
% src: src/mootcourt/decision_value.py _matched_value (the chosen set's own value when the score fixes the arm stratum); $RUN6/backtest-asos-20261103-6bb5a15c5ff0.json tables.decision_value diff_matched 0 for every rung
\item ``Best'' decision value is over the models the archive runs. On ASOS the features are metric and number of variants, so each test's score is a function of its number of arms and the difference from the order matched on arms is zero by construction. No text classifier runs there because the archive has no text.
\item On the synthetic archive the prior's value comes from cells that differ in their means, not from text; P0, P1 and P2 tie.
\item The decile row is n/a on ASOS and synthetic because the ratio of realized to planned standard error is 1.00 throughout.
\item The positive control runs on its own generator: its recovered fractions are in Table~\ref{tab:positive-control-gain}.
\end{itemize}}
\end{table*}

The paper makes four contributions.
\begin{enumerate}
\item A protocol for pre-launch forecasts of experiment outcomes (Sections~\ref{sec:setting} and~\ref{sec:protocol}). A forecast sees only rows from tests past their assigned end. Its standard error is meant to come from planned exposure; no archive records a plan, so realized exposure stands in, and on the headline archive we measure the main scores with that figure fixed (Table~\ref{tab:fmt-exposure-fixed}). Only listed pre-launch features can be read, and an audit checks for leakage after each run. Forecasts are judged by proper scores shown beside a bound on the best attainable value, by calibration tests that allow for the comparisons of one test moving together, and by decision value against a random order matched on the number of arms. On the two public archives, one gate rule is simulated at each group's own within-test correlation.
% src: $RUN6/ablations/backtest-upworthy-20261106-2dc461f52ea4--effect-scale-relative+s-star-realized.json tables.aggregate[cell=ALL,rung=P1].brier_skill -0.1036 (tables/fmt-cautionary.tex); $RUN6/paper-extras-upworthy-20261106-2dc461f52ea4.json effect_scale_choices (tables/fmt-leak-choices.tex)
\item A methodological finding from our own first run: giving a pre-launch forecaster a test's realized precision, the standard error the finished test reported, puts in-test information into its forecast. Unlike most leakage, this lowers the score. Rerunning that setup on the final code, with relative lift on the headline archive, gives a Brier skill of minus 0.104, and forecasts and outcomes move in opposite directions across levels of precision, which only tables by precision reveal. Changing one choice at a time shows that realized precision is what pushes ranking below chance; relative lift alone does not. Our main analysis removes the realized click-through rate from the forecast's scale but still uses realized traffic, as contribution 1 states (Section~\ref{sec:cautionary}).
% src: $RUN6/backtest-upworthy-20261106-2dc461f52ea4.json tables.aggregate[cell=ALL,rung=P2].logscore_vs_p1 0.0001 (-0.0004, 0.0006), rung=P2 brier_skill 0.0032 against rung=P1 0.0023; $RUN6/paper-extras-upworthy-20261106-2dc461f52ea4.json ranking_paired (tables/fmt-ranking-paired.tex); $RUN6/paper-extras-synthetic-20261103-bf4fd286e2b2.json positive_control_rungs.runs[n_tests=11100].by_cell ALL rung=P2 recovered_fraction_logscore 0.0919, headline_only rung=P2 recovered_fraction_logscore 0.0319 (tables/positive-control-rungs.tex), by_cell ALL oracle tests 9883, headline_only oracle rows 6354, tests 3177; $RUN6/backtest-upworthy-20261106-2dc461f52ea4.json tables.aggregate[cell=headline_only,rung=P1] rows 5966, tables/fmt-p1-cells.tex tests 1,770
\item A baseline suite a new method has to beat (Section~\ref{sec:method}): a base rate, a fitted prior, precedents behind an admission guard and three learned challengers, with four abstention triggers and a frozen card that traces each forecast to its data and model. Its first results are summarized above. They changed the design in one way: we built abstention as a filter on quality, and because issued rows score above abstained rows on the log score and not on Brier skill, on rows that differ in period and group, it now stays only as a label of evidence (Section~\ref{sec:abstention-results}). The planted-signal control recovers 0.09 of an oracle's log-score gain pooled and 0.03 in a cell with about as many rows as the real headline-only cell but nearly twice its tests, so the backtest does not rule out a signal of that size (Section~\ref{sec:synthetic-results}). Sizing from the forecast is reported in Section~\ref{sec:sizing-results} and Table~\ref{tab:fmt-sizing}.
% src: $RUN6/backtest-asos-20261103-6bb5a15c5ff0.json header.tests_scored 68; $RUN6/paper-extras-asos-20261103-6bb5a15c5ff0.json gates_by_stratum.strata.ALL.own_gate.n_400 275 (tables/glance.tex)
\item An open-source package, Moot Court, that implements the benchmark (Section~\ref{sec:repro}). It regenerates every results table and figure from the report files with one command each, and reruns the reports from the public archives. It runs on an organization's own ledger of aggregates, which never leaves the organization. It has not been deployed in production. A team should run its own backtest first: that says whether forecasts can be trusted on its ledger and how many tests it takes to tell, and it does not judge any single change. On the retailer's archive the backtest returns \emph{undecidable}, as its size made certain (68 scored tests against a gate of 275 for a skill of 0.10), so this paper promises no improvement on product experiments.
\end{enumerate}
One boundary holds throughout. A forecast is advice before launch: it may help order and size tests, though the evidence that it orders them well is weak and sizing shows no advantage over a base rate. The live test still decides. Moot Court never changes the decision to launch, the traffic split or when a test stops.

Section~\ref{sec:hypotheses} states six hypotheses, written after two earlier runs on the same split, and Section~\ref{sec:verdicts} reads them back part by part.
\begin{itemize}
% src: paper/plan.md line 30 (H4: abstention rate below 30 percent after burn-in); $RUN6/backtest-upworthy-20261106-2dc461f52ea4.json tables.abstention[cell=ALL].rate (tables/verdicts.tex, H4 part 3)
\item \textbf{Held:} the prior beats the base rate on log score at every origin after the first (burn-in); abstention stays under its limit of 30 percent after burn-in; the text classifier captures more lift than the matched random order, though it fails the planted control; the retailer's archive stays undecidable, as its size made certain.
\item \textbf{Failed:} calibration of issued forecasts; issued rows scoring above all rows on skill; the prior ordering tests better than the matched random order.
\item \textbf{Not shown:} pooled skill on which test will be significant; the other models' lift over the matched random order, though each beats the plain random order the plan named; precedents beating the prior on tests that vary only the headline, a gain the planted control says this backtest could not have seen.
\end{itemize}

\section{Related Work}
\label{sec:related}

\paragraph{Empirical Bayes over many experiments.}
Shrinking parallel estimates toward a common prior \cite{efron1975stein}, random-effects meta-analysis \cite{dersimonian1986meta} and partial pooling \cite{gelman2007multilevel} underlie the rungs above the base rate. In online experimentation the prior has been learned from thousands of historical tests \cite{deng2015objective}, used to separate true effects from lucky draws \cite{azevedo2019empirical}, argued to be heavy-tailed \cite{azevedo2020fat}, fit nonparametrically \cite{guo2020empirical}, shown to improve launch decisions \cite{coey2019improving}, and combined with an experiment-specific shrinkage factor against the winner's curse \cite{mudd2026breaking}. Selecting on significance distorts estimates \cite{deng2021postselection}, and a historical prior loses calibration when the corpus is winner-selected or pooled across programs \cite{schultzberg2026bayesian}; our admission rules and per-cell priors answer both. The program's success rate as a prior was compiled across programs by Kohavi, Deng and Vermeer \cite{kohavi2022intuition}, following \cite{deng2015objective}; P0 makes it cell-specific and scores it, and scoring every rung against it follows the forecasting convention of skill against climatology, whose uncertainty term is the floor we report \cite{murphy1973vector}. These models take each row's standard error as given, and the in-test treatment of ratio metrics is the delta method \cite{deng2018delta}; we report a within-row coupling of a ratio metric with its own standard error that neither treats.

\paragraph{Forecasting before or during a test, and screening before it.}
Surrogate indices \cite{athey2019surrogate}, validated as a decision rule on 200 tests \cite{zhang2023evaluating}, always-valid inference \cite{johari2015always,johari2022always,howard2021timeuniform}, Bayesian optional stopping \cite{deng2016continuous} and learned variance reduction \cite{guo2021machine} sharpen a running test from its own data; our card is issued at time zero. Three 2026 industrial systems screen candidates before a live test: an archive-trained ranker that chooses which generated ad creatives enter a 50-arm test and reports offline recall and one live rank \cite{lee2026offline}, a deployed ranker that prioritizes variants from historical outcomes \cite{hu2026experimentation}, and a back-testing engine that replays historical campaign data through production allocation code \cite{mazzanti2026backtesting}; none scores an outcome distribution against live results, and none abstains. Winner predictors on the Upworthy archive are the closest text baselines: marginally better than chance at naming the winner of 17,681 headline tests \cite{ye2025lola}, about twice random on decisive tests after retraining on revealed outcomes \cite{clark2026predict}, and 0.689 pairwise accuracy in domain, chance out of it \cite{khimani2026finetuning}; our pairwise text classifier is that baseline on matched splits. Simulators that stand in for the audience report direction only \cite{li2026simgym,hut2026can}; LLM-generated outcomes recover human effects only after calibration against human data \cite{persson2026statistical,ankargren2026when,ashokkumar2026large}, persona-prompted models drift \cite{bisbee2024synthetic}, and respondent grounding buys fidelity on self-report rather than on live experiments \cite{park2024llm}; our second opinion therefore enters at weight zero. A trial-level surrogacy gate on 183 Upworthy interventions sends every one to a live test \cite{schultzberg2026when}, the bar our gates are the analogue of.

\paragraph{Human forecasts, win rates and decision value.}
Expert forecasts of treatment effects help when averaged \cite{dellavigna2018predicting}, should be recorded before results are known \cite{dellavigna2019predict}, and failed to pick the best of 54 interventions in a megastudy \cite{milkman2021megastudies}; P3 therefore freezes the forecast before the result is admitted. Published success rates run from about a third to about a tenth of ideas \cite{kohavi2022intuition,kohavi2009controlled,kohavi2013online,larsen2022statistical}, and 18 to 25 percent of significant results are estimated to be nulls \cite{berman2022false}, a computation that descends from the false-positive report probability of \cite{wacholder2004assessing} and the low-prior regime of \cite{ioannidis2005why}; P0 measures these rates per cell. The launch decision has been cast as a Bayes decision \cite{goldberg2017decision}, test-and-roll derives test sizes from the cost of exposure and rollout \cite{feit2019test}, single tests can be badly underpowered \cite{lewis2015unfavorable}, a program optimizing returns should test more ideas with smaller allocations \cite{sudijono2024optimizing}, and the caution that surprising results are usually wrong \cite{kohavi2020trustworthy} becomes explicit abstention here. The card's significance probability is an assurance in the sense of \cite{ohagan2005assurance}, a success probability averaged over the prior rather than a power at one assumed effect. The nearest decision-value evaluation scores launch rules by the returns they would have produced on 123 historical tests \cite{chou2025evaluating}; those rules consume each experiment's own data, ours only the ledger, and our oracle's value, the best positive observed lift among a test's arms, carries the winner's curse of \cite{lee2018winners} and is read as an inflated upper bound.

\paragraph{Scoring, calibration, abstention and the protocol.}
The logarithmic score, Brier score and CRPS are strictly proper \cite{gneiting2007strictly,matheson1976scoring}, forecasts should maximize sharpness subject to calibration \cite{gneiting2007probabilistic}, and the rolling origin is the forecasting literature's out-of-sample protocol \cite{tashman2000outofsample}, applied here to versioned ledger snapshots; leakage is used in the sense of \cite{kaufman2012leakage}, a failure of evaluation design, and the paired, clustered comparisons of log score behind the admission guard and the guarded-against-unguarded reading are tests of equal predictive accuracy in the line of \cite{diebold1995comparing,giacomini2006tests}, read here as sizing rules. Consistency bars for reliability diagrams \cite{brocker2007increasing} and the pitfalls of binned calibration estimators \cite{vaicenavicius2019evaluating} motivate the clustered null and the equal-mass ECE. Abstention descends from Chow's reject rule \cite{chow1970optimum} and selective classification \cite{geifman2017selective}; a 2026 multimodal judge of winners reaches a kappa of 0.14 on 330 real tests and only a vote-margin gate that abstains on half of them lifts it to 0.31 \cite{dooskin2026judge}; our triggers extend the gate to support, novelty, reliability and decision instability. Interval coverage is a conformal-style validity check \cite{vovk2005algorithmic,angelopoulos2021gentle} without exchangeability across cells, hence the novelty trigger.

\paragraph{The two archives, local pooling, and what appears new.}
The Upworthy Research Archive records every headline test from January 2013 to April 2015 in exploratory, confirmatory and hold-out splits \cite{matias2021upworthy}; prior analyses explain its effects after the fact \cite{banerjee2025language,robertson2023negativity,ottaviani2024testing}. The ASOS dataset releases daily checkpoints from 78 real experiments \cite{liu2021datasets,liu2024statistical}; its documented uses \cite{che2024optimization,kim2026odds} concern sequential testing and bandits. The closest prior art to P2 is local empirical Bayes \cite{li2025learning}; Stein-type shrinkage with an auxiliary estimate \cite{rosenman2020combining} and crossed random effects \cite{baayen2008mixed} are the estimator analogues. The claim that survives the survey is narrow: no prior system issues a pre-launch outcome distribution with a win probability for a single proposed experiment, abstains under named triggers, and scores it with proper scores, gates and a leakage audit against live results on public archives; the screening systems choose candidates by offline ranking or replay, the simulators report direction, and the estimators correct a completed test. The nearest neighbours \cite{li2025learning,ye2025lola,clark2026predict,khimani2026finetuning} are compared on matched splits.

\paragraph{Tooling landscape.}
\label{sec:landscape}
A fetch-verified survey of thirteen commercial experimentation platforms, seventeen open-source projects, sixteen internal platforms described in public papers and the 2023 to 2026 simulation and language-model tools found no product that issues a pre-launch outcome distribution or win probability for a proposed experiment, abstains, scores its forecasts against live results, or backtests over an archive; Appendix~\ref{app:landscape} and Table~\ref{tab:landscape} give the survey and its gap table.


\section{Problem Setting and Data}
\label{sec:setting}

\paragraph{The task.} Before a test launches, forecast what each of its comparisons will show, using only the results of tests that had already concluded.
\begin{itemize}
\item A \emph{contrast} is one treatment arm compared with the reference arm of one test, with an observed effect $y$ and a standard error $s$. Both are read at the test's \emph{assigned end}, a fixed date per test rather than the date the data stop.
\item A \emph{cell} groups contrasts by the kind of change.
\item An \emph{origin} $t$ is a point in time at which forecasts are made. A contrast is \emph{visible} at $t$ when it is admitted and its result was usable by then; on both public archives that is its assigned end. Every quantity computed at $t$ uses visible rows only.
\item The \emph{ledger} holds one row per concluded contrast (Table~\ref{tab:schema}, Appendix~\ref{app:tables}). \emph{Admission} is a status, not a filter: a row that fails a rule stays, with its reason. Snapshots of the ledger are immutable, each identified by a digest of its rows.
\item Three counts of tests are kept apart: tests \emph{in the ledger}; tests \emph{admitted}, with at least one admitted contrast; and tests \emph{scored}, with at least one contrast forecast by the precedents (P2), the main model. Admitted tests that are not scored started before the first origin with enough history (Table~\ref{tab:fmt-archives} gives both counts).
\end{itemize}

\paragraph{Upworthy Research Archive.} The archive records the headline and image tests the publisher ran \cite{matias2021upworthy}; Table~\ref{tab:fmt-archives} gives its period and counts. Each arm is a \emph{package}, one headline and image shown to a randomly assigned share of readers. We use only the public exploratory split, with impressions, clicks, headline text and an image identifier per package. The archive asks that analyses of its other splits be registered first, and we have not opened them. The reference arm is the earliest-created package, and every other package is one contrast. The effect $y$ is the difference in click-through rate in percentage points, with its binomial standard error; the relative lift is kept for display only (Section~\ref{sec:cautionary}). The cell is the kind of change: headline only, image only, or both.
% src: DATA.md, Upworthy statuses: low_exposure below 1000 impressions in either arm, low_clicks below 10 clicks in either arm, missing_text an empty headline, duplicate_arm same headline and image as the reference
Admission requires 1,000 impressions and 10 clicks in both arms, non-empty headlines and an arm that differs from the reference.
The clicks rule reads an outcome, so Section~\ref{sec:ablations} also scores the archive without it. Every contrast of a test shares one assigned end, seven days after the creation of its latest arm. That is our modelling choice, since the shared reference arm keeps collecting until the last arm ends. The pre-launch fields are the two headlines, the kind of change, the number of arms and the planned exposure per arm. The archive records no plan, so realized impressions stand in for it. That is the one post-launch quantity a forecast reads, and Table~\ref{tab:fmt-exposure-fixed} measures the main scores with it held fixed. All other post-launch fields sit in a block the feature builder cannot read, which the leakage audit checks after each run.

\paragraph{The randomization notice.}
% src: $RUN6/paper-extras-upworthy-20261106-2dc461f52ea4.json randomization_window.notice (June 2024; about 22 percent of tests; 25 June 2013 to 10 January 2014); DATA.md, The archive's 2024 randomization notice
The archive's maintainers report evidence of randomization problems in about 22 percent of tests, and discourage causal or confirmatory use of tests from 25 June 2013 to 10 January 2014 \cite{upworthy2024notice}.
% src: DATA.md, The archive's 2024 randomization notice: 1,050 of the 4,873 tests started in the window; $RUN6/paper-extras-upworthy-20261106-2dc461f52ea4.json randomization_window.tests_in_ledger 4873
The file we read has no flag column, and 1,050 of its 4,873 tests started in that window.
We cannot tell which tests are affected, so the main analysis leaves out every test that started in the window: its rows stay in the ledger with that reason (status \emph{imbalance\_flag}) and are not scored. This rule was fixed before any number of the run reported here existed. Including the window is a reported ablation (Table~\ref{tab:fmt-window}), which regenerates the earlier run exactly. Release 1.0.0 of the package and the first version of this paper scored these tests.
% src: $RUN6/paper-extras-upworthy-20261106-2dc461f52ea4.json randomization_window.counts.main.origins_scored 7, randomization_window.counts.include_randomization_window.origins_scored 9
The rule has two consequences for the calendar. Seven of the nine quarterly origins score tests, since the origins of July and October 2013 fall inside the window. The January 2014 origin abstains on every row: each cell's last report on how its forecasts had held up dates from before the gap, so the reliability trigger fires (Section~\ref{sec:method}). That follows from the rule and says nothing about the tests of that quarter.

\paragraph{ASOS Digital Experiments Dataset.} The dataset holds real experiments from an online fashion retailer, with several variants per experiment, several decision metrics and cumulative checkpoints of each group's count, mean and variance \cite{liu2021datasets}. It is the only public archive of product experiments we know of, and it is small (Table~\ref{tab:fmt-archives}). Each contrast is one treatment variant on one metric, read at the final checkpoint; the earlier checkpoints are daily looks a test should not act on, so they are never outcomes. The effect is the standardized mean difference (the mean difference over the pooled standard deviation); the effect in the metric's own units stays in the post-launch block. The cell is the metric, since the content of each experiment is not released.
% src: src/mootcourt/sources/asos.py MIN_USERS 1000 (admission: both arms, finite effect and standard error)
Admission requires 1,000 users in both arms and a finite effect and standard error. The dataset carries no dates, so the loader places experiments on a nominal calendar, sorted by a hash of their identifier and spread evenly over two years. The order is arbitrary: the rolling origin tests the mechanics of the protocol, not drift,
% src: $RUN6/backtest-asos-20261103-6bb5a15c5ff0.json header.tests_scored 68; $RUN6/paper-extras-asos-20261103-6bb5a15c5ff0.json gates_by_stratum.strata.ALL.own_gate.n_400 275
and with 68 scored tests against gates of hundreds, the gates are expected to return \emph{undecidable}.

\begin{table}[t]\centering\scriptsize
% src: $RUN6/backtest-upworthy-20261106-2dc461f52ea4.json header; $RUN6/backtest-asos-20261103-6bb5a15c5ff0.json header; DATA.md, counts produced by the loaders; README.md, License and citations (CC BY 4.0 for both archives)
\caption{The two public archives side by side. Rows: the period, the units the source releases, the three counts of tests and of contrasts (Section~\ref{sec:setting} defines them), the rows each status keeps out of admission, the origins scored, and the license. Counts come from each backtest report; the packages and the ASOS exclusion come from the data documentation. The synthetic archive is described below and in Appendix~\ref{app:synthetic}.}
\label{tab:fmt-archives}
\begin{tabular}{p{0.26\columnwidth}p{0.30\columnwidth}p{0.30\columnwidth}}
\toprule
 & Upworthy Research Archive & ASOS Digital Experiments Dataset \\
\midrule
% src: DATA.md: created from 26 Jan 2013 to 29 Apr 2015; the ASOS dataset carries no dates, $RUN6/backtest-asos-20261103-6bb5a15c5ff0.json header.first_origin 2019-04-01, header.last_origin 2020-10-01 are the loader's nominal calendar
period & January 2013 to April 2015 & none in the data; the loader assigns a nominal calendar \\
% src: DATA.md line 15: 22,666 packages in the exploratory file; ASOS two to five variants per experiment, four decision metrics, metric_1 to metric_4
source units & 22,666 packages & two to five variants per experiment, four decision metrics \\
% src: $RUN6/backtest-upworthy-20261106-2dc461f52ea4.json header.tests_in_ledger 4873; $RUN6/backtest-asos-20261103-6bb5a15c5ff0.json header.tests_in_ledger 78
tests in the ledger & 4,873 & 78 \\
% src: $RUN6/backtest-upworthy-20261106-2dc461f52ea4.json header.tests_admitted 3022; $RUN6/backtest-asos-20261103-6bb5a15c5ff0.json header.tests_admitted 78
tests admitted & 3,022 & 78 \\
% src: $RUN6/backtest-upworthy-20261106-2dc461f52ea4.json header.tests_scored 2949; $RUN6/backtest-asos-20261103-6bb5a15c5ff0.json header.tests_scored 68
tests scored & 2,949 & 68 \\
% src: $RUN6/backtest-upworthy-20261106-2dc461f52ea4.json header.ledger_rows 17793; $RUN6/backtest-asos-20261103-6bb5a15c5ff0.json header.ledger_rows 396
contrasts in the ledger & 17,793 & 396 \\
% src: $RUN6/backtest-upworthy-20261106-2dc461f52ea4.json header.admitted_rows 10903; $RUN6/backtest-asos-20261103-6bb5a15c5ff0.json header.admitted_rows 381
contrasts admitted & 10,903 & 381 \\
% src: $RUN6/backtest-upworthy-20261106-2dc461f52ea4.json header.scored_rows 10697; $RUN6/backtest-asos-20261103-6bb5a15c5ff0.json header.scored_rows 337
contrasts scored & 10,697 & 337 \\
% src: $RUN6/paper-extras-upworthy-20261106-2dc461f52ea4.json randomization_window.ledger_statuses_main duplicate_arm 3321, imbalance_flag 2587, low_clicks 975, low_exposure 7; DATA.md status counts: ASOS low_exposure (missing variance) 15
exclusions by status & duplicate arm 3,321; randomization window 2,587; low clicks 975; low exposure 7 & low exposure (missing variance) 15 \\
% src: $RUN6/backtest-upworthy-20261106-2dc461f52ea4.json header.origins_scored 7; $RUN6/paper-extras-upworthy-20261106-2dc461f52ea4.json randomization_window.counts.include_randomization_window.origins_scored 9; $RUN6/backtest-asos-20261103-6bb5a15c5ff0.json header.origins_scored 7
origins scored & 7 of 9 quarterly & 7 quarterly, nominal \\
% src: README.md, License and citations: Upworthy Research Archive CC BY 4.0; ASOS Digital Experiments Dataset CC BY 4.0
license & CC BY 4.0 & CC BY 4.0 \\
\bottomrule
\end{tabular}
\end{table}

\paragraph{Synthetic archive.} A generator writes a third ledger with known cell means and spreads and a standard error fixed before launch. It is the only source on which the true effects are known, so it checks that the fitted prior recovers them. With a text signal planted, it is the positive control: it tests whether the backtest can find a signal of known size (Section~\ref{sec:synthetic-results}).
% src: $RUN6/backtest-synthetic-20261103-bf4fd286e2b2.json header.tests_in_ledger 1000, header.tests_scored 907
Without the signal it has 1,000 tests, and 907 are scored (Appendix~\ref{app:synthetic}).

\section{Method: The Baseline Suite}
\label{sec:method}

\subsection{Predictive distributions on a grid}
% src: src/mootcourt/predictors/base.py p_win delta 0.0; src/mootcourt/config.py Z 1.96
Each prior rung returns a Student-t for the true effect $\theta$ on a fixed grid; P0 and the text classifier are described below. From that distribution the card derives three quantities:
\begin{itemize}
\item $\pwin = P(\theta > \delta)$, with $\delta = 0$ throughout, the probability that the variant truly wins;
\item the predictive of the estimate the new test will report, $y^\ast = \theta + \varepsilon$ with $\varepsilon \sim \mathcal{N}(0, s_\ast^2)$;
\item $\psig = P(y^\ast > z\, s_\ast)$ with $z = 1.96$, the probability of a significant positive result at the planned size.
\end{itemize}
The predictive of $y^\ast$ is a scale mixture evaluated analytically, so no grid truncation enters a score, and P2's weight on precedents is chosen, and admitted by the guard, under the same density it is scored on. Here $s_\ast$ is the standard error expected before launch: $s_\ast = 100\sqrt{c(1-c)(1/n_{\mathrm{ref}} + 1/n_{\mathrm{var}})}$ on the headline archive, with $c$ the median reference click-through rate of the cell's visible rows, and $\sqrt{1/n_{\mathrm{ref}} + 1/n_{\mathrm{var}}}$ on the standardized scale of the retailer's archive. The level $c$ is the cell's, not the test's own realized rate; reading the realized rate is the leak of Section~\ref{sec:cautionary}. The exposures $n_{\mathrm{ref}}$ and $n_{\mathrm{var}}$ should be planned. No archive records a plan, so realized exposure stands in: impressions on the headline archive, users on the retailer's (Section~\ref{sec:setting}).

\subsection{The ladder}
% src: src/mootcourt/config.py Z 1.96
\paragraph{P0, the base rate.} Per cell and origin, P0 forecasts from the visible history alone:
\begin{itemize}
\item $\psig$ is the share of visible contrasts with a significant positive result ($y > 1.96\,s$);
\item the direction probability is the share of visible estimates above zero;
\item the intervals are the empirical quantiles of the visible $y$;
\item the log-score density is a normal with the visible mean and variance of $y$.
\end{itemize}
Every Brier skill is measured against P0's $\psig$, so P0's skill is zero by construction. P0 is the scoring baseline, not a special case of the rungs above it: P1 adds a fitted, partially pooled, heavy-tailed prior that accounts for each row's measurement error, and P2 adds precedents behind an admission guard; with zero weight on precedents, P2 is P1 exactly.

\paragraph{P1, empirical-Bayes prior.} Per cell, a Student-t prior $t_\nu(\mu_c, \tau_c)$ on $\theta$ is fit by maximizing the marginal likelihood
\begin{equation}
\sum_{i \in c} \log \int \mathcal{N}(y_i \mid \theta, s_i^2)\, t_\nu(\theta \mid \mu_c, \tau_c)\, d\theta ,
\end{equation}
% src: mootcourt-data models 938be3f002b0/model.json and d3e813b146ce/model.json, p1 nu 3 at every level (the Upworthy and ASOS bundles); src/mootcourt/predictors/p2.py docstring: tau sqrt(nu / (nu - 2))
with the integral on the grid and $\nu$ chosen from 3, 5, 10, 30 and the normal limit. Its standard deviation is $\sigma_c = \tau_c\sqrt{\nu/(\nu-2)}$.
% src: src/mootcourt/config.py KAPPA 20, NU_GRID (3, 5, 10, 30, inf); src/mootcourt/predictors/p1.py MIN_PERIOD_ROWS 30
Cell parameters are partially pooled toward the global fit with weight $n_c / (n_c + \kappa)$, $\kappa = 20$:
\begin{itemize}
\item once along the kind of change and once along the calendar period, and the two shrunk values are averaged;
\item the period component enters as half its deviation from the global fit;
\item a period with fewer than 30 visible rows, or one ending after the fit's data, contributes nothing;
\item a fit inside a quarter takes the last complete quarter's level (at the quarterly origins of the backtest the two rules coincide).
\end{itemize}
% src: mootcourt-data models 938be3f002b0/model.json and d3e813b146ce/model.json, p1 nu 3 at every level
The fits for both archives select $\nu = 3$, the heaviest tail on the grid; the pooling of the scale needs $\nu > 2$. The data therefore prefer tails at least this heavy, and a heavier tail was not tried.

% src: src/mootcourt/features.py PAIRWISE_SVD_DIMS 32; src/mootcourt/harness.py run loop (every model at an origin, the feature space included, is fit on visible_rows(adm, t))
\paragraph{P2, weighted precedents.} \emph{Features.} The features $x$ of a contrast describe the pair. One TF-IDF and truncated SVD, fit at each origin on the visible rows' headlines, gives 32 columns each for the variant headline, the reference headline and their difference. To those are added the length differences, the share of words in common, the kind of change, the number of arms and the log planned exposure; realized exposure stands in for the plan (Section~\ref{sec:setting}), so this is the one feature that reads a post-launch quantity.

\emph{Estimator.} With a kernel $k(x, x_i) = \exp(-d^2/2h^2)$ on standardized features and weights $w_i = k(x, x_i)/(s_i^2 + \sigma_c^2)$ over visible rows,
\begin{align}
m(x) &= \frac{\sum_i w_i y_i}{\sum_i w_i}, \qquad
v_m(x) = \frac{\sum_i w_i^2 (s_i^2 + \sigma_c^2)}{(\sum_i w_i)^2}, \\
\hat\theta(x) &= \lambda\, m(x) + (1-\lambda)\, \mu_c, \\
v(x) &= \lambda \left(v_m(x) + \tau_{\mathrm{res},c}^2\right) + (1-\lambda)\, \sigma_c^2 ,
\end{align}
where $\tau_{\mathrm{res},c}^2$ is the residual between-test variance after kernel regression. The working predictive is a Student-t of the cell's own $\nu$ with location $\hat\theta(x)$ and scale $\sqrt{v(x)(\nu-2)/\nu}$, so its variance is $v(x)$, and $\lambda = 0$ returns P1 exactly. A weight above zero moves the location and the scale, never the family.

\emph{Selection.} The weight $\lambda$ and the bandwidth multiplier (a multiple of the median heuristic \cite{garreau2017large}; Table~\ref{tab:fmt-protocol} lists both grids) are chosen per cell by the out-of-time log score on forward-chained validation blocks. Each block is scored with the precedents created before it, so a validation row never borrows a precedent from its own test or its future (leave-one-test-out is kept as an ablation). Every candidate, $\lambda = 0$ included, is scored with the same marginal density as the P1 likelihood.

% src: src/mootcourt/predictors/p2.py _LAMBDA_TIE 1e-3 (0.001 nats); src/mootcourt/config.py P2_GUARD_MULT 2.0 (two standard errors by test); Table tab:fmt-protocol, admission guard row
\emph{Admission guard.} The guard limits how much noise the selection can fit. A candidate with $\lambda > 0$ is admitted only when its per-row log-score gain over $\lambda = 0$ on the validation rows exceeds 0.001 nats and two standard errors clustered by test; otherwise $\lambda = 0$, and P2 equals P1 for that cell and origin. The guard does not promise that P2 beats P1 on later tests; Sections~\ref{sec:ladder-results}, \ref{sec:synthetic-results} and~\ref{sec:ablations} measure what it does.

\paragraph{Designs with no result here.} The package holds two more parts of the design, and neither has a result in this paper. P3 is a probability and a point estimate the team enters before the card is seen, frozen with a timestamp; no public archive holds historical forecasts. The \emph{second opinion} is an optional language-model forecast $p_{\mathrm{AI}}$, stacked as $p_{\mathrm{stack}} = (1 - w)\, \psig + w\, p_{\mathrm{AI}}$ with $w = 0$ until the stacked Brier score beats $\psig$ out of time at the gate's sample size.

\paragraph{Challengers.} Three learned models run on the same origins and rows. Each is reported raw and after a Platt recalibration \cite{platt2000probabilities} fit on out-of-fold predictions of the visible rows only.
\begin{itemize}
\item The \emph{boosted model} (\texttt{GBM} in the tables) is a gradient-boosted regressor on P2's pairwise features, exposure included. It is weighted by $1/(s^2 + \tau^2)$, with $\tau^2$ the DerSimonian-Laird between-test variance, and uses a time-ordered search over its learning rate, leaves and boosting rounds.
% src: src/mootcourt/config.py LOCAL_EB_K 50
\item The \emph{local comparison set} (\texttt{local\_eb}) is an empirical-Bayes prior fit on the 50 nearest visible rows of the query's cell, in the same features, with no shrinkage toward the cell prior, the shape of \cite{li2025learning}.
% src: src/mootcourt/harness.py TextLR._design (TF-IDF of the headlines and one column per cell; no exposure)
\item The \emph{text classifier} (\texttt{text\_lr}) is a logistic regression on the TF-IDF of the variant headline, its difference from the reference headline, and one column per cell. It reads no exposure, so its ranking does not come from the traffic stand-in. It is fit once to $y > 0$ for its win probability, and once to the significant-positive event for its significance probability. Its win probability therefore estimates $P(y > 0)$, not $P(\theta > \delta)$, and it issues no density. The plan named a text-only model; the pairwise design and the cell columns came after the plan.
\end{itemize}
% src: src/mootcourt/config.py LOCAL_EB_K 50 and GBM_CONFIGS (four configurations); src/mootcourt/harness.py TextLR C 1.0, min_df 3; $RUN6/ablations/backtest-upworthy-20261106-2dc461f52ea4--prior-normal.json tables.aggregate[cell=ALL,rung=P1] logscore -0.8316 against the main run's -0.7566 (0.075 nats)
Only the boosted model is searched. The classifier keeps scikit-learn's default regularization, and the local comparison set a neighbour count fixed by hand, so the text result is untuned and a search could move it. The boosted model and the local comparison set issue a normal predictive. Their log scores likely pay a penalty for that family: for P1, replacing the Student-t with a normal costs 0.075 nats (Section~\ref{sec:ablations}). Platt recalibration corrects their probabilities, not their densities, so the recalibrated rows keep that penalty.
The boosted model, like the text classifier, receives one indicator column per cell, because the ladder fits each cell apart (Appendix~\ref{app:walltime} records how we found this was needed). Two direction baselines, always the reference arm and always the variant, are scored beside every rung; the reference is the earliest-created arm, so the first measures the creation-order effect.

\subsection{Abstention, the card and the loop}
% src: src/mootcourt/config.py SUPPORT_MIN 30, NOVELTY_PERCENTILE 5, SAME_SIDE_FRACTION 0.80, ACTION_THRESHOLD 0.35
The card abstains, listing every reason, when any of four triggers fires:
\begin{itemize}
\item \textbf{support:} fewer than 30 visible rows in the cell after pooling;
\item \textbf{novelty:} the kernel density of $x$ among the cell's visible rows falls below the fifth percentile of their leave-one-out densities;
\item \textbf{reliability:} the cell's reliability report is missing, stale, future-dated or drifted;
\item \textbf{decision instability:} in a bootstrap over precedents, resampled by test, $\psig$ lands on the same side of the action threshold (Table~\ref{tab:fmt-protocol}) in 80 percent or fewer of resamples (off by construction where $\lambda = 0$).
\end{itemize}
An abstained card still shows the base rate, labelled as what it declined to improve on. It keeps the withheld forecast in the frozen file, which the join scores and no interface shows. Abstention was designed as a filter on quality. Because issued rows did not score better on skill, this paper reads it only as a label of evidence (Section~\ref{sec:abstention-results}).

The card is issued before launch and frozen with:
\begin{itemize}
\item the identifiers of the snapshot, model and feature set;
\item the predictive itself, so the join scores the outcome with the card's own object;
\item the prediction: $\psig$, a point estimate, intervals for the true effect and the observed estimate, and $\pwin$;
\item the frozen team forecast;
\item its cell's latest Brier skill read against the gate, with its interval and test count (the pooled figure, labelled, where the cell is undecidable);
\item the abstention reasons, and a fixed sentence stating that the live test decides.
\end{itemize}
In use, the package runs a loop, nightly and weekly; Table~\ref{tab:fmt-protocol} gives its constants. The backtest applies the reliability, drift and staleness rules below. The weekly join is replayed in Appendix~\ref{app:replay}. The promotion rule is design only and has no result here.
\begin{itemize}
% src: docs/method.md calls 220 to 225, 248 (the seed dated at the fit clock) and 252 (the doctor's count of cards on an unadmitted test); src/mootcourt/registry.py paired_comparison (90 percent interval by test, harness_stats.DEFAULT_B 500 resamples) and date_reliability_seed, src/mootcourt/config.py PROMOTION_MARGIN 0.0, drift margin 0.1 nats, RELIABILITY_MAX_AGE_DAYS 14; src/mootcourt/loop.py card_timing and attach_outcomes (admitted rows only); src/mootcourt/card.py backtest_view
\item \textbf{Nightly fit.} A fit on the rows usable at its clock writes a versioned, hashed bundle. A refit of the same method on newer data replaces the current model once its backtest shows a skill interval not entirely below the current model's; until then it waits as a pending refresh. A different method replaces the current one only when, on common rows, the lower bound of the 90 percent by-test interval of its paired Brier skill over the current model is above zero, on at least the gate's number of tests.
\item \textbf{Weekly join.} The join attaches outcomes to cards whose assigned end has passed, issued and abstained alike. Each row is scored on one card: the earliest whose origin and issue time are at or before the row's start; a later card is a reissue. A card issued during or after the test is left unscored, a card whose test is not admitted is never scored, and an origin after the issue time is refused. A card's stated times are the team's word, so the join checks them against the contrast's own start: on the headline archive, the later of its two arms' creation times. A card for a variant added to a running test is therefore pre-launch evidence for that contrast, not for the test.
% src: src/mootcourt/config.py DRIFT_LOGSCORE_MARGIN 0.10; src/mootcourt/harness.py build_reliability in the backtest loop
\item \textbf{Reliability and drift.} The join classifies every miss and rewrites a per-cell reliability report over a trailing window. A cell drifts when its rolling log score falls more than 0.1 nats below the base rate's at consecutive joins with enough outcomes; its cards then abstain until a join clears it.
% src: src/mootcourt/config.py RELIABILITY_MAX_AGE_DAYS 14
\item \textbf{Staleness.} The reliability trigger also has a staleness limit of 14 days. The backtest dates each cell's report at the ledger's last result, moved to the fit clock when that result is within the limit. A ledger whose last result is older than the limit therefore issues no card until the first join brings outcomes. This cold start is a known limit.
\end{itemize}
Cards are served as a command, a read-only service and a tool interface, each keyed on the ledger row.

\section{Evaluation Protocol}
\label{sec:protocol}
\paragraph{Rolling origin.} At each quarterly origin $t_k$, the predictors are fit on the rows visible at $t_k$ and score every admitted contrast that started in $[t_k, t_{k+1})$, with the pre-launch $s_\ast$ of Section~\ref{sec:method} as the forecast's standard error \cite{tashman2000outofsample}. The realized standard error enters only the outcome, never the forecast: it defines the significant-positive event (the observed estimate exceeds 1.96 times it), our main definition, and the floor's ingredients. The same event defined on the planned standard error is scored beside it (Table~\ref{tab:event-definition}). A guard around every fit, the challengers and their calibrators included, stops the run if any scored row was visible to the model that scored it. After the run, an audit recomputes that no fit row was usable after its origin and that $s_\ast$ equals the pre-launch value.

\paragraph{Scores.} Every forecast is scored five ways:
\begin{itemize}
\item \textbf{log score} of the $y^\ast$ predictive at the observed $y$, in nats per row, higher better; and the \textbf{CRPS}, in the units of $y$, lower better;
\item \textbf{Brier score} of $\psig$ against the significant-positive event, beside its irreducible floor $\mathbb{E}[p_\theta (1 - p_\theta)]$, what a forecaster who knew each true effect would score;
\item \textbf{Brier skill} $1 - \mathrm{Brier}/\mathrm{Brier}_{\mathrm{base}}$ against the cell's visible event rate, the quantity P0 forecasts, so skill above zero means beating the base rate;
% src: src/mootcourt/config.py Z 1.96
\item \textbf{direction accuracy} against its ceiling $\mathbb{E}[\Phi(|\theta|/s)]$, on all rows and on the decisive rows ($|y| > 1.96\,s$) that the winner-prediction literature scores;
% src: src/mootcourt/harness_stats.py RESOLUTION_BINS 10 (ten bins); murphy_resolution docstring, bias (n_bins - 1)/N times obar(1 - obar), so 9/N
\item \textbf{ranking:} the area under the receiver operating characteristic curve (AUC) of each probability, and its Murphy resolution term over ten equal-mass bins. Resolution is biased upward by about $9/N$ times the rate times one minus the rate, so it is read as a ranking of rungs on the same rows.
\end{itemize}
Floor and ceiling integrate the shrunken posterior of each row's true effect under the pooled prior, with the Student-t replaced by the normal of the same standard deviation, a stated approximation. They sit beside every score because there is no ground truth for $\theta$. The floor is an in-hindsight bound, since it reads the realized estimate the forecast could not.

% src: src/mootcourt/harness_stats.py LEVEL 0.90 (90 percent); $RUN6/backtest-upworthy-20261106-2dc461f52ea4.json header.bootstrap_resamples 500
\paragraph{Intervals and families.} Every score has a 90 percent percentile bootstrap interval in its report, from 500 resamples of whole tests, so the contrasts of one test stay together. The tables print only the intervals the text reads; the rest are in the reports. The calibration null and the gate simulation are simulations, not resamples, and use the same 90 percent by-test criterion.

% src: src/mootcourt/harness_stats.py apply_bonferroni, normal quantile of 1 - 0.05/family (LEVEL 0.90, two tails); tables/ece-upworthy.tex adjusted bounds (m)
One table holds many intervals, and some would exclude zero by chance alone. So the aggregate, year, calibration and decision-value tables also print Bonferroni-adjusted bounds, chosen so that all the intervals in the table cover together with probability at least 90 percent. Each bound is the bootstrap standard error times the normal quantile for $1 - 0.05/m$, where $m$ is the number of intervals in the table. Intervals in the other tables are read as descriptive, and so are per-cell and per-year intervals unless a table prints their adjusted bounds. The pooled interval is the primary claim, a rule fixed in this paper after the data were seen (Section~\ref{sec:hypotheses}).

\paragraph{Calibration and coverage.} Three checks:
\begin{itemize}
\item \textbf{Calibration.} The expected calibration error (ECE) is the gap between forecast probability and observed frequency, averaged over bins of the forecast with each bin weighted by its rows. It is computed over equal-width and equal-mass bins. We compare it with its distribution under a perfect-calibration null, which draws each outcome from the forecast's own probability. Contrasts of one test share a reference arm, so their outcomes move together. The null therefore draws them together through a one-factor Gaussian copula. Its correlation is solved so that the null's design effect, the factor by which this sharing inflates the variance, equals the one measured on the rows. The design effect and the null with independent outcomes are reported beside it.
% src: src/mootcourt/harness_stats.py interval_coverage of lo80/hi80 and lo90/hi90 (coverage_80, coverage_90)
\item \textbf{Coverage} of the 80 and 90 percent intervals for $y$.
\item \textbf{By precision.} Calibration within deciles of planned exposure, and coverage within deciles of the ratio of realized to planned standard error. The forecaster never reads the realized standard error, so this second check shows whether intervals fail where the test turned out noisier than planned.
\end{itemize}
% src: src/mootcourt/harness_stats.py NULL_DRAWS 500, the percentile resolves at most (draws - 1) / draws; null_percentile, 100 when above every draw
The number of null draws sets the resolution of the percentile (Table~\ref{tab:fmt-protocol}). A percentile of 100 means the observed ECE is above every draw, and a percentile above 95 rejects calibration.

% src: src/mootcourt/simulate.py DEFAULT_TARGET 0.90 (power), DEFAULT_LEVEL 0.90 (interval), GATE_SKILLS n_200 0.20, n_400 0.10, RESEARCH_TARGET 0.80
\paragraph{Gates.} A skill claim about a group of tests, a \emph{stratum}, is decidable when its scored tests reach that stratum's gate. The gate is the number of tests at which a simulated calibrated forecaster with a given Brier skill has power 0.90 to clear zero, that is, to put the lower end of its 90 percent by-test bootstrap interval above zero. There are two gates, one for a skill of 0.20 and one for a skill of 0.10. A research screen at power 0.80 is reported beside them. One rule applies to every stratum and rung.

The simulation copies the stratum: its base rate, its rows per test and the within-test correlation of outcomes measured on its own rows (Table~\ref{tab:fmt-gates-strata}). Each simulated test gets one calibrated probability, drawn from a Beta distribution whose mean is the base rate and whose variance gives the stated Brier skill exactly.

A gate is a sizing convention. It assumes calibration, which Section~\ref{sec:calibration-results} rejects, so a test count above it says only that the interval is worth reading; the interval is the evidence.

% src: $RUN6/gates.json by_base_rate["0.09"].icc 0.13, by_base_rate["0.30"].icc 0.13, dependence_source "the public archive nearest this base rate"; $RUN6/paper-extras-upworthy-20261106-2dc461f52ea4.json gates_by_stratum.strata.ALL.dependence.icc 0.197; docs/method.md call 183 (values measured on earlier backtests, kept in simulate.SOURCE_DEPENDENCE)
The package itself still labels its reports with an older simulation at three fixed base rates (Table~\ref{tab:gates}). There, the correlation for each rate was taken from the nearest public archive on an earlier backtest. At the headline archive's base rate it assumed 0.13, where this run's rows measure 0.20. Its verdicts agree with the rule above on every stratum of both archives.

% src: src/mootcourt/decision_value.py BUDGETS (0.25, 0.10, 0.50), main 0.25 (Table tab:fmt-protocol); POLICIES random, random_matched, arm_count, p_win, oracle (five)
\paragraph{Decision value.} At each origin, a team with a budget of $K$ tests picks which of the period's $n$ tests to run. The budget is 25 percent of the period's tests, with 10 and 50 percent reported beside it (Table~\ref{tab:fmt-protocol}). A test's value is its best positive observed lift, or zero if every contrast lost, in the units of $y$. Five orders pick the $K$ tests:
\begin{itemize}
\item \textbf{random}, whose expected capture is exactly $K/n$ of the period's value, by linearity;
\item \textbf{random matched on arms}: random within groups of tests with the same number of arms, matched to the arm counts of the set a rung picks;
\item \textbf{arm count}: the tests with the most contrasts, reading no forecast;
\item \textbf{$\pwin$}: a test's score is its largest contrast's $\pwin$ under a rung;
\item \textbf{oracle}: the tests with the highest value, which needs the outcome.
\end{itemize}
% src: src/mootcourt/decision_value.py LEVEL 0.90, BOOTSTRAP 500 resamples of tests within each origin
Decision value is the share of the period's value an order captures, read as a difference from random order or from random order matched on arms. A test with more arms has a larger best lift by chance alone, and its arm count is known before launch. So the matched order is the baseline a rung must beat (Section~\ref{sec:decision-results}). Each order is reported with its regret against the oracle and a 90 percent bootstrap interval that resamples tests within each origin. A rung whose $\pwin$ is constant within a cell (P0, and P1 and P2 wherever $\lambda = 0$) chooses by cell and then at random, so its capture differs from random order only through differences between cells.

\paragraph{Test size.} A concluded test ran once, at one size, so sizing cannot be replayed. We measure two things instead, on the headline archive only.

% src: src/mootcourt/sizing.py sizing_assurance n_bins 5 (five equal-mass bins of exposure); docs/superpowers/plans/2026-11-06-x10-sizing-design.md "upworthy only"
First, we check whether the forecast probability of a significant positive result, at the size a test actually had, matches the realized rate across five equal-mass bins of exposure (Table~\ref{tab:fmt-assurance}).

Second, we simulate the test-and-roll problem \cite{feit2019test}, in three steps:
\begin{enumerate}
% src: src/mootcourt/sizing.py AUDIENCE_MULTIPLES 5, 20, 100, PRIMARY_MULTIPLE 20; n_med median exposure of the cell's visible rows
\item A page has an audience of 20 times the median exposure of the cell's visible rows, with 5 and 100 times beside it.
\item A rule sets the impressions per arm.
\item The arm with the higher observed rate is shown to the rest of the audience.
\end{enumerate}
Gain is counted in clicks, with no money attached. Each rule reads only data visible at the origin.
\begin{itemize}
\item The \emph{prior rule} maximizes expected gain under P1's prior.
\item The \emph{base-rate rule} maximizes it under a normal prior whose spread reproduces the cell's visible rate of significant positive results.
% src: src/mootcourt/sizing.py POWER_LEVEL 0.80, POWER_ALPHA 0.05, POWER_REL_LIFT 0.10
\item The \emph{fixed power rule} takes the size with 80 percent power at the 5 percent level for a 10 percent relative lift.
\end{itemize}
% src: src/mootcourt/sizing.py WORLD_SCALES 1.0, 0.67, 1.5 (effects theta times the scale); npmle (Kiefer and Wolfowitz); design note, relative regret
A rule's regret is what it loses against an oracle that shows the better arm to the whole audience. The primary figure is relative regret: regret summed over origins and cells, divided by the summed oracle gain. Regret is measured in a world that is not the prior being tested: the nonparametric maximum-likelihood distribution of effects fitted to the rows scored at that origin and cell, with its effects scaled by 0.67 and 1.5 beside it (Table~\ref{tab:fmt-sizing}). A design note fixed all of this before any sizing number existed.

\begin{table}[t]\centering\scriptsize
% src: src/mootcourt/config.py; src/mootcourt/reliability.py MIN_DRIFT_N 10; src/mootcourt/harness_stats.py NULL_DRAWS 500; $RUN6/gates.json target_power 0.9, research_target 0.8; $RUN6/backtest-upworthy-20261106-2dc461f52ea4.json tables.decision_value budget 0.1, 0.25, 0.5; src/mootcourt/sizing.py
\caption{Protocol settings: the defaults and constants of the method, the loop and the evaluation, defined in words in the text and read by every section after it. Rows: each setting, grouped by the part of the system that reads it. Sources: \texttt{config.py} for the forecast, the rung, the triggers and the loop, \texttt{reliability.py} for the outcome minimum of the drift rule, \texttt{harness\_stats.py} for the null, \texttt{gates.json} for the gate, \texttt{tables.decision\_value} of the backtest report for the budgets and \texttt{sizing.py} for test size. Nothing here is an estimate, so no row carries an interval.}
\label{tab:fmt-protocol}
\begin{tabular}{p{0.46\columnwidth}p{0.46\columnwidth}}
\toprule
setting & value \\
\midrule
\multicolumn{2}{l}{\emph{Forecast}} \\
% src: src/mootcourt/predictors/base.py p_win delta 0.0
win margin $\delta$ & 0 \\
% src: src/mootcourt/config.py Z 1.96
significance multiplier $z$ & 1.96 \\
% src: src/mootcourt/config.py ACTION_THRESHOLD 0.35
action threshold on $\psig$ & 0.35 \\
\midrule
\multicolumn{2}{l}{\emph{Precedent rung (P2)}} \\
% src: src/mootcourt/config.py LAMBDA_GRID 0.0 to 1.0 in steps of 0.1
mixing weight $\lambda$, chosen per cell & 0, 0.1, and so on to 1 \\
% src: src/mootcourt/config.py BANDWIDTH_MULTIPLIERS 1.0, 0.5, 0.25, 0.125
bandwidth multiplier on the median heuristic & 1, 1/2, 1/4, 1/8 \\
% src: docs/method.md line 436, tie tolerance of 0.001 nats; src/mootcourt/config.py P2_GUARD_MULT 2.0
admission guard, per-row gain on validation rows & above 0.001 nats and two standard errors by test \\
\midrule
\multicolumn{2}{l}{\emph{Abstention triggers}} \\
% src: src/mootcourt/config.py SUPPORT_MIN 30
support, visible rows in the cell & fewer than 30 \\
% src: src/mootcourt/config.py NOVELTY_PERCENTILE 5
novelty, density among the cell's visible rows & below the fifth percentile \\
% src: src/mootcourt/config.py SAME_SIDE_FRACTION 0.80
decision instability, resamples on one side & 80 percent or fewer \\
\midrule
\multicolumn{2}{l}{\emph{The loop}} \\
% src: src/mootcourt/config.py PROMOTION_MARGIN 0.0; src/mootcourt/registry.py paired_comparison (90 percent interval by test)
promotion margin on the lower bound of the 90 percent interval & 0 \\
% src: src/mootcourt/config.py TRAILING_WEEKS 8
reliability report window & trailing 8 weeks \\
% src: src/mootcourt/config.py DRIFT_LOGSCORE_MARGIN 0.10, DRIFT_CONSECUTIVE 2; src/mootcourt/reliability.py MIN_DRIFT_N 10
drift rule & 0.1 nats below the base rate at 2 consecutive joins, 10 outcomes at least \\
% src: src/mootcourt/config.py RELIABILITY_MAX_AGE_DAYS 14
staleness limit of a reliability report & 14 days \\
\midrule
\multicolumn{2}{l}{\emph{Calibration and coverage}} \\
% src: src/mootcourt/harness_stats.py ece_equal_mass n_bins 10
expected calibration error bins & 10 equal-width and 10 equal-mass \\
% src: $RUN6/backtest-upworthy-20261106-2dc461f52ea4.json tables.aggregate keys coverage_80, coverage_90
coverage levels & 80 and 90 percent \\
% src: src/mootcourt/harness_stats.py NULL_DRAWS 500; the percentile resolves at most (draws - 1) / draws, 499 / 500 = 99.8 percent
null draws & 500, resolving percentiles to 99.8 \\
% src: Section 5, the rejection rule
percentile that rejects calibration & above 95 \\
\midrule
\multicolumn{2}{l}{\emph{Gates, decision value and test size}} \\
% src: $RUN6/gates.json decision_target 0.9, research_target 0.8; interval.level 0.9 by test
gate power on a 90 percent by-test interval & 0.90 for release, 0.80 reported beside it \\
% src: $RUN6/backtest-upworthy-20261106-2dc461f52ea4.json tables.decision_value budget 0.1, 0.25, 0.5
decision budgets, share of the period's tests & 10, 25 and 50 percent \\
% src: src/mootcourt/sizing.py AUDIENCE_MULTIPLES 5, 20, 100, PRIMARY_MULTIPLE 20
audience multiple in the sizing simulation & 20, with 5 and 100 beside it \\
\bottomrule
\end{tabular}
\end{table}

\paragraph{Hypotheses and their reading.}\label{sec:hypotheses}
The plan's six hypotheses, verbatim (Section~\ref{sec:verdicts} gives when they were written and reads them back):
\begin{enumerate}
% src: paper/plan.md line 27, verbatim
\item[H1] P1 beats P0 on log score at every origin after burn-in on Upworthy; P2 beats P1 on the headline\_only cell where text features carry signal, and does not beat P1 on image\_only.
% src: paper/plan.md line 28, verbatim
\item[H2] Brier skill of P2 against the base rate is positive and its 90 percent bootstrap interval clears zero at the final origin on Upworthy; the released gate (power 0.90) is decidable there.
% src: paper/plan.md line 29 (80 and 90 percent, 5 points), verbatim
\item[H3] ECE of issued cards falls inside the central 90 percent of the perfect-calibration null at most origins; 80 and 90 percent coverage are within 5 points of nominal.
% src: paper/plan.md line 30 (30 percent), verbatim
\item[H4] Issued cards have higher skill than the full set; abstained rows have lower skill when forced; abstention rate is below 30 percent after burn-in.
% src: paper/plan.md line 31 (K = 25 percent), verbatim
\item[H5] Prioritizing by p\_win captures more positive lift than random at K = 25 percent of tests, with regret against the oracle reported as is.
% src: paper/plan.md line 32, verbatim
\item[H6] On ASOS every gate is undecidable and P2 equals P1 (cell-only features); the paper reports this as the expected outcome.
\end{enumerate}
The plan's words map to this paper's as follows: \emph{burn-in} is the first origin; \emph{issued cards} are rows whose card does not abstain; \emph{the released gate (power 0.90)} is the gate defined above; \emph{the final origin} is read two ways in Section~\ref{sec:verdicts}, as the last calendar origin and as the last origin with enough tests to clear a gate.

We read the hypotheses by rules fixed in this paper after the data were seen:
\begin{itemize}
\item Each hypothesis is split into claim parts, and each part is read from the pooled interval of the table Section~\ref{sec:verdicts} names. Per-cell and per-year intervals are descriptive.
\item A clause that states when a part can be read is not a part. In H2, ``the released gate is decidable there'' says when skill can be read; it is not a claim that can hold or fail. The same goes for H5's ``regret reported as is''.
\item H2's final-origin part and H3's most-origins part are read pooled. Section~\ref{sec:verdicts} also gives each as written.
\item H5 names random order. Its verdict reads the stricter order matched on the number of arms, and Section~\ref{sec:verdicts} gives both.
\end{itemize}
A part takes one of three words, and a part read two ways takes one for each:
\begin{itemize}
\item \emph{Held}: the reading meets the part's claim.
\item \emph{Failed}: the estimate lies on the wrong side of the claim, or a test of the claim rejects it.
\item \emph{Not shown}: the estimate lies on the claim's side, but its interval also contains no effect, so this run cannot tell the claim from its absence. A part not shown has not failed.
\end{itemize}
A hypothesis whose parts differ is marked \emph{in part}.

\section{Results}
\label{sec:results}
Every number in this section is rendered from a report JSON file by \texttt{paper/build\_tables.py}; the file path and snapshot identifier are written as a comment at the top of each table.

\subsection{A cautionary run: ratio metrics and realized precision}
\label{sec:cautionary}
The first run of the protocol defined the effect as relative lift against the earliest-created arm and handed the forecaster each row's realized standard error. Both choices are common in practice, and both let post-launch information into a pre-launch question. Table~\ref{tab:cautionary} shows what that does, on the same 13,247 contrasts and nine origins, reproduced by one flag of the released code. By decile of the realized standard error the mean effect runs from minus 39 to plus 80 percent and the significant-positive rate from 0.5 to 21 percent, because a ratio against a noisy reference arm moves with the reference's realized click-through rate. The empirical-Bayes rung, which assumes the effect is independent of its precision, predicts the opposite gradient: a 16 percent chance of a significant positive result in the most precise decile, where it happened 0.5 percent of the time, and 3 percent in the least precise, where it happened 21 percent of the time. Its 90 percent intervals cover 65 percent of outcomes at both extremes. Its Brier skill against the base rate was minus 0.107, a number that a reader without this table would have taken as evidence that the ledger has no signal. Nothing in that run was a coding error; the protocol was wrong. Every number after this subsection uses the pre-launch-clean protocol of Section~\ref{sec:setting}: an absolute effect, a planned-exposure standard error from the cell's trailing level, and a feature allow-list that cannot read anything measured during the test.

\begin{table}[t]\centering\scriptsize
\caption{The cautionary run on the headline archive (relative lift, realized precision), by decile of realized standard error.}
\label{tab:cautionary}
\resizebox{\columnwidth}{!}{% generated by paper/build_tables.py from paper-extras-upworthy-20261106-2dc461f52ea4.json
\begin{tabular}{rrrrrr}
\toprule
decile of realized $s$ & median $s$ & mean $y$ & event rate & P1 mean $\psig$ & P1 cov.\ 90 \\
\midrule
1 & 9.8 & -42.5 & 0.001 & 0.141 & 0.62 \\
2 & 12.7 & -29.2 & 0.001 & 0.104 & 0.89 \\
3 & 14.9 & -23.7 & 0.009 & 0.085 & 0.97 \\
4 & 17.0 & -17.2 & 0.009 & 0.071 & 0.99 \\
5 & 19.3 & -10.2 & 0.017 & 0.061 & 0.99 \\
6 & 21.9 & -4.8 & 0.022 & 0.052 & 0.98 \\
7 & 25.1 & 3.3 & 0.032 & 0.045 & 0.96 \\
8 & 29.1 & 14.7 & 0.065 & 0.038 & 0.92 \\
9 & 36.2 & 30.7 & 0.091 & 0.032 & 0.86 \\
10 & 52.0 & 84.4 & 0.200 & 0.026 & 0.66 \\
\bottomrule
\end{tabular}
}
\end{table}

\subsection{Skill ladder on the headline archive}
\label{sec:ladder-results}
Table~\ref{tab:main-upworthy} gives the pooled scores over nine quarterly origins and 3,797 held-out tests. The empirical-Bayes prior (P1) improves the logarithmic score over the base-rate rung by 0.145 nats and is calibrated: its 80 and 90 percent intervals cover 81 and 89 percent of outcomes, and its Brier skill against the cell base rate is 0.001 with a 90 percent interval from minus 0.002 to 0.005, which is to say indistinguishable from zero. The ladder's direction accuracy of 0.573 against a ceiling of 0.823 equals that of the \emph{always\_reference} baseline (0.578): on this archive the only direction signal the ledger carries is that the earliest-created arm tends to win, and P1 learns exactly that through the cell mean. Every learned challenger loses to P1: the gradient-boosted model has Brier skill of minus 0.137 and the local empirical-Bayes comparison set minus 0.015. The precedent rung (P2) runs with the admission guard of Section~\ref{sec:method}: a precedent weight above zero is admitted only when its out-of-time gain over the prior clears two standard errors across tests. The guard leaves lambda at zero in 93 percent of cell-origin pairs (the guard itself holds it there in 52 percent, the selection in the rest), and where P2 does speak it does not help: its pooled logarithmic score is minus 0.896 against P1's minus 0.881 and its Brier skill minus 0.008 (minus 0.012 to minus 0.003). Without the guard the same rung falls to minus 0.935 and a Brier skill of minus 0.030. The guard recovers most of the loss and none of the gain, because on this archive the precedents carry no gain to recover.

\begin{table*}[t]\centering\scriptsize
\caption{Headline archive, pooled over origins and cells. Intervals are bootstrap by test.}
\label{tab:main-upworthy}
\resizebox{\textwidth}{!}{% generated by paper/build_tables.py from backtest-upworthy-20261103-4bfb354b87e3.json (source upworthy, snapshot 20261103-4bfb354b87e3, model dd73c75e92bd, 13247 scored rows over 9 origins; gate decision target 0.9)
\begin{tabular}{llrrrrrrrrrr}
\toprule
kind & model & tests & log score & minus P1 (90\%) & CRPS & Brier / base / floor & Brier skill (90\%) & dir.\ acc.\ / ceiling & cov.\ 80 / 90 & abst. & decid. \\
\midrule
challenger & GBM & 3797 & -0.989 & -0.1325 (-0.1450, -0.1195) & 0.34 & 0.0905 / 0.0832 / 0.0553 & -0.088 (-0.102, -0.074) & 0.574 / 0.823 & 0.90 / 0.95 & 0.141 & yes \\
challenger & GBM+cal & 3797 &  &  &  & 0.0831 / 0.0832 / 0.0553 & 0.000 (-0.003, 0.004) &  / 0.823 &  & 0.141 & yes \\
ladder & P0 & 3797 & -1.006 & -0.1493 (-0.1621, -0.1358) & 0.34 & 0.0832 / 0.0832 / 0.0553 & 0.000 (0.000, 0.000) & 0.576 / 0.823 & 0.85 / 0.93 & 0.141 & yes \\
ladder & P1 & 3797 & -0.857 & 0.0000 (0.0000, 0.0000) & 0.33 & 0.0830 / 0.0832 / 0.0553 & 0.001 (-0.002, 0.005) & 0.573 / 0.823 & 0.81 / 0.89 & 0.141 & yes \\
ladder & P2 & 3797 & -0.856 & 0.0004 (-0.0003, 0.0012) & 0.33 & 0.0830 / 0.0832 / 0.0553 & 0.002 (-0.002, 0.005) & 0.575 / 0.823 & 0.81 / 0.89 & 0.141 & yes \\
baseline & always\_reference & 3797 &  &  &  &  &  & 0.578 / 0.823 &  & 0.141 & yes \\
baseline & always\_variant & 3797 &  &  &  &  &  & 0.422 / 0.823 &  & 0.141 & yes \\
challenger & local\_eb & 3797 & -0.995 & -0.1387 (-0.1614, -0.1152) & 0.33 & 0.0838 / 0.0832 / 0.0553 & -0.008 (-0.014, -0.002) & 0.560 / 0.823 & 0.81 / 0.88 & 0.141 & yes \\
challenger & local\_eb+cal & 3797 &  &  &  & 0.0831 / 0.0832 / 0.0553 & 0.001 (-0.002, 0.003) &  / 0.823 &  & 0.141 & yes \\
challenger & text\_lr & 3797 &  &  &  & 0.0828 / 0.0832 / 0.0553 & 0.005 (-0.001, 0.011) & 0.555 / 0.823 &  & 0.141 & yes \\
challenger & text\_lr+cal & 3797 &  &  &  & 0.0828 / 0.0832 / 0.0553 & 0.005 (0.000, 0.009) &  / 0.823 &  & 0.141 & yes \\
\bottomrule
\end{tabular}
}
\end{table*}

\subsection{Calibration by precision}
Table~\ref{tab:calibration-upworthy} is the diagnostic the cautionary run lacked. By decile of planned exposure, P1's mean $\psig$ runs from 0.076 to 0.152 while the realized significant-positive rate runs from 0.065 to 0.110; the forecaster rises with exposure somewhat faster than the data, which is the whole of its miscalibration, and 90 percent coverage stays between 0.83 and 0.94 across deciles. P2's columns track P1's to within 0.015 in the first eight deciles; in the top two, where it admits precedents most often, its mean $\psig$ of 0.136 and 0.167 overshoots event rates of 0.081 and 0.110, and its 90 percent coverage runs from 0.83 to 0.95.

\begin{table}[t]\centering\scriptsize
\caption{Calibration by decile of planned exposure per arm (headline archive).}
\label{tab:calibration-upworthy}
\resizebox{\columnwidth}{!}{% generated by paper/build_tables.py from backtest-upworthy-20261106-2dc461f52ea4.json
\begin{tabular}{rrrrrrr}
\toprule
decile & \multicolumn{3}{c}{P1: mean $\psig$ / event rate / cov.\ 90} & \multicolumn{3}{c}{P2: mean $\psig$ / event rate / cov.\ 90} \\
\midrule
1 & 0.074 & 0.064 & 0.86 & 0.072 & 0.064 & 0.86 \\
2 & 0.093 & 0.077 & 0.87 & 0.091 & 0.077 & 0.87 \\
3 & 0.092 & 0.074 & 0.87 & 0.088 & 0.074 & 0.87 \\
4 & 0.085 & 0.105 & 0.88 & 0.084 & 0.105 & 0.88 \\
5 & 0.085 & 0.086 & 0.89 & 0.084 & 0.086 & 0.89 \\
6 & 0.100 & 0.085 & 0.89 & 0.098 & 0.085 & 0.89 \\
7 & 0.106 & 0.101 & 0.91 & 0.104 & 0.101 & 0.90 \\
8 & 0.106 & 0.089 & 0.89 & 0.105 & 0.089 & 0.89 \\
9 & 0.119 & 0.084 & 0.94 & 0.119 & 0.084 & 0.94 \\
10 & 0.133 & 0.089 & 0.95 & 0.132 & 0.089 & 0.95 \\
\bottomrule
\end{tabular}
}
\end{table}

\subsection{Decision value}
Table~\ref{tab:decision-upworthy} answers the question the launching team actually asks. Choosing a quarter of each period's tests at random captures 0.250 of the realized positive lift; choosing by P1's $\pwin$ captures 0.250 and by P2's 0.255; an oracle that knows the outcomes captures 0.686. The ledger, on this archive, does not tell the team which headline tests to run; it tells them how their results will be distributed. We report this separately from calibration on purpose: a calibrated forecaster with no discrimination has zero decision value, and the two properties should never be summed into one number.

\begin{table}[t]\centering\small
\caption{Lift captured by prioritization policies at $K$ equal to a quarter of each period's tests (headline archive).}
\label{tab:decision-upworthy}
\resizebox{\columnwidth}{!}{% generated by paper/build_tables.py from backtest-upworthy-20261106-2dc461f52ea4.json (tables.decision_value; each mean is over the 7 per-origin rows of the policy and excludes the pooled origin=ALL row)
\begin{tabular}{llrr}
\toprule
policy & model & captured fraction (mean over origins) & regret (mean) \\
\midrule
random & all & 0.251 & 60.5 \\
p\_win & GBM & 0.278 & 56.8 \\
p\_win & P0 & 0.251 & 60.6 \\
p\_win & P1 & 0.249 & 60.7 \\
p\_win & P2 & 0.289 & 57.0 \\
p\_win & local\_eb & 0.291 & 56.5 \\
p\_win & text\_lr & 0.324 & 52.6 \\
arm\_count & all & 0.311 & 52.3 \\
oracle & all & 0.686 & 0.000 \\
random\_matched & GBM & 0.258 & 59.2 \\
random\_matched & P0 & 0.250 & 60.3 \\
random\_matched & P1 & 0.251 & 60.2 \\
random\_matched & P2 & 0.258 & 59.5 \\
random\_matched & local\_eb & 0.267 & 59.1 \\
random\_matched & text\_lr & 0.265 & 59.0 \\
\bottomrule
\end{tabular}
}
\end{table}

\subsection{Learning curve and gates}
Subsampling the ledger visible at the final origin (Table~\ref{tab:learning-upworthy}) shows P1 reaching a Brier skill near 0.03 on the final quarter's held-out tests with intervals above zero from five percent of the ledger onward, while the gradient-boosted challenger stays negative at every fraction and the text classifier stays at zero. The gate sizes from the power simulation are 175 and 375 rows for skills of 0.20 and 0.10 at the release target of 0.90 (125 and 300 at 0.80, Table~\ref{tab:gates}); decidability is reported by tests, not rows, because contrasts cluster within tests.

\begin{table*}[t]\centering\scriptsize
\caption{Brier skill against the base rate by fraction of the visible ledger at the final origin (headline archive).}
\label{tab:learning-upworthy}
\resizebox{\textwidth}{!}{% generated by paper/build_tables.py from learning-curve-upworthy-20261106-2dc461f52ea4.json (held-out rows 1143, seeds 20, gates {'n_200': 175, 'n_400': 250, 'by_target': {'0.80': {'n_200': 125, 'n_400': 200}, '0.90': {'n_200': 175, 'n_400': 250}}, 'decision_target': 0.9})
\begin{tabular}{rrrrrrr}
\toprule
fraction of ledger & GBM skill & P0 skill & P1 skill & P2 skill & local\_eb skill & text\_lr skill \\
\midrule
0.01 & -0.058 (-0.094, -0.033) & -0.042 (-0.053, -0.034) & 0.009 (-0.004, 0.027) & 0.009 (-0.005, 0.027) & -0.011 (-0.020, 0.000) & -0.025 (-0.028, -0.021) \\
0.02 & -0.094 (-0.145, -0.058) & -0.027 (-0.034, -0.022) & 0.012 (-0.002, 0.030) & 0.011 (-0.002, 0.028) & -0.014 (-0.024, -0.003) & -0.014 (-0.019, -0.008) \\
0.05 & -0.088 (-0.143, -0.048) & -0.019 (-0.027, -0.014) & 0.018 (0.004, 0.037) & 0.018 (0.004, 0.037) & -0.005 (-0.019, 0.008) & -0.008 (-0.013, 0.000) \\
0.10 & -0.077 (-0.128, -0.039) & -0.008 (-0.013, -0.005) & 0.020 (0.008, 0.037) & 0.020 (0.007, 0.038) & 0.003 (-0.007, 0.017) & -0.006 (-0.015, 0.003) \\
0.25 & -0.054 (-0.096, -0.021) & -0.003 (-0.005, -0.002) & 0.024 (0.011, 0.043) & 0.025 (0.010, 0.046) & 0.012 (-0.003, 0.030) & 0.003 (-0.008, 0.014) \\
0.50 & -0.039 (-0.079, -0.010) & 0.000 (-0.002, 0.002) & 0.024 (0.011, 0.042) & 0.027 (0.012, 0.049) & 0.019 (0.004, 0.039) & 0.006 (-0.010, 0.021) \\
1.00 & -0.041 (-0.079, -0.015) & 0.000 (0.000, 0.000) & 0.024 (0.011, 0.042) & 0.029 (0.014, 0.051) & 0.013 (-0.011, 0.037) & 0.017 (-0.007, 0.037) \\
\bottomrule
\end{tabular}
}
\end{table*}

\begin{table}[t]\centering\small
\caption{Gate sizes from the power simulation.}
\label{tab:gates}
\resizebox{\columnwidth}{!}{% generated by paper/build_tables.py from /Volumes/Crucial X10/mootcourt-data/reports/gates.json (base rates ['0.09', '0.14', '0.30'], interval 0.9 by test, 3.5 rows per test, reps 1000, bootstrap 500, seed 20261102)
\begin{tabular}{rlrr}
\toprule
base rate & power target & tests for skill 0.20 & tests for skill 0.10 \\
\midrule
0.09 & 0.80 & 75 & 125 \\
0.09 & 0.90 & 100 & 175 \\
0.14 & 0.80 & 50 & 100 \\
0.14 & 0.90 & 75 & 125 \\
0.30 & 0.80 & 50 & 75 \\
0.30 & 0.90 & 50 & 100 \\
\bottomrule
\end{tabular}
}
\end{table}

\subsection{Abstention}
Table~\ref{tab:abstention-upworthy} counts the triggers. The reliability trigger fires only at the burn-in origin, where no earlier report exists; support fires in the smallest cell; novelty fires on a few hundred rows; the decision-instability trigger never fires on this archive. The pooled abstention rate is 8.7 percent. \todo{Compare skill on issued versus forced rows once the ablation that removes one trigger at a time is in the report.}

\begin{table}[t]\centering\small
\caption{Abstention triggers (headline archive).}
\label{tab:abstention-upworthy}
\resizebox{\columnwidth}{!}{% generated by paper/build_tables.py from backtest-upworthy-20261103-4bfb354b87e3.json
\begin{tabular}{lrrrrrr}
\toprule
cell & n & rate & support & novelty & reliability & instability \\
\midrule
ALL & 13247 & 0.141 & 304 & 1121 & 781 & 0 \\
both & 471 & 0.221 & 68 & 36 & 68 & 0 \\
headline\_only & 7417 & 0.115 & 0 & 375 & 477 & 0 \\
image\_only & 5359 & 0.171 & 236 & 710 & 236 & 0 \\
\bottomrule
\end{tabular}
}
\end{table}

\subsection{The deployment-shaped archive}
Table~\ref{tab:main-asos} reports the same protocol on 68 experiments. Every metric cell is undecidable at both gate targets, the abstention rate is 0.552 because the support trigger fires in most cells, and P1's Brier skill is minus 0.107 with an interval that reaches zero. The reader should take exactly one thing from this table: the protocol says \emph{undecidable} when the ledger is too small, instead of reporting a number that looks like a result. The nominal calendar of this archive means its rolling origin tests the mechanics, not drift.

\begin{table*}[t]\centering\scriptsize
\caption{Deployment-shaped archive, pooled over origins and metric cells.}
\label{tab:main-asos}
\resizebox{\textwidth}{!}{% generated by paper/build_tables.py from /Volumes/Crucial X10/mootcourt-data/reports/backtest-asos-20261103-6bb5a15c5ff0.json (source asos, snapshot 20261103-6bb5a15c5ff0, model 548c646b9cb1, 337 scored rows over 7 origins; gate decision target 0.9)
\begin{tabular}{llrrrrrrrr}
\toprule
kind & model & tests & log score & CRPS & Brier skill (90\%) & dir.\ acc.\ / ceiling & cov.\ 80 / 90 & abst. & decid. \\
\midrule
ladder & P0 & 68 & 4.383 & 0.00 & -0.443 (-0.836, -0.185) & 0.499 / 0.806 & 0.91 / 0.93 & 0.448 & no \\
ladder & P1 & 68 & 4.610 & 0.00 & -0.107 (-0.237, -0.002) & 0.522 / 0.806 & 0.78 / 0.86 & 0.448 & no \\
ladder & P2 & 68 & 4.610 & 0.00 & -0.107 (-0.237, -0.002) & 0.522 / 0.806 & 0.78 / 0.86 & 0.448 & no \\
baseline & always\_reference & 68 &  &  &  & 0.496 / 0.806 &  & 0.448 & no \\
baseline & always\_variant & 68 &  &  &  & 0.504 / 0.806 &  & 0.448 & no \\
challenger & GBM & 68 & 3.629 & 0.00 & -0.714 (-1.259, -0.377) & 0.493 / 0.806 & 0.99 / 1.00 & 0.448 & no \\
challenger & local\_eb & 68 & 4.428 & 0.00 & -0.383 (-0.736, -0.138) & 0.487 / 0.806 & 0.89 / 0.90 & 0.448 & no \\
\bottomrule
\end{tabular}
}
\end{table*}

% Results, second half: abstention, drift, the deployment-shaped archive and the synthetic archive. Source comments follow the convention of results.tex: a "% src:" line names the report and key behind the one line that follows it.
\subsection{Abstention}
\label{sec:abstention-results}
A card is the frozen forecast a team sees before launch. It abstains when one of four triggers fires, and an abstained card keeps the forecast it withheld, which is scored but never shown. Table~\ref{tab:fmt-abstention-triggers} condenses the trigger counts of Table~\ref{tab:abstention-upworthy} (Appendix~\ref{app:tables}), and Table~\ref{tab:fmt-abstention-composition} gives the rate after the burn-in origin, the first origin, at which no reliability report exists yet.
% src: $RUN6/paper-extras-upworthy-20261106-2dc461f52ea4.json abstention_composition.P2 abstained_rows 2722, rows 10018 (0.272); $RUN6/backtest-upworthy-20261106-2dc461f52ea4.json tables.origins origin=2014-01-01 n_scored 1236, abstention_rate 1.0
The rate after burn-in is 27.2 percent, and that includes January 2014, where every row abstains. This is a consequence of the rule on the randomization window: the gap in results leaves the reliability reports stale at that origin. Without that origin the rate is lower, and by cell the novelty trigger makes it highest in the image-only cell.

\begin{table}[t]\centering\scriptsize
% table: tab:abstention-upworthy; tables/fmt-abstention-triggers.tex
\caption{The abstention triggers on the headline archive, from Table~\ref{tab:abstention-upworthy} (Appendix~\ref{app:tables}), all cells. Rows: the trigger, then any trigger. Columns: the rows withheld, their share of all rows, and when the trigger fires. A row can trip more than one trigger, so the counts do not add to the total. These are counts, so no row carries an interval.}
\label{tab:fmt-abstention-triggers}
% generated by paper/build_tables.py from the report keys named above each row (results_b)
\begin{tabular}{lrrp{0.5\columnwidth}}
\toprule
trigger & rows & share & when it fires \\
\midrule
% src: $RUN6/backtest-upworthy-20261103-4bfb354b87e3.json tables.abstention[cell=ALL].reliability
reliability & 781 &  & on the rows of the burn-in origin and never after \\
% src: $RUN6/backtest-upworthy-20261103-4bfb354b87e3.json tables.abstention[cell=ALL].support; $RUN6/backtest-upworthy-20261103-4bfb354b87e3.json tables.abstention[cell=both].support; $RUN6/backtest-upworthy-20261103-4bfb354b87e3.json tables.abstention[cell=headline_only].support; $RUN6/backtest-upworthy-20261103-4bfb354b87e3.json tables.abstention[cell=image_only].support
support & 304 &  & on rows of the burn-in origin, in the two cells with fewer than 30 visible rows \\
% src: $RUN6/backtest-upworthy-20261103-4bfb354b87e3.json tables.abstention[cell=ALL].novelty; $RUN6/backtest-upworthy-20261103-4bfb354b87e3.json tables.abstention[cell=ALL].novelty; $RUN6/backtest-upworthy-20261103-4bfb354b87e3.json tables.abstention[cell=ALL].n
novelty & 1,121 & 8.5\% & the only trigger that fires after burn-in \\
% src: $RUN6/backtest-upworthy-20261103-4bfb354b87e3.json tables.abstention[cell=ALL].instability; $RUN6/backtest-upworthy-20261103-4bfb354b87e3.json tables.p2_selection_by_cell[cell=both].share_lambda_zero; $RUN6/backtest-upworthy-20261103-4bfb354b87e3.json tables.p2_selection_by_cell[cell=both].origins; $RUN6/backtest-upworthy-20261103-4bfb354b87e3.json tables.p2_selection_by_cell[cell=headline_only].share_lambda_zero; $RUN6/backtest-upworthy-20261103-4bfb354b87e3.json tables.p2_selection_by_cell[cell=headline_only].origins; $RUN6/backtest-upworthy-20261103-4bfb354b87e3.json tables.p2_selection_by_cell[cell=image_only].share_lambda_zero; $RUN6/backtest-upworthy-20261103-4bfb354b87e3.json tables.p2_selection_by_cell[cell=image_only].origins
decision instability & 0 &  & never, being off by construction in the 19 pairs where the precedent weight is zero \\
\midrule
% src: $RUN6/paper-extras-upworthy-20261103-4bfb354b87e3.json issued_row_brier_skill.rungs.P2.abstained.n; $RUN6/backtest-upworthy-20261103-4bfb354b87e3.json tables.abstention[cell=ALL].rate
any trigger & 1,874 & 14.1\% & the pooled abstention rate \\
\bottomrule
\end{tabular}

\end{table}

Issued rows do not have the higher skill. With the burn-in origin included the abstained rows score above the issued rows on Brier skill (Table~\ref{tab:fmt-abstention-rows}), and the abstained rows carry that origin and a higher event rate. With it excluded the two groups do not differ measurably, pooled or within the strata that hold both groups, whose interval spans zero (Table~\ref{tab:fmt-abstention-composition}).

\begin{table}[t]\centering\scriptsize
% src: $RUN6/paper-extras-upworthy-20261106-2dc461f52ea4.json issued_row_brier_skill bootstrap_resamples 500, level 0.9; $RUN6/backtest-upworthy-20261106-2dc461f52ea4.json tables.aggregate[cell=ALL,rung=P2] logscore, logscore_nonabstained, n, n_nonabstained; tables/fmt-abstention-rows.tex
\caption{Issued against abstained against all rows, P2 on the headline archive, burn-in origin included. The abstained rows are scored on the forecast the rule withheld. Columns: rows, Brier skill against the row set's own trailing base rate, pooled log score per row (the abstained figure derived from the other two, so it has no interval), and event rate.}
\label{tab:fmt-abstention-rows}
% generated by paper/build_tables.py from the report keys named above each row (results_b)
\begin{tabular}{lrrrr}
\toprule
row set & rows & Brier skill (90\%) & log score & event rate \\
\midrule
% src: $RUN6/paper-extras-upworthy-20261103-4bfb354b87e3.json issued_row_brier_skill.rungs.P2.issued.n; $RUN6/paper-extras-upworthy-20261103-4bfb354b87e3.json issued_row_brier_skill.rungs.P2.issued.brier_skill; $RUN6/paper-extras-upworthy-20261103-4bfb354b87e3.json issued_row_brier_skill.rungs.P2.issued.brier_skill_lo; $RUN6/paper-extras-upworthy-20261103-4bfb354b87e3.json issued_row_brier_skill.rungs.P2.issued.brier_skill_hi; $RUN6/backtest-upworthy-20261103-4bfb354b87e3.json tables.aggregate[cell=ALL,rung=P2].logscore_nonabstained; $RUN6/paper-extras-upworthy-20261103-4bfb354b87e3.json issued_row_brier_skill.rungs.P2.issued.event_rate
issued & 11,373 & -0.002 (-0.006, 0.002) & -0.854 & 0.09 \\
% src: $RUN6/paper-extras-upworthy-20261103-4bfb354b87e3.json issued_row_brier_skill.rungs.P2.abstained.n; $RUN6/paper-extras-upworthy-20261103-4bfb354b87e3.json issued_row_brier_skill.rungs.P2.abstained.brier_skill; $RUN6/paper-extras-upworthy-20261103-4bfb354b87e3.json issued_row_brier_skill.rungs.P2.abstained.brier_skill_lo; $RUN6/paper-extras-upworthy-20261103-4bfb354b87e3.json issued_row_brier_skill.rungs.P2.abstained.brier_skill_hi; $RUN6/paper-extras-upworthy-20261103-4bfb354b87e3.json issued_row_brier_skill.rungs.P2.abstained.event_rate
abstained & 1,874 & 0.020 (0.008, 0.034) &  & 0.11 \\
% src: $RUN6/paper-extras-upworthy-20261103-4bfb354b87e3.json issued_row_brier_skill.rungs.P2.all.n; $RUN6/paper-extras-upworthy-20261103-4bfb354b87e3.json issued_row_brier_skill.rungs.P2.all.brier_skill; $RUN6/paper-extras-upworthy-20261103-4bfb354b87e3.json issued_row_brier_skill.rungs.P2.all.brier_skill_lo; $RUN6/paper-extras-upworthy-20261103-4bfb354b87e3.json issued_row_brier_skill.rungs.P2.all.brier_skill_hi; $RUN6/backtest-upworthy-20261103-4bfb354b87e3.json tables.aggregate[cell=ALL,rung=P2].logscore; $RUN6/paper-extras-upworthy-20261103-4bfb354b87e3.json issued_row_brier_skill.rungs.P2.all.event_rate
all rows & 13,247 & 0.002 (-0.002, 0.005) & -0.856 & 0.09 \\
\bottomrule
\end{tabular}

\end{table}

\begin{table}[t]\centering\scriptsize
% src: $RUN6/paper-extras-upworthy-20261106-2dc461f52ea4.json abstention_composition (burn_in_origin 2013-04-01, bootstrap_resamples 500); $RUN6/backtest-upworthy-20261106-2dc461f52ea4.json tables.origins; tables/fmt-abstention-composition.tex
\caption{Abstained against issued rows on the headline archive with the burn-in origin excluded. Brier skill is one minus the group's summed Brier over its summed base Brier, the base being the trailing rate of the row's cell at its origin. The within row reads the difference only inside strata that hold both groups; the rest, January 2014's among them, drop out.}
\label{tab:fmt-abstention-composition}
% generated by paper/build_tables.py from the report keys named above each row (experiments)
\begin{tabular}{lrr}
\toprule
 & P1 & P2 \\
\midrule
% src: $RUN6/paper-extras-upworthy-20261106-2dc461f52ea4.json abstention_composition.P1.issued.rows; $RUN6/paper-extras-upworthy-20261106-2dc461f52ea4.json abstention_composition.P1.issued.tests; $RUN6/paper-extras-upworthy-20261106-2dc461f52ea4.json abstention_composition.P2.issued.rows; $RUN6/paper-extras-upworthy-20261106-2dc461f52ea4.json abstention_composition.P2.issued.tests
issued, rows / tests & 7,296 / 2,021 & 7,296 / 2,021 \\
% src: $RUN6/paper-extras-upworthy-20261106-2dc461f52ea4.json abstention_composition.P1.issued.skill; $RUN6/paper-extras-upworthy-20261106-2dc461f52ea4.json abstention_composition.P1.issued.lo; $RUN6/paper-extras-upworthy-20261106-2dc461f52ea4.json abstention_composition.P1.issued.hi; $RUN6/paper-extras-upworthy-20261106-2dc461f52ea4.json abstention_composition.P2.issued.skill; $RUN6/paper-extras-upworthy-20261106-2dc461f52ea4.json abstention_composition.P2.issued.lo; $RUN6/paper-extras-upworthy-20261106-2dc461f52ea4.json abstention_composition.P2.issued.hi
issued, Brier skill (90\%) & -0.0004 (-0.0040, 0.0033) & 0.0004 (-0.0035, 0.0045) \\
% src: $RUN6/paper-extras-upworthy-20261106-2dc461f52ea4.json abstention_composition.P1.abstained.rows; $RUN6/paper-extras-upworthy-20261106-2dc461f52ea4.json abstention_composition.P1.abstained.tests; $RUN6/paper-extras-upworthy-20261106-2dc461f52ea4.json abstention_composition.P2.abstained.rows; $RUN6/paper-extras-upworthy-20261106-2dc461f52ea4.json abstention_composition.P2.abstained.tests
abstained, rows / tests & 2,722 / 835 & 2,722 / 835 \\
% src: $RUN6/paper-extras-upworthy-20261106-2dc461f52ea4.json abstention_composition.P1.abstained.skill; $RUN6/paper-extras-upworthy-20261106-2dc461f52ea4.json abstention_composition.P1.abstained.lo; $RUN6/paper-extras-upworthy-20261106-2dc461f52ea4.json abstention_composition.P1.abstained.hi; $RUN6/paper-extras-upworthy-20261106-2dc461f52ea4.json abstention_composition.P2.abstained.skill; $RUN6/paper-extras-upworthy-20261106-2dc461f52ea4.json abstention_composition.P2.abstained.lo; $RUN6/paper-extras-upworthy-20261106-2dc461f52ea4.json abstention_composition.P2.abstained.hi
abstained, Brier skill (90\%) & -0.0017 (-0.0066, 0.0031) & -0.0013 (-0.0060, 0.0041) \\
\midrule
% src: $RUN6/paper-extras-upworthy-20261106-2dc461f52ea4.json abstention_composition.P1.pooled_issued_minus_abstained_skill; $RUN6/paper-extras-upworthy-20261106-2dc461f52ea4.json abstention_composition.P2.pooled_issued_minus_abstained_skill
issued minus abstained, pooled & 0.0013 & 0.0017 \\
% src: $RUN6/paper-extras-upworthy-20261106-2dc461f52ea4.json abstention_composition.P1.stratified_issued_minus_abstained_skill.estimate; $RUN6/paper-extras-upworthy-20261106-2dc461f52ea4.json abstention_composition.P1.stratified_issued_minus_abstained_skill.lo; $RUN6/paper-extras-upworthy-20261106-2dc461f52ea4.json abstention_composition.P1.stratified_issued_minus_abstained_skill.hi; $RUN6/paper-extras-upworthy-20261106-2dc461f52ea4.json abstention_composition.P2.stratified_issued_minus_abstained_skill.estimate; $RUN6/paper-extras-upworthy-20261106-2dc461f52ea4.json abstention_composition.P2.stratified_issued_minus_abstained_skill.lo; $RUN6/paper-extras-upworthy-20261106-2dc461f52ea4.json abstention_composition.P2.stratified_issued_minus_abstained_skill.hi
issued minus abstained, within (90\%) & -0.0010 (-0.0384, 0.0125) & -0.0023 (-0.0459, 0.0097) \\
% src: $RUN6/paper-extras-upworthy-20261106-2dc461f52ea4.json abstention_composition.P1.stratified_issued_minus_abstained_skill.strata_with_both_groups; $RUN6/paper-extras-upworthy-20261106-2dc461f52ea4.json abstention_composition.P1.stratified_issued_minus_abstained_skill.strata_total; $RUN6/paper-extras-upworthy-20261106-2dc461f52ea4.json abstention_composition.P2.stratified_issued_minus_abstained_skill.strata_with_both_groups; $RUN6/paper-extras-upworthy-20261106-2dc461f52ea4.json abstention_composition.P2.stratified_issued_minus_abstained_skill.strata_total
strata holding both groups & 13 of 18 & 13 of 18 \\
\midrule
% src: $RUN6/paper-extras-upworthy-20261106-2dc461f52ea4.json abstention_composition.P1.abstention_rate_by_cell.headline_only.rate; $RUN6/paper-extras-upworthy-20261106-2dc461f52ea4.json abstention_composition.P1.abstention_rate_by_cell.headline_only.rows; $RUN6/paper-extras-upworthy-20261106-2dc461f52ea4.json abstention_composition.P2.abstention_rate_by_cell.headline_only.rate; $RUN6/paper-extras-upworthy-20261106-2dc461f52ea4.json abstention_composition.P2.abstention_rate_by_cell.headline_only.rows
abstention rate, headline-only, of rows & 0.160 of 5,567 & 0.160 of 5,567 \\
% src: $RUN6/paper-extras-upworthy-20261106-2dc461f52ea4.json abstention_composition.P1.abstention_rate_by_cell.image_only.rate; $RUN6/paper-extras-upworthy-20261106-2dc461f52ea4.json abstention_composition.P1.abstention_rate_by_cell.image_only.rows; $RUN6/paper-extras-upworthy-20261106-2dc461f52ea4.json abstention_composition.P2.abstention_rate_by_cell.image_only.rate; $RUN6/paper-extras-upworthy-20261106-2dc461f52ea4.json abstention_composition.P2.abstention_rate_by_cell.image_only.rows
abstention rate, image-only, of rows & 0.432 of 4,148 & 0.432 of 4,148 \\
% src: $RUN6/paper-extras-upworthy-20261106-2dc461f52ea4.json abstention_composition.P1.abstention_rate_by_cell.both.rate; $RUN6/paper-extras-upworthy-20261106-2dc461f52ea4.json abstention_composition.P1.abstention_rate_by_cell.both.rows; $RUN6/paper-extras-upworthy-20261106-2dc461f52ea4.json abstention_composition.P2.abstention_rate_by_cell.both.rate; $RUN6/paper-extras-upworthy-20261106-2dc461f52ea4.json abstention_composition.P2.abstention_rate_by_cell.both.rows
abstention rate, smallest, of rows & 0.135 of 303 & 0.135 of 303 \\
\bottomrule
\end{tabular}

\end{table}

% src: $RUN6/backtest-upworthy-20261106-2dc461f52ea4.json tables.origins (abstention rate 1.0 at 2014-01-01); tables/fmt-abstention-rows.tex; tables/fmt-abstention-released.tex
The log score reads differently, and it is not stratified. Issued rows score above abstained rows, and above each group of withheld rows scored with its trigger switched off (Table~\ref{tab:fmt-abstention-released}, condensed from Table~\ref{tab:trigger-released}, Appendix~\ref{app:tables}). These row sets differ in origin and in cell, and no report compares them within strata, so this is not evidence that abstention selected the worse forecasts. So issued rows are better on the log score and not on Brier skill, and abstention is not shown to filter quality on this archive. We keep it as a label of evidence, a purpose we revised after seeing this result. Direction accuracy points the other way for the novelty trigger, the only one that fires after January 2014: its released rows are above the issued rows on direction accuracy and below them on the log score. Support and instability release no rows.

\begin{table}[t]\centering\scriptsize
% src: $RUN6/ablations/backtest-upworthy-20261106-2dc461f52ea4--drop-trigger-novelty.json tables.released_rows cell=ALL rung=P2 n 1486; tables/trigger-released-upworthy.tex; tables/fmt-abstention-released.tex
\caption{The rows a trigger had withheld, scored with that trigger switched off, P2 on the headline archive, from Table~\ref{tab:trigger-released} (Appendix~\ref{app:tables}), with the main run's issued rows for comparison (Table~\ref{tab:trigger-ablation}). Support and instability release no rows. Columns: released rows, pooled log score per row and direction accuracy.}
\label{tab:fmt-abstention-released}
\resizebox{\columnwidth}{!}{% generated by paper/build_tables.py from the report keys named above each row (results_b)
\begin{tabular}{lrrr}
\toprule
rows & count & log score (90\%) & direction accuracy (90\%) \\
\midrule
% src: $RUN6/ablations/backtest-upworthy-20261103-4bfb354b87e3--drop-trigger-novelty.json tables.released_rows[cell=ALL,rung=P2].n; $RUN6/ablations/backtest-upworthy-20261103-4bfb354b87e3--drop-trigger-novelty.json tables.released_rows[cell=ALL,rung=P2].logscore; $RUN6/ablations/backtest-upworthy-20261103-4bfb354b87e3--drop-trigger-novelty.json tables.released_rows[cell=ALL,rung=P2].logscore_lo; $RUN6/ablations/backtest-upworthy-20261103-4bfb354b87e3--drop-trigger-novelty.json tables.released_rows[cell=ALL,rung=P2].logscore_hi; $RUN6/ablations/backtest-upworthy-20261103-4bfb354b87e3--drop-trigger-novelty.json tables.released_rows[cell=ALL,rung=P2].direction_accuracy; $RUN6/ablations/backtest-upworthy-20261103-4bfb354b87e3--drop-trigger-novelty.json tables.released_rows[cell=ALL,rung=P2].direction_accuracy_lo; $RUN6/ablations/backtest-upworthy-20261103-4bfb354b87e3--drop-trigger-novelty.json tables.released_rows[cell=ALL,rung=P2].direction_accuracy_hi
released by novelty & 1,093 & -0.824 (-0.895, -0.750) & 0.560 (0.528, 0.590) \\
% src: $RUN6/ablations/backtest-upworthy-20261103-4bfb354b87e3--drop-trigger-reliability.json tables.released_rows[cell=ALL,rung=P2].n; $RUN6/ablations/backtest-upworthy-20261103-4bfb354b87e3--drop-trigger-reliability.json tables.released_rows[cell=ALL,rung=P2].logscore; $RUN6/ablations/backtest-upworthy-20261103-4bfb354b87e3--drop-trigger-reliability.json tables.released_rows[cell=ALL,rung=P2].logscore_lo; $RUN6/ablations/backtest-upworthy-20261103-4bfb354b87e3--drop-trigger-reliability.json tables.released_rows[cell=ALL,rung=P2].logscore_hi; $RUN6/ablations/backtest-upworthy-20261103-4bfb354b87e3--drop-trigger-reliability.json tables.released_rows[cell=ALL,rung=P2].direction_accuracy; $RUN6/ablations/backtest-upworthy-20261103-4bfb354b87e3--drop-trigger-reliability.json tables.released_rows[cell=ALL,rung=P2].direction_accuracy_lo; $RUN6/ablations/backtest-upworthy-20261103-4bfb354b87e3--drop-trigger-reliability.json tables.released_rows[cell=ALL,rung=P2].direction_accuracy_hi
burn-in rows released by reliability alone & 477 & -0.956 (-1.065, -0.852) & 0.516 (0.475, 0.554) \\
\midrule
% src: $RUN6/backtest-upworthy-20261103-4bfb354b87e3.json tables.aggregate[cell=ALL,rung=P2].n_nonabstained; $RUN6/backtest-upworthy-20261103-4bfb354b87e3.json tables.aggregate[cell=ALL,rung=P2].logscore_nonabstained; $RUN6/backtest-upworthy-20261103-4bfb354b87e3.json tables.aggregate[cell=ALL,rung=P2].logscore_nonabstained_lo; $RUN6/backtest-upworthy-20261103-4bfb354b87e3.json tables.aggregate[cell=ALL,rung=P2].logscore_nonabstained_hi; $RUN6/backtest-upworthy-20261103-4bfb354b87e3.json tables.aggregate[cell=ALL,rung=P2].direction_accuracy_nonabstained; $RUN6/backtest-upworthy-20261103-4bfb354b87e3.json tables.aggregate[cell=ALL,rung=P2].direction_accuracy_nonabstained_lo; $RUN6/backtest-upworthy-20261103-4bfb354b87e3.json tables.aggregate[cell=ALL,rung=P2].direction_accuracy_nonabstained_hi
issued rows, main run & 11,373 & -0.854 (-0.877, -0.832) & 0.582 (0.572, 0.592) \\
\bottomrule
\end{tabular}
}
\end{table}

\subsection{Drift and the weekly replay}
\label{sec:drift-results}
% src: $RUN6/backtest-upworthy-20261106-2dc461f52ea4.json tables.score_by_year cell=ALL rung=P1 brier_skill_lo, brier_skill_hi, adjusted_lo, adjusted_hi, event_rate, years 2013, 2014, 2015
Table~\ref{tab:fmt-drift-year} condenses Table~\ref{tab:drift} (Appendix~\ref{app:tables}) by the calendar year in which the test started. P1's skill shows no trend: its unadjusted interval lies above zero in 2013 and 2015 and below zero in 2014, and every adjusted interval spans zero. The significant-positive rate falls, a shift that P0's trailing base rate lags and that the drift rule, which watches the log score against P0 alone, cannot see.

\begin{table}[t]\centering\scriptsize
% src: $RUN6/backtest-upworthy-20261106-2dc461f52ea4.json tables.score_by_year cell=ALL rung=P1 family_size 156 (tables/drift-upworthy.tex); tables/fmt-drift-year.tex
\caption{The prior by calendar year of the test's start on the headline archive, from Table~\ref{tab:drift} (Appendix~\ref{app:tables}). Columns: tests, P1's Brier skill, that interval adjusted for the table's 156 intervals, and the significant-positive rate. The tests of 2013 are the burn-in origin's, with the excluded window removed.}
\label{tab:fmt-drift-year}
% generated by paper/build_tables.py from the report keys named above each row (results_b)
\begin{tabular}{rrrrr}
\toprule
year & tests & P1 Brier skill (90\%) & adjusted interval & significant-positive rate \\
\midrule
% src: $RUN6/backtest-upworthy-20261103-4bfb354b87e3.json tables.score_by_year[cell=ALL,rung=P1,year=2013].year; $RUN6/backtest-upworthy-20261103-4bfb354b87e3.json tables.score_by_year[cell=ALL,rung=P1,year=2013].n_tests; $RUN6/backtest-upworthy-20261103-4bfb354b87e3.json tables.score_by_year[cell=ALL,rung=P1,year=2013].brier_skill; $RUN6/backtest-upworthy-20261103-4bfb354b87e3.json tables.score_by_year[cell=ALL,rung=P1,year=2013].brier_skill_lo; $RUN6/backtest-upworthy-20261103-4bfb354b87e3.json tables.score_by_year[cell=ALL,rung=P1,year=2013].brier_skill_hi; $RUN6/backtest-upworthy-20261103-4bfb354b87e3.json tables.score_by_year[cell=ALL,rung=P1,year=2013].event_rate
2013 & 1,068 & 0.001 (-0.008, 0.011) &  & 0.117 \\
% src: $RUN6/backtest-upworthy-20261103-4bfb354b87e3.json tables.score_by_year[cell=ALL,rung=P1,year=2014].year; $RUN6/backtest-upworthy-20261103-4bfb354b87e3.json tables.score_by_year[cell=ALL,rung=P1,year=2014].n_tests; $RUN6/backtest-upworthy-20261103-4bfb354b87e3.json tables.score_by_year[cell=ALL,rung=P1,year=2014].brier_skill; $RUN6/backtest-upworthy-20261103-4bfb354b87e3.json tables.score_by_year[cell=ALL,rung=P1,year=2014].brier_skill_lo; $RUN6/backtest-upworthy-20261103-4bfb354b87e3.json tables.score_by_year[cell=ALL,rung=P1,year=2014].brier_skill_hi; $RUN6/backtest-upworthy-20261103-4bfb354b87e3.json tables.score_by_year[cell=ALL,rung=P1,year=2014].event_rate
2014 & 2,328 & -0.002 (-0.005, 0.002) &  & 0.087 \\
% src: $RUN6/backtest-upworthy-20261103-4bfb354b87e3.json tables.score_by_year[cell=ALL,rung=P1,year=2015].year; $RUN6/backtest-upworthy-20261103-4bfb354b87e3.json tables.score_by_year[cell=ALL,rung=P1,year=2015].n_tests; $RUN6/backtest-upworthy-20261103-4bfb354b87e3.json tables.score_by_year[cell=ALL,rung=P1,year=2015].brier_skill; $RUN6/backtest-upworthy-20261103-4bfb354b87e3.json tables.score_by_year[cell=ALL,rung=P1,year=2015].brier_skill_lo; $RUN6/backtest-upworthy-20261103-4bfb354b87e3.json tables.score_by_year[cell=ALL,rung=P1,year=2015].brier_skill_hi; $RUN6/backtest-upworthy-20261103-4bfb354b87e3.json tables.score_by_year[cell=ALL,rung=P1,year=2015].adjusted_lo; $RUN6/backtest-upworthy-20261103-4bfb354b87e3.json tables.score_by_year[cell=ALL,rung=P1,year=2015].adjusted_hi; $RUN6/backtest-upworthy-20261103-4bfb354b87e3.json tables.score_by_year[cell=ALL,rung=P1,year=2015].event_rate
2015 & 401 & 0.030 (0.016, 0.048) & (-0.004, 0.064) & 0.057 \\
\bottomrule
\end{tabular}

\end{table}

% src: $RUN6/paper-extras-upworthy-20261106-2dc461f52ea4.json within_origin_auc_by_year.years["2015"] P0, P1, P2, text_lr auc_within_origin, auc_within_origin_and_cell; tables/fmt-auc-by-year.tex
Table~\ref{tab:fmt-auc-by-year} asks whether the 2015 skill is better ranking. Within origin P1's AUC clears one half that year, but so does P0's, and P0 ranks only by cell. Within origin and cell P1's interval spans one half. So the order P1 gives in 2015 is largely between cells, and its skill that year is not shown to be better ranking inside them, while P2 and the text classifier clear one half within origin and cell in that year. The part of that skill due to P0 lagging a falling rate is not separated here.

\begin{table}[t]\centering\scriptsize
% src: $RUN6/paper-extras-upworthy-20261106-2dc461f52ea4.json within_origin_auc_by_year (500 resamples); tables/fmt-auc-by-year.tex
\caption{The event AUC by calendar year on the headline archive, within origin (where P0 still differs between cells) and within origin and cell (where P0 is one number).}
\label{tab:fmt-auc-by-year}
\resizebox{\columnwidth}{!}{% generated by paper/build_tables.py from the report keys named above each row (experiments)
\begin{tabular}{lrrrrr}
\toprule
year & tests & P1 AUC within & P2 AUC within & P1 Brier & P2 Brier \\
 & & origin (90\%) & origin (90\%) & skill & skill \\
\midrule
% src: $RUN6/paper-extras-upworthy-20261106-2dc461f52ea4.json within_origin_auc_by_year.years["2013"].P1.tests; $RUN6/paper-extras-upworthy-20261106-2dc461f52ea4.json within_origin_auc_by_year.years["2013"].P1.auc_within_origin.estimate; $RUN6/paper-extras-upworthy-20261106-2dc461f52ea4.json within_origin_auc_by_year.years["2013"].P1.auc_within_origin.lo; $RUN6/paper-extras-upworthy-20261106-2dc461f52ea4.json within_origin_auc_by_year.years["2013"].P1.auc_within_origin.hi; $RUN6/paper-extras-upworthy-20261106-2dc461f52ea4.json within_origin_auc_by_year.years["2013"].P2.auc_within_origin.estimate; $RUN6/paper-extras-upworthy-20261106-2dc461f52ea4.json within_origin_auc_by_year.years["2013"].P2.auc_within_origin.lo; $RUN6/paper-extras-upworthy-20261106-2dc461f52ea4.json within_origin_auc_by_year.years["2013"].P2.auc_within_origin.hi; $RUN6/paper-extras-upworthy-20261106-2dc461f52ea4.json within_origin_auc_by_year.years["2013"].P1.brier_skill; $RUN6/paper-extras-upworthy-20261106-2dc461f52ea4.json within_origin_auc_by_year.years["2013"].P2.brier_skill
2013 & 245 & 0.528 (0.459, 0.591) & 0.541 (0.468, 0.598) & 0.0334 & 0.0373 \\
% src: $RUN6/paper-extras-upworthy-20261106-2dc461f52ea4.json within_origin_auc_by_year.years["2014"].P1.tests; $RUN6/paper-extras-upworthy-20261106-2dc461f52ea4.json within_origin_auc_by_year.years["2014"].P1.auc_within_origin.estimate; $RUN6/paper-extras-upworthy-20261106-2dc461f52ea4.json within_origin_auc_by_year.years["2014"].P1.auc_within_origin.lo; $RUN6/paper-extras-upworthy-20261106-2dc461f52ea4.json within_origin_auc_by_year.years["2014"].P1.auc_within_origin.hi; $RUN6/paper-extras-upworthy-20261106-2dc461f52ea4.json within_origin_auc_by_year.years["2014"].P2.auc_within_origin.estimate; $RUN6/paper-extras-upworthy-20261106-2dc461f52ea4.json within_origin_auc_by_year.years["2014"].P2.auc_within_origin.lo; $RUN6/paper-extras-upworthy-20261106-2dc461f52ea4.json within_origin_auc_by_year.years["2014"].P2.auc_within_origin.hi; $RUN6/paper-extras-upworthy-20261106-2dc461f52ea4.json within_origin_auc_by_year.years["2014"].P1.brier_skill; $RUN6/paper-extras-upworthy-20261106-2dc461f52ea4.json within_origin_auc_by_year.years["2014"].P2.brier_skill
2014 & 2,303 & 0.541 (0.519, 0.563) & 0.545 (0.523, 0.566) & -0.0036 & -0.0035 \\
% src: $RUN6/paper-extras-upworthy-20261106-2dc461f52ea4.json within_origin_auc_by_year.years["2015"].P1.tests; $RUN6/paper-extras-upworthy-20261106-2dc461f52ea4.json within_origin_auc_by_year.years["2015"].P1.auc_within_origin.estimate; $RUN6/paper-extras-upworthy-20261106-2dc461f52ea4.json within_origin_auc_by_year.years["2015"].P1.auc_within_origin.lo; $RUN6/paper-extras-upworthy-20261106-2dc461f52ea4.json within_origin_auc_by_year.years["2015"].P1.auc_within_origin.hi; $RUN6/paper-extras-upworthy-20261106-2dc461f52ea4.json within_origin_auc_by_year.years["2015"].P2.auc_within_origin.estimate; $RUN6/paper-extras-upworthy-20261106-2dc461f52ea4.json within_origin_auc_by_year.years["2015"].P2.auc_within_origin.lo; $RUN6/paper-extras-upworthy-20261106-2dc461f52ea4.json within_origin_auc_by_year.years["2015"].P2.auc_within_origin.hi; $RUN6/paper-extras-upworthy-20261106-2dc461f52ea4.json within_origin_auc_by_year.years["2015"].P1.brier_skill; $RUN6/paper-extras-upworthy-20261106-2dc461f52ea4.json within_origin_auc_by_year.years["2015"].P2.brier_skill
2015 & 401 & 0.609 (0.545, 0.675) & 0.634 (0.568, 0.703) & 0.0243 & 0.0297 \\
\bottomrule
\end{tabular}
}
\end{table}

The weekly replay (Table~\ref{tab:replay}, Appendix~\ref{app:tables}) runs the released loop under its own rules: one scored card per admitted row, issued at the week start at or before the test's start, with a bundle refit at every quarterly origin. It joins outcomes at the first weekly join on or after the horizon, scores each with the card's own predictive (abstained cards on the withheld forecast) and evaluates the drift rule weekly.
\begin{itemize}
% src: $RUN6/replay-upworthy-20261106-2dc461f52ea4.json header.n_weeks 111, header.n_cards 10697 (tables/replay-upworthy.tex)
\item Span: 111 weeks and 10,697 cards.
% src: $RUN6/replay-upworthy-20261106-2dc461f52ea4.json totals.abstained 2022, totals.abstained_by_trigger reliability 271; $RUN6/paper-extras-upworthy-20261106-2dc461f52ea4.json issued_row_brier_skill.rungs.P2.abstained.n 3401; $RUN6/backtest-upworthy-20261106-2dc461f52ea4.json tables.abstention[cell=ALL].reliability 1915 (tables/replay-upworthy.tex)
\item Abstained: 2,022 cards, against 3,401 rows in the backtest. The replay refreshes reliability at every weekly join, so a stale report clears within weeks (271 cards on that trigger); the backtest reads it at the quarterly origin, so a stale report holds the whole quarter (1,915 rows).
% src: $RUN6/replay-upworthy-20261106-2dc461f52ea4.json totals.outcomes_joined 10697, header.n_outcomes 10697 (tables/replay-upworthy.tex)
\item Outcomes joined: every card (10,697).
% src: $RUN6/replay-upworthy-20261106-2dc461f52ea4.json totals.coverage_90_all 0.8965 (coverage_90_all_n 8675, the outcomes of issued cards: 10697 minus 2022 abstained)
\item Coverage: 89.6 percent of the 8,675 outcomes of issued cards fell in the card's 90 percent interval; an abstained card shows no interval.
% src: $RUN6/replay-upworthy-20261106-2dc461f52ea4.json totals.n_flagged_weeks 6; drift_onsets both 2014-03-31 and 2015-04-20; weeks[].flagged 2014-03-31, 2014-04-07, 2015-04-20 to 2015-05-11 (tables/replay-upworthy.tex)
\item Weeks flagged: 6, all in the smallest cell, two from 31 March 2014 and four from 20 April 2015 to the end of the archive (Appendix~\ref{app:replay}).
\end{itemize}

\subsection{The deployment-shaped archive}
\label{sec:asos-results}
Table~\ref{tab:fmt-asos} condenses Table~\ref{tab:main-asos} (Appendix~\ref{app:tables}), the same protocol on the deployment-shaped archive.
% src: $RUN6/paper-extras-asos-20261103-6bb5a15c5ff0.json gates_by_stratum.strata (ALL and metric_1 to metric_4, every verdict undecidable at this n); $RUN6/backtest-asos-20261103-6bb5a15c5ff0.json header.tests_scored 68
Every metric cell and the pooled table are undecidable at their own gates, a result 68 experiments could not have avoided, so this archive cannot say whether the protocol's findings hold on product experiments.
The abstention rate is high mostly because the support trigger fires at the early origins.
What the archive does show is that the prior is overconfident there: the skill interval lies below zero before adjustment and spans it after, and the expected calibration error is past the rejection line.

\begin{table}[t]\centering\scriptsize
% table: tab:main-asos; tables/fmt-asos.tex
\caption{The deployment-shaped archive, from Table~\ref{tab:main-asos} (Appendix~\ref{app:tables}). Rows: one figure each. The adjusted interval is for the table's intervals. The calibration percentile is against the clustered null.}
\label{tab:fmt-asos}
\resizebox{\columnwidth}{!}{% generated by paper/build_tables.py from the report keys named above each row (results_b)
\begin{tabular}{lr}
\toprule
figure & value \\
\midrule
% src: $RUN6/backtest-asos-20261103-6bb5a15c5ff0.json header.scored_rows
scored contrasts & 337 \\
% src: $RUN6/backtest-asos-20261103-6bb5a15c5ff0.json header.tests_scored
experiments scored & 68 \\
% src: $RUN6/backtest-asos-20261103-6bb5a15c5ff0.json header.origins_scored
nominal quarterly origins & seven \\
% src: $RUN6/paper-extras-asos-20261103-6bb5a15c5ff0.json gates_by_stratum.strata.ALL.dependence.base_rate; $RUN6/paper-extras-asos-20261103-6bb5a15c5ff0.json gates_by_stratum.strata.ALL.own_gate.n_200; $RUN6/paper-extras-asos-20261103-6bb5a15c5ff0.json gates_by_stratum.strata.ALL.own_gate.n_400
own gate at skill 0.20 and 0.10, tests (base rate 0.14) & 175 and 275 \\
% src: $RUN6/backtest-asos-20261103-6bb5a15c5ff0.json tables.aggregate[cell=ALL,rung=P1].abstention_rate
abstention rate & 0.457 \\
% src: $RUN6/backtest-asos-20261103-6bb5a15c5ff0.json tables.abstention[cell=ALL].support
support trigger, rows at the early origins & 126 \\
% src: $RUN6/backtest-asos-20261103-6bb5a15c5ff0.json tables.aggregate[cell=ALL,rung=P1].brier_skill; $RUN6/backtest-asos-20261103-6bb5a15c5ff0.json tables.aggregate[cell=ALL,rung=P1].brier_skill_lo; $RUN6/backtest-asos-20261103-6bb5a15c5ff0.json tables.aggregate[cell=ALL,rung=P1].brier_skill_hi
P1 Brier skill (90\%) & -0.1073 (-0.2366, -0.0020) \\
% src: $RUN6/backtest-asos-20261103-6bb5a15c5ff0.json tables.aggregate[cell=ALL,rung=P1].adjusted_lo; $RUN6/backtest-asos-20261103-6bb5a15c5ff0.json tables.aggregate[cell=ALL,rung=P1].adjusted_hi
P1 Brier skill, adjusted interval & (-0.3258, 0.1113) \\
% src: $RUN6/backtest-asos-20261103-6bb5a15c5ff0.json tables.aggregate[cell=ALL,rung=P1].ece
expected calibration error & 0.127 \\
% src: $RUN6/backtest-asos-20261103-6bb5a15c5ff0.json tables.aggregate[cell=ALL,rung=P1].ece_percentile
calibration percentile & 98.8 \\
% src: $RUN6/backtest-asos-20261103-6bb5a15c5ff0.json tables.aggregate[cell=ALL,rung=P1].design_effect
design effect & 3.8 \\
% src: $RUN6/backtest-asos-20261103-6bb5a15c5ff0.json tables.aggregate[cell=ALL,rung=P1].coverage_80; $RUN6/backtest-asos-20261103-6bb5a15c5ff0.json tables.aggregate[cell=ALL,rung=P1].coverage_90
coverage of the 80 and 90 percent intervals & 0.78 and 0.86 \\
% src: $RUN6/backtest-asos-20261103-6bb5a15c5ff0.json tables.aggregate[cell=ALL,rung=GBM+cal].brier_skill; $RUN6/backtest-asos-20261103-6bb5a15c5ff0.json tables.aggregate[cell=ALL,rung=GBM+cal].brier_skill_lo; $RUN6/backtest-asos-20261103-6bb5a15c5ff0.json tables.aggregate[cell=ALL,rung=GBM+cal].brier_skill_hi
boosted model, recalibrated Brier skill (90\%) & -0.0090 (-0.0756, 0.0446) \\
% src: $RUN6/backtest-asos-20261103-6bb5a15c5ff0.json tables.aggregate[cell=ALL,rung=local_eb+cal].brier_skill; $RUN6/backtest-asos-20261103-6bb5a15c5ff0.json tables.aggregate[cell=ALL,rung=local_eb+cal].brier_skill_lo; $RUN6/backtest-asos-20261103-6bb5a15c5ff0.json tables.aggregate[cell=ALL,rung=local_eb+cal].brier_skill_hi
local comparison set, recalibrated Brier skill (90\%) & 0.0150 (-0.0351, 0.0691) \\
% src: $RUN6/backtest-asos-20261103-6bb5a15c5ff0.json tables.decision_value[origin=ALL,budget=0.25,policy=p_win,rung=GBM].diff_random; $RUN6/backtest-asos-20261103-6bb5a15c5ff0.json tables.decision_value[origin=ALL,budget=0.25,policy=p_win,rung=GBM].diff_random_lo; $RUN6/backtest-asos-20261103-6bb5a15c5ff0.json tables.decision_value[origin=ALL,budget=0.25,policy=p_win,rung=GBM].diff_random_hi
boosted model, lift minus random order, quarter of tests (90\%) & 0.0420 (-0.0243, 0.1352) \\
% src: $RUN6/backtest-asos-20261103-6bb5a15c5ff0.json tables.decision_value[origin=ALL,budget=0.25,policy=p_win,rung=local_eb].diff_random; $RUN6/backtest-asos-20261103-6bb5a15c5ff0.json tables.decision_value[origin=ALL,budget=0.25,policy=p_win,rung=local_eb].diff_random_lo; $RUN6/backtest-asos-20261103-6bb5a15c5ff0.json tables.decision_value[origin=ALL,budget=0.25,policy=p_win,rung=local_eb].diff_random_hi
local comparison set, the same (90\%) & 0.0044 (-0.0660, 0.0934) \\
% src: $RUN6/backtest-asos-20261103-6bb5a15c5ff0.json tables.decision_value[origin=ALL,budget=0.25,policy=p_win,rung=GBM].diff_matched; $RUN6/backtest-asos-20261103-6bb5a15c5ff0.json tables.decision_value[origin=ALL,budget=0.25,policy=p_win,rung=local_eb].diff_matched
both, lift minus the matched order, quarter of tests & 0.000 and 0.000 \\
% src: $RUN6/paper-extras-asos-20261103-6bb5a15c5ff0.json impressions_effect_correlation.corr_log_impressions_abs_z
exposure stand-in, correlation with $|$standardized effect$|$ & 0.215 \\
% src: $RUN6/paper-extras-upworthy-20261106-2dc461f52ea4.json impressions_effect_correlation.corr_log_impressions_abs_z
the same correlation on the headline archive & 0.087 \\
\bottomrule
\end{tabular}
}
\end{table}

% src: $RUN6/backtest-asos-20261103-6bb5a15c5ff0.json tables.p2_selection_by_cell share_lambda_zero 1.00 in all four cells; tables.decision_value origin=ALL budget=0.25 policy=p_win diff_random, diff_matched (src/mootcourt/decision_value.py _matched_value)
P2 equals P1 because the guard holds its weight at zero everywhere. The two challengers that run without text on this archive sit at the base rate once recalibrated. Their captured lift spans zero against random order and is exactly zero against the matched order, which on this archive returns the value of the chosen set itself (Appendix~\ref{app:asos}).

\subsection{Synthetic truth and the positive control}
\label{sec:synthetic-results}
On the synthetic archive with no text signal (Appendix~\ref{app:synthetic}) the pooled prior recovers the generator's cell means and spreads within a few tenths, and the guard holds the precedent weight at zero everywhere.
\begin{itemize}
% src: $RUN6/backtest-synthetic-20261103-bf4fd286e2b2.json tables.aggregate[cell=ALL] rung=P1 coverage_80 0.7834
\item Coverage of the 80 percent interval: 78 percent.
% src: $RUN6/backtest-synthetic-20261103-bf4fd286e2b2.json tables.aggregate[cell=ALL] rung=P1 coverage_90 0.8908
\item Coverage of the 90 percent interval: 89 percent.
% src: $RUN6/backtest-synthetic-20261103-bf4fd286e2b2.json tables.aggregate[cell=ALL] rung=P1 brier_skill 0.0023 (-0.0128, 0.0168)
\item Brier skill: 0.002 (minus 0.013 to 0.017).
% src: $RUN6/backtest-synthetic-20261103-bf4fd286e2b2.json tables.aggregate[cell=ALL] rung=P1 ece_percentile 95.2
\item Expected calibration error: percentile 95.2 of the clustered null, just past the rejection line of 95.
% src: $RUN6/backtest-synthetic-20261103-bf4fd286e2b2.json tables.aggregate[cell=ALL] rung=P1 ece_percentile_independent 97.8
\item Expected calibration error under the independent null: percentile 97.8.
% src: $RUN6/backtest-synthetic-20261103-bf4fd286e2b2.json tables.decision_value origin=ALL budget=0.25 P1 captured_fraction 0.4453, random 0.2502, diff_random 0.1951 (0.1732, 0.2156)
\item Decision value: ranking by $\pwin$ gains plus 0.195 (0.173 to 0.216) for the prior itself, because the cells differ in their means.
\end{itemize}

% src: $RUN6/positive-control/backtest-synthetic-20261106-8cc34d3af720.json header.admitted_rows 4000 (text_effect_strength 1.0, a third of the headline_only prior sd of 3.0); $RUN6/paper-extras-synthetic-20261103-bf4fd286e2b2.json positive_control_sensitivity.rows strength_over_prior_sd headline_only 0.333
The positive control plants a text lift of 1.0 percent between headlines with and without a signal word, a third of the prior standard deviation in the headline-only cell. It runs under two designs, read apart in Table~\ref{tab:fmt-control} (from Tables~\ref{tab:positive-control}, \ref{tab:positive-control-gain} and~\ref{tab:positive-control-rungs}, Appendix~\ref{app:tables}). The end-of-ledger split holds out 15 percent of a ledger's tests. The rolling backtest is the design the real archive gets. Each recovered figure is a fraction of an oracle's gain in log score, not a fraction of the lift, and the oracle is P1 told the planted term.

\begin{table}[t]\centering\scriptsize
% src: $RUN6/paper-extras-synthetic-20261103-bf4fd286e2b2.json positive_control_sensitivity.rows (12 seeds, 36 cells per point, held_out_share_of_tests 0.15); positive_control_rungs.runs; tables/positive-control-gain.tex; tables/fmt-control.tex
\caption{The positive control by ledger size at a planted lift of one percent, under the two designs. Rows per cell count ledger rows in the first panel and scored rows in the second. The guard's admission counts cells in the first panel and cell-origin pairs in the second, and the rolling backtest's fractions carry no interval in the report.}
\label{tab:fmt-control}
\resizebox{\columnwidth}{!}{% generated by paper/build_tables.py from the report keys named above each row (results_b)
\begin{tabular}{rrlrr}
\toprule
tests & rows per cell & guard admits the lift & recovered (90\%) & headline-only recovered \\
\midrule
% src: $RUN6/paper-extras-synthetic-20261103-bf4fd286e2b2.json positive_control_sensitivity.rows[n_tests=300,text_effect_strength=1.0].n_tests; $RUN6/paper-extras-synthetic-20261103-bf4fd286e2b2.json positive_control_sensitivity.rows[n_tests=300,text_effect_strength=1.0].rows_per_cell; $RUN6/paper-extras-synthetic-20261103-bf4fd286e2b2.json positive_control_sensitivity.rows[n_tests=300,text_effect_strength=1.0].share_admitted; $RUN6/paper-extras-synthetic-20261103-bf4fd286e2b2.json positive_control_sensitivity.rows[n_tests=300,text_effect_strength=1.0].recovered_fraction; $RUN6/paper-extras-synthetic-20261103-bf4fd286e2b2.json positive_control_sensitivity.rows[n_tests=300,text_effect_strength=1.0].recovered_fraction_lo; $RUN6/paper-extras-synthetic-20261103-bf4fd286e2b2.json positive_control_sensitivity.rows[n_tests=300,text_effect_strength=1.0].recovered_fraction_hi
300 & 200 & 6 percent of cells & 0.00 (0.00, 0.00) &  \\
% src: $RUN6/paper-extras-synthetic-20261103-bf4fd286e2b2.json positive_control_sensitivity.rows[n_tests=1000,text_effect_strength=1.0].n_tests; $RUN6/paper-extras-synthetic-20261103-bf4fd286e2b2.json positive_control_sensitivity.rows[n_tests=1000,text_effect_strength=1.0].rows_per_cell; $RUN6/paper-extras-synthetic-20261103-bf4fd286e2b2.json positive_control_sensitivity.rows[n_tests=1000,text_effect_strength=1.0].share_admitted; $RUN6/paper-extras-synthetic-20261103-bf4fd286e2b2.json positive_control_sensitivity.rows[n_tests=1000,text_effect_strength=1.0].recovered_fraction; $RUN6/paper-extras-synthetic-20261103-bf4fd286e2b2.json positive_control_sensitivity.rows[n_tests=1000,text_effect_strength=1.0].recovered_fraction_lo; $RUN6/paper-extras-synthetic-20261103-bf4fd286e2b2.json positive_control_sensitivity.rows[n_tests=1000,text_effect_strength=1.0].recovered_fraction_hi
1,000 & 667 & 19 percent of cells & 0.02 (0.00, 0.03) &  \\
% src: $RUN6/paper-extras-synthetic-20261103-bf4fd286e2b2.json positive_control_rungs.runs[n_tests=2000].n_tests; $RUN6/positive-control/backtest-synthetic-20261106-8cc34d3af720.json header.admitted_rows; $RUN6/paper-extras-synthetic-20261103-bf4fd286e2b2.json positive_control_rungs.runs[n_tests=2000].rows_per_cell_scored; $RUN6/positive-control/backtest-synthetic-20261106-8cc34d3af720.json tables.p2_selection_by_cell[cell=headline_only].share_lambda_zero
2,000 (4,000 contrasts) & about 1,200 & none in the planted cell &  &  \\
% src: $RUN6/paper-extras-synthetic-20261103-bf4fd286e2b2.json positive_control_sensitivity.rows[n_tests=11100,text_effect_strength=1.0].n_tests; $RUN6/paper-extras-synthetic-20261103-bf4fd286e2b2.json positive_control_sensitivity.rows[n_tests=11100,text_effect_strength=1.0].rows_per_cell; $RUN6/paper-extras-synthetic-20261103-bf4fd286e2b2.json positive_control_sensitivity.rows[n_tests=11100,text_effect_strength=1.0].share_admitted; $RUN6/paper-extras-synthetic-20261103-bf4fd286e2b2.json positive_control_sensitivity.rows[n_tests=11100,text_effect_strength=1.0].recovered_fraction; $RUN6/paper-extras-synthetic-20261103-bf4fd286e2b2.json positive_control_sensitivity.rows[n_tests=11100,text_effect_strength=1.0].recovered_fraction_lo; $RUN6/paper-extras-synthetic-20261103-bf4fd286e2b2.json positive_control_sensitivity.rows[n_tests=11100,text_effect_strength=1.0].recovered_fraction_hi; $RUN6/paper-extras-synthetic-20261103-bf4fd286e2b2.json positive_control_sensitivity.rows[n_tests=11100,text_effect_strength=1.0].by_cell[cell=headline_only].recovered_fraction
11,100 & 7,400 & 97 percent of cells & 0.18 (0.16, 0.20) & 0.10 \\
\bottomrule
\end{tabular}
}
\end{table}

What the guard admits:
\begin{itemize}
% src: $RUN6/positive-control/backtest-synthetic-20261106-8cc34d3af720.json tables.p2_selection cell=headline_only: lambda 0.0 at all 8 origins; tables.aggregate[cell=headline_only] rung=P2 logscore -2.6348 equal to rung=P1; $RUN6/paper-extras-synthetic-20261103-bf4fd286e2b2.json positive_control_rungs.runs[n_tests=2000].rows_per_cell_scored 1192
\item In the rolling backtest at 2,000 tests, with 1,192 scored rows per cell, the guard admits no weight in the planted headline-only cell at any of the eight origins, so P2 equals P1 in that cell.
% src: $RUN6/positive-control/backtest-synthetic-20261106-8cc34d3af720.json tables.p2_selection cell=headline_only: gain -0.0002 to 0.0093
\item In the same backtest the in-selection gains in that headline-only cell run from minus 0.0002 to 0.0093 nats, none past two standard errors.
% src: $RUN6/paper-extras-synthetic-20261103-bf4fd286e2b2.json positive_control_sensitivity.rows n_tests 11100 rows_per_cell 7400 text_effect_strength 0.0 share_admitted 0.000
\item With nothing planted, at 7,400 ledger rows per cell, the guard admits nothing.
\end{itemize}

What P2 recovers under the end-of-ledger split:
\begin{itemize}
% src: $RUN6/paper-extras-synthetic-20261103-bf4fd286e2b2.json positive_control_sensitivity.rows n_tests 11100 text_effect_strength 1.0 realized_gain 0.00369 (0.00334, 0.00401)
\item At 11,100 tests and a lift of 1 percent, P2's held-out gain over P1 is 0.0037 nats per row (0.0033 to 0.0040).
% src: $RUN6/paper-extras-synthetic-20261103-bf4fd286e2b2.json positive_control_sensitivity.rows n_tests 11100 text_effect_strength 1.0 oracle_gain 0.02042 (0.01890, 0.02209)
\item At 11,100 tests and a lift of 1 percent, the oracle gains 0.0204 (0.0189 to 0.0221).
% src: $RUN6/paper-extras-synthetic-20261103-bf4fd286e2b2.json positive_control_sensitivity.rows n_tests 11100 text_effect_strength 1.0 recovered_fraction 0.180, by_cell headline_only recovered_fraction 0.0995
\item At 11,100 tests and a lift of 1 percent, P2 recovers about a fifth of that gain pooled and a tenth in the headline-only cell.
% src: $RUN6/paper-extras-synthetic-20261103-bf4fd286e2b2.json positive_control_sensitivity.rows n_tests 11100 text_effect_strength 2.0 recovered_fraction 0.274 and 4.0 0.254 (share_admitted 1.000 at both); tables/positive-control-gain.tex
\item At 11,100 tests and lifts of 2 and 4 percent, where the guard admits in every cell, P2 recovers about a quarter.
% src: $RUN6/paper-extras-synthetic-20261103-bf4fd286e2b2.json positive_control_sensitivity.rows n_tests 300 and 1000 at text_effect_strength 0.5 and 1.0: recovered_fraction 0.000, 0.000, 0.014, 0.016 (tables/positive-control-gain.tex)
\item At 300 and 1,000 tests and lifts of 1 percent and below, P2 recovers at most 0.02.
\end{itemize}
An earlier version of this curve exposed the period-rule defect of Section~\ref{sec:guard-period}.

What the rolling backtest returns, the figures to set against the real archive:
\begin{itemize}
% src: $RUN6/paper-extras-synthetic-20261103-bf4fd286e2b2.json positive_control_rungs.runs[n_tests=11100].by_cell ALL rungs[rung=P2].recovered_fraction_logscore 0.0919; headline_only 0.0319
\item At 11,100 tests P2 recovers 0.09 of the oracle's gain pooled and 0.03 in the headline-only cell.
% src: $RUN6/paper-extras-synthetic-20261103-bf4fd286e2b2.json positive_control_rungs.runs[n_tests=11100].by_cell headline_only rungs[rung=P2].rows 6354; $RUN6/backtest-upworthy-20261106-2dc461f52ea4.json tables.abstention[cell=headline_only].n 5966
\item At 11,100 tests that planted cell has 6,354 scored rows and the real headline-only cell 5,966, so it is the cell of the real one's size. The match is in size, not in the effects.
% src: $RUN6/paper-extras-synthetic-20261103-bf4fd286e2b2.json positive_control_rungs.runs[n_tests=11100].by_cell headline_only rungs[rung=P2].diff_logscore_vs_p1 0.0004 (0.0001, 0.0007); $RUN6/backtest-upworthy-20261106-2dc461f52ea4.json tables.aggregate[cell=headline_only,rung=P2] logscore_vs_p1 0.0002 (-0.0007, 0.0010)
\item At 11,100 tests P2 minus P1 in that planted cell is 0.0004 nats, which lies inside the real headline-only interval of minus 0.0007 to 0.0010.
\end{itemize}
So the real interval contains what a planted signal of that size returns, and the level result on the real archive does not refute one. The control shows that this backtest is weak at this size: it bounds what a precedent signal of this form could have shown, and it does not separate an archive with no such signal from one whose signal P2 recovers this little of.

\begin{table}[t]\centering\scriptsize
% src: $RUN6/paper-extras-synthetic-20261103-bf4fd286e2b2.json positive_control_rungs.runs (planted lift 1.0) n_tests 2000 and 11100 (tables/positive-control-rungs.tex); tables/fmt-control-rungs.tex
\caption{Rungs on the planted ledger at the two sizes under the rolling backtest, pooled unless the row names the headline-only cell. A negative Brier difference favours the rung; a positive log score difference does.}
\label{tab:fmt-control-rungs}
\resizebox{\columnwidth}{!}{% generated by paper/build_tables.py from the report keys named above each row (results_b)
\begin{tabular}{lrr}
\toprule
rung and measure & 2,000 tests & 11,100 tests \\
\midrule
% src: $RUN6/paper-extras-synthetic-20261103-bf4fd286e2b2.json positive_control_rungs.runs[n_tests=2000].by_cell["ALL"].rungs[rung=text_lr].brier_skill; $RUN6/paper-extras-synthetic-20261103-bf4fd286e2b2.json positive_control_rungs.runs[n_tests=2000].by_cell["ALL"].rungs[rung=text_lr].brier_skill_lo; $RUN6/paper-extras-synthetic-20261103-bf4fd286e2b2.json positive_control_rungs.runs[n_tests=2000].by_cell["ALL"].rungs[rung=text_lr].brier_skill_hi; $RUN6/paper-extras-synthetic-20261103-bf4fd286e2b2.json positive_control_rungs.runs[n_tests=11100].by_cell["ALL"].rungs[rung=text_lr].brier_skill; $RUN6/paper-extras-synthetic-20261103-bf4fd286e2b2.json positive_control_rungs.runs[n_tests=11100].by_cell["ALL"].rungs[rung=text_lr].brier_skill_lo; $RUN6/paper-extras-synthetic-20261103-bf4fd286e2b2.json positive_control_rungs.runs[n_tests=11100].by_cell["ALL"].rungs[rung=text_lr].brier_skill_hi
text classifier, Brier skill & -0.014 (-0.027, 0.001) & -0.017 (-0.022, -0.011) \\
% src: $RUN6/paper-extras-synthetic-20261103-bf4fd286e2b2.json positive_control_rungs.runs[n_tests=2000].by_cell["headline_only"].rungs[rung=GBM].diff_brier_vs_p1; $RUN6/paper-extras-synthetic-20261103-bf4fd286e2b2.json positive_control_rungs.runs[n_tests=2000].by_cell["headline_only"].rungs[rung=GBM].diff_brier_vs_p1_lo; $RUN6/paper-extras-synthetic-20261103-bf4fd286e2b2.json positive_control_rungs.runs[n_tests=2000].by_cell["headline_only"].rungs[rung=GBM].diff_brier_vs_p1_hi
boosted model, headline-only, Brier minus P1 & -0.0044 (-0.0078, -0.0005) &  \\
% src: $RUN6/paper-extras-synthetic-20261103-bf4fd286e2b2.json positive_control_rungs.runs[n_tests=2000].by_cell["headline_only"].rungs[rung=GBM].diff_logscore_vs_p1
boosted model, headline-only, log score minus P1 & -0.020 &  \\
% src: $RUN6/paper-extras-synthetic-20261103-bf4fd286e2b2.json positive_control_rungs.runs[n_tests=2000].by_cell["ALL"].rungs[rung=P2].diff_logscore_vs_p1; $RUN6/paper-extras-synthetic-20261103-bf4fd286e2b2.json positive_control_rungs.runs[n_tests=2000].by_cell["ALL"].rungs[rung=P2].diff_logscore_vs_p1_lo; $RUN6/paper-extras-synthetic-20261103-bf4fd286e2b2.json positive_control_rungs.runs[n_tests=2000].by_cell["ALL"].rungs[rung=P2].diff_logscore_vs_p1_hi; $RUN6/paper-extras-synthetic-20261103-bf4fd286e2b2.json positive_control_rungs.runs[n_tests=11100].by_cell["ALL"].rungs[rung=P2].diff_logscore_vs_p1; $RUN6/paper-extras-synthetic-20261103-bf4fd286e2b2.json positive_control_rungs.runs[n_tests=11100].by_cell["ALL"].rungs[rung=P2].diff_logscore_vs_p1_lo; $RUN6/paper-extras-synthetic-20261103-bf4fd286e2b2.json positive_control_rungs.runs[n_tests=11100].by_cell["ALL"].rungs[rung=P2].diff_logscore_vs_p1_hi
P2, log score minus P1 (nats) & 0.0000 (-0.0002, 0.0003) & 0.0018 (0.0014, 0.0023) \\
% src: $RUN6/paper-extras-synthetic-20261103-bf4fd286e2b2.json positive_control_rungs.runs[n_tests=11100].by_cell["ALL"].rungs[rung=P2].recovered_fraction_logscore
P2, share of the oracle's gain &  & 0.09 \\
\bottomrule
\end{tabular}
}
\end{table}

The other rungs on the planted ledgers, under the rolling backtest (Table~\ref{tab:fmt-control-rungs}):
\begin{itemize}
% src: $RUN6/paper-extras-synthetic-20261103-bf4fd286e2b2.json positive_control_rungs.runs[n_tests=2000 and 11100].by_cell["ALL"].rungs rung=text_lr, P1 brier_skill; tables/fmt-control-rungs.tex
\item At both sizes the text classifier's Brier skill is at or below zero and below P1's.
% src: $RUN6/paper-extras-synthetic-20261103-bf4fd286e2b2.json positive_control_decision_value n2000, n11100 paired.budgets["0.25"].differences.text_lr_minus_P1 adjusted_lo, adjusted_hi; tables/fmt-ranking-paired.tex
\item At both sizes and a quarter of the tests the text classifier ranks below P1, its adjusted interval clear of zero (Table~\ref{tab:fmt-ranking-paired}).
% src: $RUN6/paper-extras-synthetic-20261103-bf4fd286e2b2.json positive_control_rungs.runs[n_tests=2000].by_cell["headline_only"].rungs[rung=GBM] diff_brier_vs_p1, diff_logscore_vs_p1; tables/fmt-control-rungs.tex
\item At 2,000 tests the boosted model beats P1 on the Brier score in the headline-only cell and loses to it on the log score there.
% src: $RUN6/paper-extras-synthetic-20261103-bf4fd286e2b2.json positive_control_rungs.runs[n_tests=11100].by_cell["headline_only"].rungs[rung=GBM] diff_brier_vs_p1, diff_logscore_vs_p1; tables/fmt-control-rungs.tex
\item At 11,100 tests the boosted model loses to P1 on the log score in the headline-only cell, and its Brier difference there spans zero.
% src: $RUN6/paper-extras-synthetic-20261103-bf4fd286e2b2.json positive_control_rungs.runs[n_tests=2000].by_cell["ALL"].rungs[rung=P2] diff_brier_vs_p1, diff_logscore_vs_p1; tables/fmt-control-rungs.tex
\item At 2,000 tests P2 is within the interval of P1 on both scores.
% src: $RUN6/paper-extras-synthetic-20261103-bf4fd286e2b2.json positive_control_rungs.runs[n_tests=11100].by_cell["ALL"].rungs diff_brier_vs_p1, diff_logscore_vs_p1; positive_control_decision_value.n11100.paired.budgets["0.25"].differences.P2_minus_P1 adjusted_lo, adjusted_hi; tables/fmt-ranking-paired.tex
\item At 11,100 tests P2 is above P1 on both scores, the only rung that is, and at a quarter of the tests its paired ranking gain over P1 spans zero after adjustment (Table~\ref{tab:fmt-ranking-paired}).
\end{itemize}
The text classifier is the rung that ranks best on the real archive, and it fails this control on skill and ranks below the prior here. The control therefore neither validates nor explains its value on the real archive: whatever it reads there is not a location shift of this form, and one exploratory split is all the evidence for it.

% $REPORTS = /Volumes/Crucial X10/mootcourt-data/reports. Ablation reports: $REPORTS/ablations/backtest-upworthy-20261103-4bfb354b87e3--<setting>.json, each with tables.aggregate[cell=ALL] rows for P1 and P2.
\section{Ablations and Learning Curve}
\label{sec:ablations}

\paragraph{Does any setting of the precedent rung beat the prior?}
Table~\ref{tab:ablations-upworthy} (Appendix~\ref{app:tables}) changes one setting of P2 at a time on the same nine origins and 13,247 rows.
% src: $REPORTS/ablations/...--lambda-0.0.json P2 logscore -0.8812 brier_skill 0.0015; --lambda-0.2.json -0.9662, -0.0558; --lambda-0.5.json -0.9725, -0.0614; --lambda-1.0.json -0.9922, -0.0787
Fixing the precedent weight makes the loss monotonic: $\lambda$ of 0, 0.2, 0.5 and 1 give logarithmic scores of minus 0.881, minus 0.966, minus 0.973 and minus 0.992 and Brier skills of 0.001, minus 0.056, minus 0.061 and minus 0.079; $\lambda = 0$ reproduces P1 to the fourth decimal, as the construction requires.
% src: --bandwidth-mult-0.125.json P2 logscore -0.8892 brier_skill -0.0012 (-0.0048, 0.0026); --bandwidth-mult-0.25.json -0.8956; --bandwidth-mult-0.5.json -0.8961; --bandwidth-mult-1.0.json -0.8963; main run -0.8957
Fixing the bandwidth multiplier at 1/8, 1/4, 1/2 or 1 gives minus 0.889, minus 0.896, minus 0.896 and minus 0.896; the narrowest kernel is the least bad (Brier skill minus 0.001, minus 0.005 to 0.003), still does not beat the prior, and is better than the selected bandwidth by about 0.007 nats, so the selection also overfits the bandwidth.
% src: --features-text.json P2 logscore -0.8961 brier_skill -0.0079; --features-structured.json -0.8965, -0.0100; main run (both) -0.8957, -0.0078
Text features alone, structured features alone and both are within 0.001 nats of each other (minus 0.8961, minus 0.8965 and minus 0.8957): the headline text adds nothing to the structured context, and neither adds anything to the prior.
% src: --p2-selection-forward.json P2 logscore -0.8909 brier_skill -0.0001 (-0.0044, 0.0037) direction_accuracy 0.5759; --p2-selection-forward+recency-select.json -0.8910, -0.0002, p2_selection_by_cell half_life_days_median 90, 180, 90; --recency-select.json -0.8957, -0.0058, half-lives 90, 135, 90
Forward-chaining selection, which validates a candidate only on rows that start after an inner cut and so never lets a validation row borrow a precedent that did not yet exist, is the most principled rule and the least bad variant: minus 0.891 and a Brier skill of minus 0.000 (minus 0.004 to 0.004), with or without a recency kernel (half-lives of 90 to 180 days where it uses one; under leave-one-test-out recency moves the score by less than 0.001 nats). It ties the prior and does not beat it, so the default stays the rule the plan named; the rule for changing it (an advantage whose 90 percent interval clears zero) is in the method notes.
% src: --p2-guard-1.0.json P2 logscore -0.9293 brier_skill -0.0251; --p2-guard-off.json -0.9360, -0.0311; main run (guard 2.0) -0.8957, -0.0078
% src: PENDING reports/paper-extras-<snapshot>.json guard_admission_rates (the 39 and 8 percent on no-signal synthetic cells; today docs/method.md call 71)
Loosening the guard to one standard error gives minus 0.929 and a Brier skill of minus 0.025; removing it, minus 0.936 and minus 0.031; the two-standard-error guard of the main run, minus 0.896 and minus 0.008. On synthetic ledgers with no precedent signal the one-standard-error guard admitted a positive weight in 39 percent of cell fits and the two-standard-error guard in 8 percent, which is why the default is two.

\paragraph{Settings of the prior.}
% src: --kappa-0.0.json P1 logscore -0.8815; --kappa-5.0.json -0.8802; --kappa-80.0.json -0.8788; main run (kappa 20) -0.8812; --prior-normal.json P1 logscore -0.9572 brier_skill -0.0158 (-0.0215, -0.0105)
The pooling strength barely matters: $\kappa$ of 0, 5, 20 and 80 give P1 logarithmic scores of minus 0.882, minus 0.880, minus 0.881 and minus 0.879. The prior's shape matters: a normal prior in place of the Student-t costs 0.076 nats (minus 0.957) and turns the prior's Brier skill negative (minus 0.0158, minus 0.0215 to minus 0.0105), so the heavy tails of headline effects are real.
% src: $REPORTS/ablations/backtest-asos-20261103-6bb5a15c5ff0--prior-normal.json P1 logscore 4.3318 against main 4.6095, brier_skill -0.3678 (-0.6886, -0.1409); --kappa-0.0.json 4.5891, --kappa-5.0.json 4.6016, --kappa-80.0.json 4.6168
The same two findings hold on the deployment-shaped archive (Table~\ref{tab:ablations-asos} in Appendix~\ref{app:tables}): the normal prior loses 0.278 nats and falls to a Brier skill of minus 0.368, and $\kappa$ moves the score by under 0.03 nats.


\paragraph{Learning curve.}
% src: $REPORTS/learning-curve-upworthy-20261103-4bfb354b87e3.json final_origin 2015-01-01, heldout_rows 1143, heldout_tests 318, fractions [0.05, 0.1, 0.25, 0.5, 1.0], seeds 3; rows rung=P1 n_fit_rows 601 (fraction 0.05) and 11997 (1.00), skill 0.0266 (0.0131, 0.0429) and 0.0295 (0.0141, 0.0492)
Table~\ref{tab:learning-upworthy} (Appendix~\ref{app:tables}) subsamples the history visible at the final origin (January 2015) by whole tests at five fractions, three seeds each, and scores the same 1,143 held-out rows in 318 tests every time. P1 sits at a Brier skill of about 0.03 at every history size, 0.027 (0.013 to 0.043) with 601 history rows and 0.030 (0.014 to 0.049) with 11,997: twenty times more history adds nothing measurable, as a prior that has learned a base rate and a spread should behave.
% src: same file, rows rung=GBM skill -0.2836 (fraction 0.05) to -0.1161 (1.00); rung=local_eb 0.0090 (-0.0029, 0.0230) to 0.0025 (-0.0220, 0.0320); rung=text_lr 0.0004 (-0.0114, 0.0103) to -0.0100 (-0.0366, 0.0126); rung=P2 equals rung=P1 at every fraction
The gradient-boosted challenger improves with history (minus 0.284 to minus 0.116) and never reaches zero; the local empirical-Bayes set and the text classifier stay at zero within their intervals at every fraction; and P2 repeats P1 because the guard holds $\lambda$ at zero at every history size.
% src: $REPORTS/gates.json by_target 0.90 n_400 375; learning-curve-upworthy rows n_fit_rows minimum 601
This quarter's skill of 0.03 clears zero while the pooled skill over nine origins (Table~\ref{tab:main-upworthy}) does not; the two do not disagree, because the pooled figure averages nine quarters whose fits saw less history and whose base rates differ, while the learning curve scores one quarter with the fullest history and its bootstrap over 318 tests does not include the noise of choosing that quarter. The smallest history (601 rows) is above the 375-row gate, so no fraction here is undecidable.
% src: $REPORTS/learning-curve-asos-20261103-6bb5a15c5ff0.json heldout_rows 56, heldout_tests 9, skipped_runs {0.05: 3}; rows rung=P1 skill -0.0736 (-0.1805, 0.0088) at fraction 0.10 and -0.1824 (-0.5001, -0.0368) at 1.00
On the deployment-shaped archive the final origin holds 56 rows in 9 tests, the five-percent subsample cannot be fit, and every interval is wide (P1 from minus 0.074, minus 0.180 to 0.009, to minus 0.182, minus 0.500 to minus 0.037); the curve is produced for completeness (Table~\ref{tab:learning-asos}, Appendix~\ref{app:tables}) and decides nothing.


% $RUN6 = paper/data/run8-randomization-window
\section{The Hypotheses Read Back}
\label{sec:verdicts}
Section~\ref{sec:hypotheses} states the six hypotheses and the three verdict words. This section says when the hypotheses were written, reads them back in Table~\ref{tab:verdicts}, and names what held and what did not.

\paragraph{When they were written.}
\begin{itemize}
% src: git log of paper/plan.md: one commit, 145eaf7 (3 November 2026), which also holds paper/data/run2-prelaunch-clean/gates.json and run-log.json; the plan's section 1a reads the first real run
\item Commit 145eaf7 (3 November 2026) holds the plan beside the reports of the second run on the same split, and the plan's section 1a reads the first. So the hypotheses postdate those two runs: the cautionary run of Section~\ref{sec:cautionary} and the pre-launch-clean run.
% src: references.bib mootcourt2026plan: tag plan-registered-2026-11-04, commit f0b088c794dcea1a1eac9872cd15aeb1c2412976, pushed 2026-11-04T17:19:12Z, plan.md sha256 e6dab9062fd1dc59734e5e4374ba429d3edc06efe681f68bb81424d88f45b613, byte-identical to paper/plan.md of 145eaf7
\item A byte-identical copy of the plan was posted publicly on 4 November 2026, after those runs and after the first review of a draft \cite{mootcourt2026plan}. The commit and the tag named in that reference are the record.
% src: paper/plan.md line 31 (H5: "more positive lift than random"), the plan's E4 and E12; docs/method.md call 201 (the arm-count-matched order); this paper's Section 5 (the reading rule)
\item Four readings postdate the data: the pooled interval as the primary claim, the arm-count-matched baseline of the fifth hypothesis, reading that hypothesis over the challengers as well as the ladder, and the verdict words themselves. On the ladder alone the fifth hypothesis does not hold under the matched baseline.
% src: docs/superpowers/plans/2026-11-06-paper-v2-plan.md D1, D2 (the randomization window, 6 November 2026)
\item The date rule of Section~\ref{sec:setting} changed the rows these hypotheses are read on after they were written, and it did not change their wording.
\end{itemize}
Written after two runs on the same split and read by rules partly fixed after the data, these hypotheses do not make this a pre-registered study, and the verdicts below are exploratory like every other number in the paper.

\begin{table*}[t]\centering\scriptsize
% src: $RUN6/backtest-upworthy-20261106-2dc461f52ea4.json and $RUN6/backtest-asos-20261103-6bb5a15c5ff0.json (the main backtests); $RUN6/paper-extras-upworthy-20261106-2dc461f52ea4.json keys p1_minus_p0_by_origin, issued_ece_by_origin, issued_row_brier_skill, abstention_composition and drift_matched_null, bootstrap_resamples 500; $RUN6/learning-curve-upworthy-20261106-2dc461f52ea4.json; $RUN6/paper-extras-synthetic-20261103-bf4fd286e2b2.json positive_control_rungs; tables.decision_value budget=0.25, family_size 6; tables/verdicts.tex
\caption{The six hypotheses read back part by part under the rule of Section~\ref{sec:hypotheses}. Rows: the parts of each hypothesis, under a line that gives its overall reading and the section that holds the result; H5 is one part, read rung by rung, and H2 is read two ways. Columns: what the plan fixed, in short (Section~\ref{sec:hypotheses} gives the wording), what was measured, and the verdict on the part. A pair of numbers after a value is its 90 percent bootstrap interval by test; an adjusted interval or bound is the Bonferroni-adjusted one of the decision table. H1 reads the log score in nats, H2 and H4 the Brier skill against the base rate, and H5 the captured lift over the baseline at a quarter of the tests.}
\label{tab:verdicts}
% generated by paper/build_tables.py from the report keys named above each row (limits)
\begin{tabular}{@{}p{0.10\textwidth}p{0.21\textwidth}p{0.45\textwidth}p{0.17\textwidth}@{}}
\toprule
\raggedright hypothesis & \raggedright what was planned & \raggedright what was measured & \raggedright verdict \tabularnewline
\midrule
\multicolumn{4}{@{}l}{\textbf{H1, in part} (see~\ref{sec:ladder-results})} \\
% src: $RUN6/backtest-upworthy-20261106-2dc461f52ea4.json tables.score_by_origin[origin=2015-01-01T00:00:00+00:00,rung=P1].logscore; $RUN6/backtest-upworthy-20261106-2dc461f52ea4.json tables.score_by_origin[origin=2015-01-01T00:00:00+00:00,rung=P0].logscore; $RUN6/paper-extras-upworthy-20261106-2dc461f52ea4.json p1_minus_p0_by_origin.by_origin[origin=2015-01-01T00:00:00+00:00].difference; $RUN6/paper-extras-upworthy-20261106-2dc461f52ea4.json p1_minus_p0_by_origin.by_origin[origin=2015-01-01T00:00:00+00:00].lo; $RUN6/paper-extras-upworthy-20261106-2dc461f52ea4.json p1_minus_p0_by_origin.by_origin[origin=2015-01-01T00:00:00+00:00].hi; $RUN6/paper-extras-upworthy-20261106-2dc461f52ea4.json p1_minus_p0_by_origin.by_origin_clear_of_zero
\raggedright H1 part 1 & \raggedright P1 beats P0 on log score at every origin after burn-in, P0 scored with a normal density, as in the code committed with the plan & \raggedright above at every origin after burn-in; the narrowest margin January 2015, minus 0.730 against minus 0.761; paired by test the difference there is plus 0.031, 0.007 to 0.055, and clear of zero at all seven origins & \raggedright held \tabularnewline
% src: $RUN6/backtest-upworthy-20261106-2dc461f52ea4.json tables.aggregate[cell=headline_only,rung=P2].logscore_vs_p1; $RUN6/backtest-upworthy-20261106-2dc461f52ea4.json tables.aggregate[cell=headline_only,rung=P2].logscore_vs_p1_lo; $RUN6/backtest-upworthy-20261106-2dc461f52ea4.json tables.aggregate[cell=headline_only,rung=P2].logscore_vs_p1_hi; $RUN6/paper-extras-synthetic-20261103-bf4fd286e2b2.json positive_control_rungs.runs[n_tests=11100].by_cell["headline_only"].rungs[rung=P2].diff_logscore_vs_p1
\raggedright H1 part 2 & \raggedright P2 beats P1 on headline\_only & \raggedright plus 0.0002, minus 0.0007 to 0.0010; the interval contains the 0.0004 that a planted lift returns in a cell of this size under the same backtest & \raggedright not shown \tabularnewline
% src: $RUN6/backtest-upworthy-20261106-2dc461f52ea4.json tables.aggregate[cell=image_only,rung=P2].logscore_vs_p1; $RUN6/backtest-upworthy-20261106-2dc461f52ea4.json tables.aggregate[cell=image_only,rung=P2].logscore_vs_p1_lo; $RUN6/backtest-upworthy-20261106-2dc461f52ea4.json tables.aggregate[cell=image_only,rung=P2].logscore_vs_p1_hi
\raggedright H1 part 3 & \raggedright P2 does not beat P1 on image\_only & \raggedright plus 0.0000 nats, 0.0000 to 0.0000: the guard admitted no weight in this cell at any origin & \raggedright held \tabularnewline
\midrule
\multicolumn{4}{@{}l}{\textbf{H2, in part: the gate claim failed; skill not shown read pooled and undecidable as written; above zero at one earlier origin, where the base rate lags}} \\
% src: $RUN6/backtest-upworthy-20261106-2dc461f52ea4.json tables.aggregate[cell=ALL,rung=P2].brier_skill; $RUN6/backtest-upworthy-20261106-2dc461f52ea4.json tables.aggregate[cell=ALL,rung=P2].brier_skill_lo; $RUN6/backtest-upworthy-20261106-2dc461f52ea4.json tables.aggregate[cell=ALL,rung=P2].brier_skill_hi
\raggedright H2, read pooled & \raggedright P2's Brier skill positive, its interval clear of zero, pooled over origins & \raggedright 0.0032, minus 0.0004 to 0.0072, not clear of zero & \raggedright not shown \tabularnewline
% src: $RUN6/learning-curve-upworthy-20261106-2dc461f52ea4.json rows[fraction=1.0,rung=P2].skill; $RUN6/learning-curve-upworthy-20261106-2dc461f52ea4.json rows[fraction=1.0,rung=P2].skill_lo; $RUN6/learning-curve-upworthy-20261106-2dc461f52ea4.json rows[fraction=1.0,rung=P2].skill_hi; $RUN6/learning-curve-upworthy-20261106-2dc461f52ea4.json heldout_tests; $RUN6/paper-extras-upworthy-20261106-2dc461f52ea4.json gates_by_stratum.strata.ALL.own_gate.n_400; $RUN6/paper-extras-upworthy-20261106-2dc461f52ea4.json drift_matched_null.rungs.P0.by_origin["2015-01-01T00:00:00+00:00"].mean_forecast; $RUN6/paper-extras-upworthy-20261106-2dc461f52ea4.json drift_matched_null.rungs.P0.by_origin["2015-01-01T00:00:00+00:00"].event_rate; $RUN6/backtest-upworthy-20261106-2dc461f52ea4.json tables.origins[origin=2015-04-01T00:00:00+00:00].n_tests
\raggedright H2, as written & \raggedright the same at the final origin & \raggedright at January 2015, the latest origin with enough held-out rows for the learning curve: 0.029, 0.014 to 0.051, on 318 tests, where the pooled gate of 300 was simulated at the pooled event rate and not at this origin's; there the base rate is 0.091 and the event rate 0.058, so the base rate lags and part of the skill may be that lag; at April 2015, the last origin of the calendar: 83 tests, below every gate & \raggedright above zero at January 2015, with that caution, not counted as holding; undecidable at April 2015 \tabularnewline
% src: $RUN6/backtest-upworthy-20261106-2dc461f52ea4.json tables.origins[origin=2015-04-01T00:00:00+00:00].n_tests; $RUN6/paper-extras-upworthy-20261106-2dc461f52ea4.json gates_by_stratum.strata.ALL.own_gate.n_400
\raggedright H2, the gate & \raggedright the released gate decidable at the final origin & \raggedright the last origin of the calendar, April 2015, holds 83 tests, below the gate of 300 for a skill of 0.10 and below every other gate & \raggedright failed \tabularnewline
\midrule
\multicolumn{4}{@{}l}{\textbf{H3, in part} (see~\ref{sec:calibration-results})} \\
% src: $RUN6/paper-extras-upworthy-20261106-2dc461f52ea4.json issued_ece_by_origin.issued.origins_inside_central_90; $RUN6/paper-extras-upworthy-20261106-2dc461f52ea4.json issued_ece_by_origin.issued.origins; $RUN6/paper-extras-upworthy-20261106-2dc461f52ea4.json issued_ece_by_origin.issued.pooled.ece_percentile
\raggedright H3 part 1 & \raggedright ECE of issued cards inside the central 90 percent of the null at most origins & \raggedright inside at 2 of the 5 origins that issue cards, not at most of them; pooled at the 100th percentile of the clustered null & \raggedright failed \tabularnewline
% src: $RUN6/backtest-upworthy-20261106-2dc461f52ea4.json tables.aggregate[cell=ALL,rung=P1].coverage_80; $RUN6/backtest-upworthy-20261106-2dc461f52ea4.json tables.aggregate[cell=ALL,rung=P1].coverage_80_lo; $RUN6/backtest-upworthy-20261106-2dc461f52ea4.json tables.aggregate[cell=ALL,rung=P1].coverage_80_hi; $RUN6/backtest-upworthy-20261106-2dc461f52ea4.json tables.aggregate[cell=ALL,rung=P1].coverage_90; $RUN6/backtest-upworthy-20261106-2dc461f52ea4.json tables.aggregate[cell=ALL,rung=P1].coverage_90_lo; $RUN6/backtest-upworthy-20261106-2dc461f52ea4.json tables.aggregate[cell=ALL,rung=P1].coverage_90_hi; $RUN6/backtest-upworthy-20261106-2dc461f52ea4.json tables.coverage_by_se_ratio[rung=P1,decile=1].coverage_90; $RUN6/backtest-upworthy-20261106-2dc461f52ea4.json tables.coverage_by_se_ratio[rung=P1,decile=10].coverage_90
\raggedright H3 part 2 & \raggedright 80 and 90 percent coverage within 5 points of nominal & \raggedright 81 percent, 80 to 82, and 90 percent, 89 to 90, within 5 points of nominal pooled; 1.00 to 0.69 across deciles (a range, not an interval) of the ratio of realized to planned standard error & \raggedright held, read pooled; the fall across deciles is descriptive \tabularnewline
\midrule
\multicolumn{4}{@{}l}{\textbf{H4, in part} (see~\ref{sec:abstention-results})} \\
% src: $RUN6/paper-extras-upworthy-20261106-2dc461f52ea4.json issued_row_brier_skill.rungs.P2.issued.brier_skill; $RUN6/paper-extras-upworthy-20261106-2dc461f52ea4.json issued_row_brier_skill.rungs.P2.issued.brier_skill_lo; $RUN6/paper-extras-upworthy-20261106-2dc461f52ea4.json issued_row_brier_skill.rungs.P2.issued.brier_skill_hi; $RUN6/paper-extras-upworthy-20261106-2dc461f52ea4.json issued_row_brier_skill.rungs.P2.all.brier_skill
\raggedright H4 part 1 & \raggedright issued cards above the full set on skill & \raggedright issued rows 0.0004, minus 0.0035 to 0.0046, against 0.0032 for the full set & \raggedright failed \tabularnewline
% src: $RUN6/paper-extras-upworthy-20261106-2dc461f52ea4.json issued_row_brier_skill.rungs.P2.abstained.brier_skill; $RUN6/paper-extras-upworthy-20261106-2dc461f52ea4.json issued_row_brier_skill.rungs.P2.abstained.brier_skill_lo; $RUN6/paper-extras-upworthy-20261106-2dc461f52ea4.json issued_row_brier_skill.rungs.P2.abstained.brier_skill_hi; $RUN6/paper-extras-upworthy-20261106-2dc461f52ea4.json abstention_composition.P2.abstained.skill; $RUN6/paper-extras-upworthy-20261106-2dc461f52ea4.json abstention_composition.P2.abstained.lo; $RUN6/paper-extras-upworthy-20261106-2dc461f52ea4.json abstention_composition.P2.abstained.hi; $RUN6/paper-extras-upworthy-20261106-2dc461f52ea4.json abstention_composition.P2.issued.skill; $RUN6/paper-extras-upworthy-20261106-2dc461f52ea4.json abstention_composition.P2.issued.lo; $RUN6/paper-extras-upworthy-20261106-2dc461f52ea4.json abstention_composition.P2.issued.hi; $RUN6/paper-extras-upworthy-20261106-2dc461f52ea4.json abstention_composition.P2.pooled_issued_minus_abstained_skill; $RUN6/paper-extras-upworthy-20261106-2dc461f52ea4.json abstention_composition.P2.stratified_issued_minus_abstained_skill.estimate; $RUN6/paper-extras-upworthy-20261106-2dc461f52ea4.json abstention_composition.P2.stratified_issued_minus_abstained_skill.lo; $RUN6/paper-extras-upworthy-20261106-2dc461f52ea4.json abstention_composition.P2.stratified_issued_minus_abstained_skill.hi
\raggedright H4 part 2 & \raggedright abstained rows lower on skill when forced & \raggedright abstained rows, scored on the withheld forecast: with the burn-in origin, above the full set at 0.008, 0.001 to 0.015; without it, minus 0.0013, minus 0.0060 to 0.0041, against 0.0004, minus 0.0035 to 0.0045 for issued rows, no measurable difference & \raggedright failed on all rows; with burn-in excluded, on the claim's side pooled with no interval and on the wrong side within origin and cell \tabularnewline
% src: $RUN6/backtest-upworthy-20261106-2dc461f52ea4.json tables.abstention[cell=ALL].rate; $RUN6/backtest-upworthy-20261106-2dc461f52ea4.json tables.abstention[cell=ALL].rate; $RUN6/backtest-upworthy-20261106-2dc461f52ea4.json tables.abstention[cell=ALL].n; $RUN6/backtest-upworthy-20261106-2dc461f52ea4.json tables.origins[origin=2013-04-01T00:00:00+00:00].n_scored; $RUN6/backtest-upworthy-20261106-2dc461f52ea4.json tables.origins[origin=2013-04-01T00:00:00+00:00].abstention_rate
\raggedright H4 part 3 & \raggedright abstention below 30 percent after burn-in & \raggedright 32 percent pooled and 27.2 percent after the burn-in origin, January 2014 included & \raggedright held \tabularnewline
\midrule
\multicolumn{4}{@{}l}{\textbf{H5, in part} (Sections~\ref{sec:hypotheses} and~\ref{sec:decision-results})} \\
% src: $RUN6/backtest-upworthy-20261106-2dc461f52ea4.json tables.decision_value[origin=ALL,budget=0.25,policy=p_win,rung=text_lr].diff_matched; $RUN6/backtest-upworthy-20261106-2dc461f52ea4.json tables.decision_value[origin=ALL,budget=0.25,policy=p_win,rung=text_lr].diff_matched_adjusted_lo; $RUN6/backtest-upworthy-20261106-2dc461f52ea4.json tables.decision_value[origin=ALL,budget=0.25,policy=p_win,rung=text_lr].diff_matched_adjusted_hi
\raggedright H5, text classifier & \raggedright more positive lift than random at K = 25 percent of tests & \raggedright against the arm-count-matched random order, plus 0.046, adjusted interval above zero & \raggedright held against the matched order \tabularnewline
% src: $RUN6/backtest-upworthy-20261106-2dc461f52ea4.json tables.decision_value[origin=ALL,budget=0.25,policy=p_win,rung=GBM].diff_matched; $RUN6/backtest-upworthy-20261106-2dc461f52ea4.json tables.decision_value[origin=ALL,budget=0.25,policy=p_win,rung=GBM].diff_matched_adjusted_lo; $RUN6/backtest-upworthy-20261106-2dc461f52ea4.json tables.decision_value[origin=ALL,budget=0.25,policy=p_win,rung=GBM].diff_matched_adjusted_hi; $RUN6/backtest-upworthy-20261106-2dc461f52ea4.json tables.decision_value[origin=ALL,budget=0.25,policy=p_win,rung=GBM].diff_random_adjusted_lo; $RUN6/backtest-upworthy-20261106-2dc461f52ea4.json tables.decision_value[origin=ALL,budget=0.25,policy=p_win,rung=GBM].diff_random_adjusted_hi
\raggedright H5, boosted model & \raggedright the same & \raggedright against the matched order, plus 0.017, adjusted interval spanning zero; against the random order the plan names, adjusted 0.001 to 0.052 & \raggedright not shown against the matched order; held against the plan's random order \tabularnewline
% src: $RUN6/backtest-upworthy-20261106-2dc461f52ea4.json tables.decision_value[origin=ALL,budget=0.25,policy=p_win,rung=P2].diff_matched; $RUN6/backtest-upworthy-20261106-2dc461f52ea4.json tables.decision_value[origin=ALL,budget=0.25,policy=p_win,rung=P2].diff_matched_adjusted_lo; $RUN6/backtest-upworthy-20261106-2dc461f52ea4.json tables.decision_value[origin=ALL,budget=0.25,policy=p_win,rung=P2].diff_matched_adjusted_hi; $RUN6/backtest-upworthy-20261106-2dc461f52ea4.json tables.decision_value[origin=ALL,budget=0.25,policy=p_win,rung=P2].diff_random_adjusted_lo; $RUN6/backtest-upworthy-20261106-2dc461f52ea4.json tables.decision_value[origin=ALL,budget=0.25,policy=p_win,rung=P2].diff_random_adjusted_hi
\raggedright H5, precedent rung & \raggedright the same & \raggedright against the matched order, plus 0.018, adjusted interval spanning zero; against the random order the plan names, adjusted 0.004 to 0.045 & \raggedright not shown against the matched order; held against the plan's random order \tabularnewline
% src: $RUN6/backtest-upworthy-20261106-2dc461f52ea4.json tables.decision_value[origin=ALL,budget=0.25,policy=p_win,rung=local_eb].diff_matched; $RUN6/backtest-upworthy-20261106-2dc461f52ea4.json tables.decision_value[origin=ALL,budget=0.25,policy=p_win,rung=local_eb].diff_matched_adjusted_lo; $RUN6/backtest-upworthy-20261106-2dc461f52ea4.json tables.decision_value[origin=ALL,budget=0.25,policy=p_win,rung=local_eb].diff_matched_adjusted_hi; $RUN6/backtest-upworthy-20261106-2dc461f52ea4.json tables.decision_value[origin=ALL,budget=0.25,policy=p_win,rung=local_eb].diff_random_adjusted_lo; $RUN6/backtest-upworthy-20261106-2dc461f52ea4.json tables.decision_value[origin=ALL,budget=0.25,policy=p_win,rung=local_eb].diff_random_adjusted_hi
\raggedright H5, local comparison set & \raggedright the same & \raggedright against the matched order, plus 0.018, adjusted interval spanning zero; against the random order the plan names, adjusted 0.002 to 0.056 & \raggedright not shown against the matched order; held against the plan's random order \tabularnewline
% src: $RUN6/backtest-upworthy-20261106-2dc461f52ea4.json tables.decision_value[origin=ALL,budget=0.25,policy=p_win,rung=P1].diff_matched; $RUN6/backtest-upworthy-20261106-2dc461f52ea4.json tables.decision_value[origin=ALL,budget=0.25,policy=p_win,rung=P1].diff_matched_lo; $RUN6/backtest-upworthy-20261106-2dc461f52ea4.json tables.decision_value[origin=ALL,budget=0.25,policy=p_win,rung=P1].diff_matched_hi
\raggedright H5, prior & \raggedright the same & \raggedright at the matched order within noise, minus 0.004, minus 0.014 to 0.007 & \raggedright failed \tabularnewline
\midrule
\multicolumn{4}{@{}l}{\textbf{H6, held} (see~\ref{sec:asos-results})} \\
% src: $RUN6/backtest-asos-20261103-6bb5a15c5ff0.json tables.aggregate[cell=ALL,rung=P2].n_tests
\raggedright H6 part 1 & \raggedright every gate undecidable on the retailer's archive & \raggedright undecidable at every gate & \raggedright held, a result 68 tests could not have avoided \tabularnewline
\raggedright H6 part 2 & \raggedright P2 equals P1 & \raggedright P2 equals P1 on this archive, because the guard holds the weight at zero & \raggedright held \tabularnewline
\bottomrule
\end{tabular}

\end{table*}

\paragraph{What held and what did not.} The prior beat the base rate on the log score at every origin after burn-in, abstention stayed under its limit after burn-in, the text classifier captured more lift than the matched order, and the retailer's archive stayed undecidable. Pooled skill at realized precision was not shown, calibration of issued cards failed, issued cards did not score above the full set on Brier skill, and the prior did not rank tests. A precedent gain in the largest cell was not shown, and the backtest could not have shown one of the size a planted control returns. Three parts hold for reasons that limit what they say, as the notes below explain.

\paragraph{Notes on the table.}
\begin{itemize}
\item H1 part 2: the interval contains what a planted signal of that size returns, so this result does not refute a precedent signal; it says this backtest cannot see one (Section~\ref{sec:synthetic-results}).
\item H1 part 3: P2 equals P1 in the image-only cell because the guard admitted no weight there, so the part holds by the guard and says nothing about what precedents carry in that cell.
% src: $RUN6/learning-curve-upworthy-20261106-2dc461f52ea4.json rows[fraction=1.0,rung=P2] skill 0.029 (0.014, 0.051); $RUN6/paper-extras-upworthy-20261106-2dc461f52ea4.json drift_matched_null.rungs.P0.by_origin["2015-01-01T00:00:00+00:00"] mean_forecast 0.091, event_rate 0.058, drift_matched_null.rungs.P2.by_origin["2015-01-01T00:00:00+00:00"] mean_forecast 0.066; within_origin_auc_by_year.years["2015"].P2.auc_within_origin 0.634 (0.568, 0.703)
\item H2: the plan says ``at the final origin'', and the verdict depends on the reading. Pooled, under the rule of Section~\ref{sec:hypotheses}, it is not shown.
% src: $RUN6/backtest-upworthy-20261106-2dc461f52ea4.json tables.aggregate[cell=ALL,rung=P2] brier_skill_planned 0.0058 (0.0018, 0.0095)
With the event at planned precision instead, P2's pooled interval lies above zero, unadjusted, at 0.0058, so this reading depends on how the event is defined and the skill is at most small either way (Section~\ref{sec:ladder-results}).
As written it holds at January 2015, the last origin with enough tests to clear a gate, and the last calendar origin is undecidable. At January 2015 the base rate of 0.091 stands above an event rate of 0.058 and P2's mean forecast of 0.066 sits near the event rate, so part of that skill may be the base rate's lag rather than foresight about any one test (Table~\ref{tab:fmt-auc-by-year} gives the ranking within that year).
\item H3: the ECE percentile rejects calibration, and Section~\ref{sec:calibration-results} shows that most of the error is a level the forecast does not anticipate.
\item H4: the hypothesis names skill. With the burn-in origin the abstained rows score above the full set on it, and without it the abstained and issued rows do not differ measurably (Table~\ref{tab:fmt-abstention-composition}). On the log score, which no part names, issued rows are better (Section~\ref{sec:abstention-results}). The 27.2 percent after burn-in includes January 2014, which abstains on every row because the date rule leaves its reliability reports stale.
\item H5: the prior ranks by cell only, its win probability one number per cell and origin. The verdict rests on the arm-count-matched order, adopted after the data, and regret against the oracle is reported as is in Section~\ref{sec:decision-results}.
\item H6: 68 tests against gates of hundreds could not have come out otherwise, and P2 equals P1 only because the guard holds the weight at zero. This hypothesis confirms the protocol's behaviour on a small archive and is no evidence about product experiments.
\end{itemize}

\section{Discussion}
\label{sec:discussion}
The answers below hold for one publisher's archive read in an exploratory way (Section~\ref{sec:limitations}). What can transfer to a team's own experiments is the protocol and not these numbers.

\paragraph{RQ1: can a team predict a new experiment's result from its concluded experiments?} Partly. Past tests predict the distribution of the next result better than a reference that uses no past effect, and a heavy-tailed base rate does this nearly as well as the fitted prior. They predict at most a small part of which test will be significant.
\begin{itemize}
% src: $RUN6/paper-extras-upworthy-20261106-2dc461f52ea4.json p1_minus_p0_by_origin.pooled difference 0.120 (0.106, 0.136), by_origin_clear_of_zero 7 (the seven scored origins; the first, 2013-04-01, is abstention_composition.burn_in_origin); tables/verdicts.tex H1 part 1; $RUN6/ablations/backtest-upworthy-20261106-2dc461f52ea4--prior-normal.json tables.aggregate[cell=ALL,rung=P1] logscore -0.8316 against $RUN6/backtest-upworthy-20261106-2dc461f52ea4.json tables.aggregate[cell=ALL,rung=P1] logscore -0.7566 (difference 0.0750, the heavy-tailed shape); $RUN6/ablations/backtest-upworthy-20261106-2dc461f52ea4--prior-normal.json tables.aggregate[cell=ALL,rung=P0] logscore_vs_p1 -0.0446 (-0.0548, -0.0336), negated for P1 minus P0 with both densities normal: 0.045 (0.034, 0.055); $RUN6/ablations/backtest-upworthy-20261106-2dc461f52ea4--base-density-t.json tables.aggregate[cell=ALL,rung=P0].logscore_vs_p1 -0.0196 (-0.0232, -0.0155), negated: P1 minus the Student-t base rate 0.020 (0.015, 0.023) (tables/fmt-base-density.tex); $RUN6/ablations/backtest-upworthy-20261106-2dc461f52ea4--base-density-t.json tables.score_by_origin (P1 above P0 at four origins, below at three); $RUN6/ablations/backtest-upworthy-20261106-2dc461f52ea4--exposure-cell-median-visible+base-density-t.json tables.aggregate[cell=ALL,rung=P0].logscore_vs_p1 -0.0069 (-0.0106, -0.0030), negated: P1 minus the Student-t base rate with exposure fixed on visible rows 0.007 (0.003, 0.011); image_only -0.0007 (-0.0081, 0.0053), both -0.0064 (-0.0308, 0.0181) (tables/fmt-base-density.tex); $RUN6/ablations/by-origin-upworthy-20261106-2dc461f52ea4--exposure-cell-median-visible+base-density-t.json leave_one_origin_out[left_out=2014-07-01] difference 0.0008 (-0.0044, 0.0060), leave_one_origin_out_above_zero 6 of 7, by_origin_above_zero 3, by_origin_below_zero 1, origin_level difference 0.0036 (-0.0120, 0.0191); $RUN6/ablations/by-origin-upworthy-20261106-2dc461f52ea4--base-density-t.json leave_one_origin_out_above_zero 7 of 7, by_origin_above_zero 4, by_origin_below_zero 1 (tables/fmt-by-origin.tex); $RUN6/ablations/by-origin-upworthy-20261106-2dc461f52ea4--base-density-t.json origin_level difference 0.0121 (-0.0060, 0.0302), origins 7 (a Student-t interval on the seven origin means, the quarters as units; it contains zero); $RUN6/by-origin-upworthy-20261106-2dc461f52ea4.json origin_level difference 0.1053 (0.0563, 0.1542) (the normal base rate: above zero with the quarters as units)
\item Against a base rate that is also given heavy tails, the prior's pooled lead on the log score is 0.020 nats, and 0.007 nats once exposure is fixed on visible rows, a setting that changes only the prior (Table~\ref{tab:fmt-base-density}). Both pooled intervals are above zero, and both resample tests with the seven scored quarters held fixed, so they describe these rows. With the quarters as the units the interval contains zero for each, so the prior's lead over a heavy-tailed base rate is not shown; seven values are few, and that interval treats them as exchangeable (Table~\ref{tab:fmt-by-origin}). By origin the paired interval against the Student-t base rate is above zero at four origins and below zero at one. With exposure fixed the pooled lead rests on July 2014: without that origin its interval spans zero. Against the base rate as first specified, scored with a normal density, the prior leads by 0.120 nats pooled (Table~\ref{tab:glance}), and that lead holds at each origin and with quarters as units (Tables~\ref{tab:verdicts} and~\ref{tab:fmt-by-origin}). Most of that gap is the shape of the density: with both densities normal the prior's gain is 0.045 nats (Table~\ref{tab:fmt-prior-settings}).
% src: $RUN6/learning-curve-upworthy-20261106-2dc461f52ea4.json rows[fraction=0.01,rung=P1] n_history_tests 26; paired_vs_full[fraction=0.01,rung=P1] logscore_minus_full -0.0829 (-0.1086, -0.0592); final_origin 2015-01-01 (one origin); paired_vs_full[fraction=0.01,rung=P0] logscore_minus_full -0.1232, and P0 below P1 at each of the six fractions, point estimates with no paired interval between the two rungs; tables/fmt-learning-vs-full.tex (the comparison was added after a review, Section 7); $RUN6/ablations/backtest-upworthy-20261106-2dc461f52ea4--base-density-t.json tables.aggregate[cell=ALL,rung=P0].logscore_vs_p1 -0.0196, negated 0.020; $RUN6/ablations/backtest-upworthy-20261106-2dc461f52ea4--exposure-cell-median+base-density-t.json tables.aggregate[cell=ALL,rung=P0].logscore_vs_p1 -0.0085, negated 0.009; $RUN6/paper-extras-upworthy-20261106-2dc461f52ea4.json exposure_fixed.p_sig_constant_within_origin_and_cell.fixed exposure_strata_with_one_value 1.0 (with exposure fixed the planned standard error does not vary within a cell and origin; tables/fmt-base-density.tex); $RUN6/no-history-upworthy-20261106-2dc461f52ea4.json P0.minus_reference difference 0.7066 (0.5779, 0.8680), P0.by_origin_above_zero 7, P1.minus_reference difference 0.8263 (0.6922, 0.9928), P1.by_origin_above_zero 7; $RUN6/ablations/no-history-upworthy-20261106-2dc461f52ea4--base-density-t.json P0.minus_reference difference 0.8067 (0.6730, 0.9761), P0.by_origin_above_zero 7 (tables/fmt-no-history.tex; the reference is Normal(0, s_star^2) and reads no ledger)
\item Both forecasters are fitted to the same history, so the differences above do not measure what history adds. One comparison does, against a floor: a reference that uses no past effect and assumes no test has one. The base rate and the prior both beat it at every origin, and nearly all of that gain is already in the base rate (Table~\ref{tab:fmt-no-history}). The reference denies that tests differ, so it is not a serious opponent, and a fixed heavy-tailed density chosen without the archive was not scored. What more history adds is measured separately, at one origin, in a comparison added after a review. There, P1 fitted on about 26 tests scores 0.083 nats per row below P1 fitted on the full history, with the interval clear of zero (Table~\ref{tab:fmt-learning-vs-full}). On point estimates the base rate loses more on the same subsamples, so that gain is not specific to the prior.
\item Its skill on which test will be significant is at most small, and it depends on how that event is defined.
% src: $RUN6/backtest-upworthy-20261106-2dc461f52ea4.json tables.aggregate[cell=ALL,rung=P1] brier_skill 0.0023 (-0.0012, 0.0061), n_tests 2949; tables.score_by_year[cell=ALL,rung=P1] year 2014 n_tests 2303, brier_skill -0.0036, adjusted (-0.0102, 0.0030); year 2013 brier_skill 0.0334, adjusted (-0.0073, 0.0742); year 2015 brier_skill 0.0243, adjusted (-0.0040, 0.0525) (tables/fmt-drift-year.tex)
\item With the event at realized precision, the main reading, the prior's pooled skill is 0.0023 with an interval spanning zero. That interval is the evidence. By year the estimate is below zero in the year that holds most of the tests, 2,303 of the 2,949. It is above zero in the first and the last scored year, and every adjusted interval by year spans zero (Table~\ref{tab:fmt-drift-year}).
% src: $RUN6/paper-extras-upworthy-20261106-2dc461f52ea4.json gates_by_stratum.strata.ALL.own_gate.n_400 300 (the gate for a skill of 0.10; src/mootcourt/card.py prints n_400 as the gate for a skill of 0.10 and n_200 for 0.20); $RUN6/paper-extras-upworthy-20261106-2dc461f52ea4.json gates_by_stratum.strata.ALL.dependence.n_tests 2949, rungs.P1.verdict "decidable at skill 0.20 and 0.10"; Section 6.2, "Decidable does not mean skilled" (the gate is sized for a skill of 0.10; the pooled skills read here are below 0.01)
The archive clears the gate for a skill of 0.10. But that gate assumes calibration, which the prior fails, and it is not sized for the far smaller skills seen here (Section~\ref{sec:gates-results}), so an interval that spans zero does not show that the skill is absent.
% src: $RUN6/backtest-upworthy-20261106-2dc461f52ea4.json tables.aggregate[cell=ALL] brier_skill_planned rung=P1 0.0046 (0.0008, 0.0083), rung=P2 0.0058 (0.0018, 0.0095); $RUN6/include-randomization-window/backtest-upworthy-20261106-4bfb354b87e3.json tables.aggregate[cell=ALL] brier_skill_planned rung=P1 -0.0002 (-0.0037, 0.0038), rung=P2 0.0004 (-0.0037, 0.0045) (tables/fmt-window.tex)
\item With the event at planned precision the unadjusted intervals of the prior and P2 lie above zero in the main run. With the flagged randomization window included both span zero, so this reading does not survive switching the date rule off (Table~\ref{tab:fmt-window}; Section~\ref{sec:limitations}). It does not survive a base rate that is given the prior's quarter term either (Table~\ref{tab:fmt-base-recency}).
% src: $RUN6/paper-extras-upworthy-20261106-2dc461f52ea4.json exposure_fixed.auc_event_within_ties_exact main.P1 0.525, fixed.P1 0.5; exposure_fixed.p_sig_constant_within_origin_and_cell fixed.P1 strata_with_one_value 21, origin_cell_strata 21 (tables/fmt-exposure-fixed.tex)
\item Its small ability to rank tests within an origin and cell comes entirely from the exposure stand-in: with exposure fixed, here in the whole-ledger diagnostic, its forecast is constant within each origin and cell and the AUC falls from 0.525 to 0.500 (Table~\ref{tab:fmt-exposure-fixed}).
% src: $RUN6/paper-extras-upworthy-20261106-2dc461f52ea4.json exposure_fixed.by_cell.ALL.fixed.P1 brier_skill 0.0059 (0.0031, 0.0090), against main.P1 brier_skill 0.0023 (-0.0012, 0.0061); exposure_fixed.by_cell.ALL.fixed.P2 brier_skill 0.0071 (0.0043, 0.0105), against main.P2 brier_skill 0.0032 (-0.0004, 0.0072); the report carries no adjusted bound for these; $RUN6/ablations/backtest-upworthy-20261106-2dc461f52ea4--exposure-cell-median-visible.json tables.aggregate[cell=ALL] rung=P1 brier_skill 0.0069 (0.0041, 0.0100), rung=P2 brier_skill 0.0081 (0.0051, 0.0114)
Brier skill rises in that run, and in the run fixed on visible rows, for the prior and for P2, each with an unadjusted interval above zero, which Table~\ref{tab:fmt-assurance} may explain and which we did not test. So the pooled reading of skill on significance depends on the exposure stand-in as well as on the event's definition.
% src: $RUN6/paper-extras-upworthy-20261106-2dc461f52ea4.json drift_matched_null.rungs.P1.pooled.current_null.percentile 100.0; calibration_within_period.rungs.P1.by_stratum_group.ALL share_of_ece_removed 0.70; calibration_within_period.rungs.P1.strata_shifted 18 (of the 21 pairs of origin and cell)
\item As issued its forecasts fail the calibration test, at the 100th percentile of the clustered null of Section~\ref{sec:protocol}, which draws the outcomes of one test together. A level shift fitted after the fact to each origin and cell with enough rows removes about 70 percent of the error. That suggests an error of level and does not test one: the shift reads the realized rate, and its constants are fitted on the rows they are judged on (Section~\ref{sec:calibration-results}).
\item On the archive of product experiments the answer is unknown, because that archive is undecidable at every gate.
\end{itemize}

\paragraph{RQ2: how, and how well?} A heavy-tailed density fitted per cell to earlier results gives most of the log-score gain measured here, and a prior fitted with measurement error adds little, an addition that is not shown once quarters are the units. The method is to replay in time order and to read each interval. Beyond the prior, the evidence is as follows.
\begin{itemize}
% src: $RUN6/backtest-upworthy-20261106-2dc461f52ea4.json tables.aggregate[cell=ALL,rung=P2] logscore_vs_p1 0.0001 (-0.0004, 0.0006), brier_skill 0.0032 (-0.0004, 0.0072); $RUN6/paper-extras-upworthy-20261106-2dc461f52ea4.json guard_pairs.summary pairs_admitted 4, pairs_in_selection 21
\item Precedents stand level with the prior on the pooled log score, where the paired difference spans zero (Section~\ref{sec:ladder-results}). On Brier skill each rung's interval against the base rate spans zero, and the main backtest holds no paired interval between the two. In most pairs of origin and cell the guard admits no weight and the rung is the prior itself, so part of that tie is by construction. In a comparison that is not matched on arm count, they rank tests a little better than the prior.
% src: $RUN6/paper-extras-upworthy-20261106-2dc461f52ea4.json ranking_paired.budgets["0.25"].differences.P2_minus_P1 estimate 0.026, adjusted_lo 0.016, adjusted_hi 0.037; budgets "0.10" adjusted (0.011, 0.028) and "0.50" adjusted (0.003, 0.020) (tables/fmt-ranking-paired.tex)
Paired by test, the precedent rung captures 0.026 more of the available lift than the prior at a quarter of the tests, and the adjusted intervals clear zero at a tenth, a quarter and a half (Table~\ref{tab:fmt-ranking-paired}).
% src: $RUN6/paper-extras-upworthy-20261106-2dc461f52ea4.json ranking_paired.budgets["0.25"].differences.P2_minus_s_star_desc -0.0104 (-0.0346, 0.0125), budgets["0.50"].differences.P2_minus_s_star_desc -0.0295 (-0.0563, -0.0016), unadjusted (adjusted_comparisons lists only the two differences from P1); differences.P2_minus_arm_count budgets["0.25"] (-0.0554, -0.0149), budgets["0.50"] (-0.0830, -0.0375); Section 6.5 (the least-precise-first direction was added after the data)
The edge may be the number of arms or traffic, and not precedents: P2 captures less than a rule that only counts arms at a quarter and at a half of the tests, and no more than taking the least precise tests first, a direction added after the data, where at a half the unadjusted interval of the difference lies just below zero (Section~\ref{sec:decision-results}).
% src: $RUN6/paper-extras-synthetic-20261103-bf4fd286e2b2.json positive_control_rungs.runs[n_tests=11100].by_cell["ALL"].rungs[rung=P2].recovered_fraction_logscore 0.09, by_cell["headline_only"].rungs[rung=P2].recovered_fraction_logscore 0.03 (inputs.seed 0, one ledger of 11,100 tests; the fractions carry no interval keys); by_cell["ALL"].oracle gain_logscore_vs_p1 0.0197 (0.0177, 0.0219), which the backtest reads with an interval clear of zero
\item The precedent rung recovers little of a planted signal. In a positive control under the same rolling design, the scoring detects a model told the signal, and P2 recovers 0.09 of that model's log-score gain pooled and 0.03 in the headline-only cell. Both figures come from one simulated ledger and carry no interval. So an interval around zero does not refute a precedent signal of that size, and it does not show that the rung could use one (Section~\ref{sec:synthetic-results}).
% src: $RUN6/backtest-upworthy-20261106-2dc461f52ea4.json tables.decision_value[origin=ALL,budget=0.25,policy=p_win,rung=text_lr] diff_matched 0.046, diff_matched_adjusted_lo 0.019, diff_matched_adjusted_hi 0.073; tables.decision_value[origin=ALL,budget=0.25,policy=arm_count] captured_fraction 0.3093 against policy=p_win rung=text_lr 0.3070; $RUN6/paper-extras-upworthy-20261106-2dc461f52ea4.json ranking_paired.budgets["0.25"].differences.text_lr_minus_arm_count -0.0023 (-0.0268, 0.0221), budgets["0.10"] (-0.0079, 0.0247), budgets["0.50"] (-0.0474, 0.0031); differences.P2_minus_arm_count budgets["0.10"] (-0.0249, 0.0008), budgets["0.25"] (-0.0554, -0.0149), budgets["0.50"] (-0.0830, -0.0375) (tables/fmt-ranking-paired.tex)
\item The text classifier ranks best among the models, at 0.046 of captured lift over the arm-count-matched random order with an adjusted interval clear of zero. It is untuned, and the number comes from one split. Against a rule that ranks by arm count and reads no forecast, the classifier's paired interval spans zero at all three budgets (Table~\ref{tab:fmt-ranking-paired}). So the classifier's value is read only against the matched order.
% src: $RUN6/backtest-upworthy-20261106-2dc461f52ea4.json tables.aggregate[cell=ALL,rung=text_lr] brier_skill 0.0070 (-0.0005, 0.0140), brier_skill_planned 0.0159 (0.0082, 0.0238); $RUN6/include-randomization-window/backtest-upworthy-20261106-4bfb354b87e3.json tables.aggregate[cell=ALL,rung=text_lr] brier_skill_planned 0.0037 (-0.0027, 0.0098) (tables/fmt-window.tex); header.challengers empty in the three base-recency reports
\item The text classifier's skill on significance has the same dependence: its interval spans zero at realized precision and lies above zero, unadjusted, at planned precision in the main run. With the randomization window included it spans zero (Table~\ref{tab:fmt-window}), and the classifier was not scored against a base rate with a quarter term (Section~\ref{sec:limitations}).
% src: $RUN6/paper-extras-synthetic-20261103-bf4fd286e2b2.json positive_control_decision_value (tables/fmt-ranking-paired.tex: planted ledgers, text classifier minus P1 -0.0119 and -0.0199)
\item The text classifier fails the planted control on skill and ranks below P1 there, so the control neither validates nor explains its value on the real archive, and what it reads there is left open.
\item Abstention is not shown to improve what is issued on this archive, so it stays as a label that names which trigger fired and claims nothing about the forecast's quality (Section~\ref{sec:abstention-results}). The card prints beside every forecast the latest gated skill of its cell, or the pooled figure, labelled, when the cell is undecidable.
% src: $RUN6/paper-extras-upworthy-20261106-2dc461f52ea4.json sizing_evaluation.test_and_roll.by_audience["20"].paired_differences.prior_minus_base_rate estimate -0.0016, adjusted_lo -0.0042, adjusted_hi 0.0010 (tables/fmt-sizing.tex)
\item Sizing was evaluated on a synthetic horizon, with no money involved (Table~\ref{tab:fmt-sizing}).
% src: $RUN6/paper-extras-upworthy-20261106-2dc461f52ea4.json sizing_evaluation.test_and_roll primary_multiple 20, by_audience["20"].paired_differences.prior_minus_base_rate estimate -0.0016, adjusted_lo -0.0042, adjusted_hi 0.0010; prior_minus_power adjusted (-1.349, -1.083), base_rate_minus_power -1.214 (-1.314, -1.102), relative regret, lower is better; tables/fmt-sizing.tex fixed power rule 1.568 at the primary multiple, equal to test the whole audience, against never test 1.000
At the primary horizon, an audience of 20 times the size of the cell's median two-arm test, the prior's size has no measurable advantage over the base rate's. Both beat a rule sized for a fixed power, but only because at that horizon the rule hits its cap and tests the whole audience, which loses more than never testing (Section~\ref{sec:sizing-results}). Against the base rate the sign changes with the horizon and with the assumed world. The paper does not claim that sizing from the forecast improves on a base rate.
\end{itemize}

\paragraph{RQ3: what improves the prediction?} Among what was tested, a few changes matter and most settings do not.
\begin{itemize}
% src: $RUN6/ablations/backtest-upworthy-20261106-2dc461f52ea4--prior-normal.json tables.aggregate[cell=ALL,rung=P1] logscore -0.8316 against the main run's -0.7566 (0.075 nats); $RUN6/backtest-upworthy-20261106-2dc461f52ea4.json tables.aggregate[cell=ALL,rung=P0].logscore -0.8762 against the base-density-t report's -0.7761 (the base rate gains more than 0.075 from the Student-t; tables/fmt-base-density.tex)
\item A heavy-tailed density gains most, on either forecaster. It is worth 0.075 nats to the prior, and more than that to the base rate (Tables~\ref{tab:fmt-prior-settings} and~\ref{tab:fmt-base-density}). So this is a statement about the shape of results on this scale and not about the prior's design. An in-sample diagnostic, which reads a rate known only after launch, suggests that much of what the Student-t captures is a mixture over the click-through level. It does not split that gain, and we make no claim that effects themselves are heavy-tailed (Table~\ref{tab:fmt-prior-family}).
\item The untuned text classifier's ranking on this one split is a lead and not a change shown to improve prediction, for the reasons RQ2 gives.
\item Anticipating the level would improve calibration, if a forecast could do it. The diagnostic of RQ1 shows the size of that error and uses the realized rate, so no method in this paper delivers the gain.
% src: $RUN6/ablations/backtest-upworthy-20261106-2dc461f52ea4--lambda-1.0.json tables.aggregate[cell=ALL,rung=P2] logscore_vs_p1 -0.01303
\item A precedent weight fixed at one makes the score worse, by 0.013 nats against the prior.
% src: $RUN6/ablations/backtest-upworthy-20261106-2dc461f52ea4--p2-guard-off.json tables.aggregate[cell=ALL,rung=P2] logscore_vs_p1 -0.00386; $RUN6/ablations/backtest-upworthy-20261106-2dc461f52ea4--p2-selection-loto+p2-guard-off.json tables.aggregate[cell=ALL,rung=P2] logscore_vs_p1 0.00015 (-0.00110, 0.00148)
Without the guard the selected weight still costs 0.004 nats under the default selection rule, and the guard removes that cost. Under leave-one-test-out selection P2's paired difference from the prior spans zero with the guard off, so the guard is needed by the default rule and not by every rule (Section~\ref{sec:ablations}).
% src: $RUN6/ablations/backtest-upworthy-20261106-2dc461f52ea4--features-text.json tables.aggregate[cell=ALL,rung=P2] logscore_vs_p1 -0.00119 (-0.00191, -0.00052); $RUN6/ablations/backtest-upworthy-20261106-2dc461f52ea4--kappa-0.0.json tables.aggregate[cell=ALL,rung=P1] logscore -0.7607 and $RUN6/ablations/backtest-upworthy-20261106-2dc461f52ea4--kappa-80.0.json -0.7549, the two ends of the grid (tables/fmt-prior-settings.tex prints -0.761 and -0.755; the report holds no paired interval between settings of kappa)
\item Bandwidth, recency and the selection rule barely move P2's score, and text features alone put it below the prior with an interval clear of zero (Table~\ref{tab:fmt-p2-settings}). Pooling strength moves the prior's log score from minus 0.761 to minus 0.755 over its grid. That is not small beside the prior's lead over a heavy-tailed base rate, and no paired interval compares the settings (Table~\ref{tab:fmt-prior-settings}).
% src: $RUN6/learning-curve-upworthy-20261106-2dc461f52ea4.json rows rung=P1 skill 0.024 at n_fit_rows 2369.2 and at 9447 (tables/fmt-learning.tex); paired_vs_full[fraction=0.25,rung=P1] brier_gain_over_full -0.000013 (-0.000069, 0.000053), logscore_minus_full -0.0051 (-0.0075, -0.0027); paired_vs_full[fraction=0.5,rung=P1] brier_gain_over_full -0.000008 (-0.000050, 0.000045), logscore_minus_full -0.0013 (-0.0025, -0.0003), logscore_seed_sd 0.0022, above the mean gap, and the interval holds the subsamples fixed (tables/fmt-learning-vs-full.tex)
\item Against the full history at one origin, the prior's paired Brier difference spans zero from a quarter of the history on, which is less than showing that the gain ends there, and on the log score the mean gap at half of it is small and not established (Tables~\ref{tab:fmt-learning} and~\ref{tab:fmt-learning-vs-full}). The skill curve is scored against the full history's base rate, which is stale at that origin, so it does not show how much skill the prior has (Section~\ref{sec:ablations}).
\item Not tested, so these are candidates and not findings: recording planned exposure, stronger text models, choosing the width of the precedent rung's forecast on held-out rows, and team forecasts frozen before launch.
\end{itemize}

%% tmlr:move-begin production
%% tmlr:note production Notes on using the package in production, where it has not been deployed, are in Appendix~\ref{sec:production}.
%% tmlr:title production Using this in production
\label{sec:production}
\paragraph{Using this in production.} This system has never been deployed in a production experimentation platform, and the evidence is a replay. What it offers a team is a way to test forecasts on its own ledger. It needs only aggregates, and the ledger never leaves the organization. The order of use is fixed.
\begin{enumerate}
% src: $RUN6/paper-extras-upworthy-20261106-2dc461f52ea4.json gates_by_stratum.strata.ALL rungs.P1.verdict "decidable at skill 0.20 and 0.10" (the two skills the gates are sized for); $RUN6/backtest-upworthy-20261106-2dc461f52ea4.json tables.aggregate[cell=ALL,rung=P1] brier_skill 0.0023 (-0.0012, 0.0061)
\item Import the concluded experiments and run the backtest on them first. That report, not this paper's numbers, shows where forecasts fail on that ledger and how many tests a claim of a stated skill needs under the gate's assumptions. It has not been shown to certify a forecast: its gates are sized for skills of 0.10 and 0.20 and assume calibration, and the skills measured here are far smaller.
% src: $RUN6/backtest-upworthy-20261106-2dc461f52ea4.json tables.decision_value[origin=ALL,policy=p_win,rung=P2] diff_matched_adjusted_lo 0.0018 at budget 0.10, -0.0026 at 0.25, -0.0203 at 0.50; $RUN6/paper-extras-upworthy-20261106-2dc461f52ea4.json ranking_paired.budgets differences.P2_minus_arm_count (-0.0249, 0.0008) at "0.10", (-0.0554, -0.0149) at "0.25", (-0.0830, -0.0375) at "0.50"
\item Where a claim is decidable, a card shows before launch what result the ledger leads one to expect. On this archive the issued probabilities failed the calibration test, so a team should read its own reliability table before acting on a card's probability. Ranking candidate tests by the card beats a random order matched on arm count at a tenth of the tests only, after adjustment. A rule that only counts arms captured more lift than the card did, with a paired interval below zero at a quarter and a half of the tests, so a team should check ranking on its own backtest before relying on it. A card can also suggest a size, and its value over a base rate is not shown.
\item Join outcomes weekly. The card abstains on thin support, novelty, a stale or drifted reliability record or an unstable decision, which labels the evidence and is not shown to improve what is issued.
% src: $RUN6/replay-upworthy-20261106-2dc461f52ea4.json totals.n_flagged_weeks 6, header.n_weeks 111, weeks[].drifted_cells both only (tables/replay-upworthy.tex); $RUN6/backtest-upworthy-20261106-2dc461f52ea4.json tables.aggregate[cell=both,rung=P1] n_tests 173, the fewest of the three cells
On this archive the weekly loop flagged drift in 6 of 111 weeks, all in the smallest cell (Section~\ref{sec:drift-results}). The rule compares a cell's log score with the base rate's, so it may not see a change that both miss together, such as the year-to-year fall of the significant-positive rate here. A team should watch each cell's event rate beside the flag.
Anything above the prior must earn its place on a paired, gated comparison.
\end{enumerate}
Four limits apply to that use.
\begin{itemize}
\item The evidence is on randomized A/B and A/B/n tests, mainly one publisher's, and a team must treat every other design as untested (Section~\ref{sec:limitations}, Scope).
\item The protocol finds weak evidence and leakage in a backtest. It does not find defects in a change: bugs, instrumentation errors and harm remain the job of tests, guardrails and the live experiment.
\item This paper makes no promise of improvement on product experiments, because the one archive of them is undecidable.
\item A team that ranks by the forecast selects its own future ledger, so rows should record whether a card influenced the test (Section~\ref{sec:limitations}).
\end{itemize}
%% tmlr:move-end production

\section{Hypotheses, Limitations and Ethics}
\label{sec:limitations}

\paragraph{The hypotheses.}
\label{sec:hypotheses}
% src: git log of paper/plan.md: one commit, 145eaf7 (3 November 2026), which also holds paper/data/run2-prelaunch-clean/gates.json and run-log.json; the registered copy (references.bib mootcourt2026plan): tag plan-registered-2026-11-04, commit f0b088c794dcea1a1eac9872cd15aeb1c2412976, pushed 2026-11-04T17:19:12Z, plan.md sha256 e6dab9062fd1dc59734e5e4374ba429d3edc06efe681f68bb81424d88f45b613, byte-identical to paper/plan.md of 145eaf7
Six hypotheses were fixed in writing before the final runs and are unchanged since; the record is commit 145eaf7 (3 November 2026), which holds the plan beside the first run's reports; a byte-identical copy was registered publicly on 4 November 2026, after round 01 \cite{mootcourt2026plan}.
% src: verdicts read from Sections 6 and 7 and their src comments: H1 $RUN6/backtest-upworthy-20261103-4bfb354b87e3.json tables.score_by_origin (P1 above P0 at all nine origins) and tables.aggregate[cell=ALL,rung=P2] logscore_vs_p1 0.0004 (-0.0003, 0.0012); H2 tables.aggregate[cell=ALL,rung=P2] brier_skill 0.0019 (-0.0018, 0.0054); H3 ece_percentile 100, coverage_80 0.8122, coverage_90 0.8945, tables.coverage_by_se_ratio 1.000 to 0.655; H4 tables.abstention rate 0.1415 and tables.released_rows of the drop-trigger reports; H5 tables.decision_value origin=ALL budget=0.25 policy=p_win diff_matched text_lr 0.0395 (adjusted 0.0158, 0.0632), GBM 0.0276 (0.0093, 0.0459), P2 0.0139 (0.0023, 0.0264; adjusted -0.0038, 0.0316), local_eb 0.0150 (adjusted -0.0086, 0.0386), P1 -0.0043; H6 $RUN6/backtest-asos-20261103-6bb5a15c5ff0.json tables.aggregate decidable_by_tests "undecidable at this n", tables.p2_selection_by_cell share_lambda_zero 1.00
Refuted: that the prior beats the base rate at every origin and the precedent rung beats the prior where text carries signal (the first half holds; the rung equals the prior, plus 0.0004 nats within the interval, with the family held fixed), and that the precedent rung's Brier skill clears zero (0.002, minus 0.002 to 0.005). Held in part: calibration (coverage holds pooled, 81 and 89 percent, and fails within strata of the realized level, 1.00 to 0.66; the probabilities sit at the 100th percentile of the clustered null); abstention (14 percent; the triggers change what is issued, not measurably how well); and decision value, that ranking by predicted win probability captures more lift than the arm-count-matched random order at a quarter of the tests, confirmed for the text classifier and the boosted model (plus 0.040 and plus 0.028, adjusted intervals above zero), unresolved for the precedent rung and the local set (plus 0.014 and plus 0.015, adjusted intervals spanning zero), failed by construction for the prior. Confirmed: that the deployment-shaped archive is undecidable at every gate.

\paragraph{The archives and the design.} Headline tests on one publisher are one kind of change in one product, and the deployment-shaped archive is too small for any gate and carries no dates, so the protocol is measured on real experiments without transfer or real drift, from one archive of each kind, one seed and one design.

\paragraph{The scores and the stand-ins.}
% src: $RUN6/paper-extras-upworthy-20261103-4bfb354b87e3.json impressions_effect_correlation corr_log_impressions_abs_z 0.0994, n_rows 13453; $RUN6/paper-extras-asos-20261103-6bb5a15c5ff0.json corr_log_impressions_abs_z 0.2152, corr_log_run_length_abs_z 0.1522
No true effect is ever observed, so every score is against the test's own estimate, with floors and ceilings beside it; two post-launch quantities stand in for pre-launch ones. Realized impressions stand in for the planned exposure the archives do not record: their log correlates 0.099 with the absolute standardized effect on the headline archive and 0.215 on the deployment-shaped one, where the stand-in is the user count at the end of a run of varying length, a measured leak a live card does not have. % src: $RUN6/backtest-upworthy-20261103-4bfb354b87e3.json tables.aggregate[cell=ALL] event_rate_planned 0.0814, brier_base 0.0832 = p(1 - p) gives the realized event rate 0.0910; brier_skill_planned rung=P1 -0.0002 (-0.0037, 0.0038), rung=P2 -0.0033 (-0.0068, 0.0001), rung=text_lr+cal 0.0028 (-0.0011, 0.0063); $RUN6/backtest-asos-20261103-6bb5a15c5ff0.json and $RUN6/backtest-synthetic-20261103-bf4fd286e2b2.json brier_skill_planned equal to brier_skill for every rung (tables/event-definition-upworthy.tex, tables/event-definition-asos.tex, tables/event-definition-synthetic.tex)
The headline significance event and the floor ingredients use the realized standard error; Table~\ref{tab:event-definition} scores the event defined on the planned one as well: the event rate falls from 0.091 to 0.081 and every skill stays within its interval (P1 minus 0.000, minus 0.004 to 0.004), and on the other two archives the two events coincide. The scale mixture over the uncertain click-through level that might close the coverage gradient of Section~\ref{sec:calibration-results} is not built.

\paragraph{The gates and the families.}
% src: $RUN6/backtest-upworthy-20261103-4bfb354b87e3.json tables.aggregate family_size 52, tables.score_by_year family_size 156, tables.decision_value family_size 6
The gate sizes assume the calibration the backtest is meant to judge; we read the gates as a sizing convention and rest every powered claim on its interval. The aggregate and year tables read 52 and 156 intervals and the decision table six per budget, with adjusted bounds beside them; the remaining tables carry none yet; the pooled interval is the primary claim.

\paragraph{The guard at the size of its cell.}
% src: $RUN6/paper-extras-synthetic-20261103-bf4fd286e2b2.json guard_period_control.rows (planted lift zero, 12 seeds, 36 cells per size, multiplier 2.0): n_tests 3700 old_rule share_admitted 0.556 mean_rows_in_period_level 61.5, old_rule_no_period share_admitted 0.028, new_rule share_admitted 0.056 mean_rows_in_period_level 798.7; n_tests 11100 old_rule 0.333, 183.3, old_rule_no_period 0.000, new_rule 0.000, 2441.3 (tables/guard-period-control.tex)
The first synthetic sensitivity runs admitted a weight in a majority of cells with no planted text effect at the two largest sizes, and the control arm of Table~\ref{tab:guard-period-control} found the cause in P1: it added half of the latest period's shrunken deviation to every cell's prior mean, and a synthetic ledger ends a week into a quarter, so that period held 62 rows of noise at 3,700 tests and 183 at 11,100, which the precedent rung undid; the guard admitted a real gain over a wrong P1 in 56 and 33 percent of cells. No period at all takes the share to 3 and 0 percent, and the rule now in force, the level of the last complete quarter (Section~\ref{sec:method}), takes it to 6 and 0 percent with 799 and 2,441 rows in the level; at the quarterly origins of the backtest the two rules coincide, so no backtest number moved.
% src: $RUN6/paper-extras-synthetic-20261103-bf4fd286e2b2.json guard_gain_diagnostic.rows ratio_of_means 0.68, 0.63, 0.71, 0.67 at n_tests 300, 1000, 3700, 11100 (tables/guard-gain-diagnostic.tex); guard_held_out_gain.rows (held_out_share_of_tests 0.15, 12 seeds, 36 cells per size): mean_bias 0.00928, 0.00203, 0.00056, 0.00000; mean_candidate_held_out_se 0.0031, 0.0009, 0.0003, 0.0001 (tables/guard-gain-held-out.tex); guard_fit_fraction_stability.rows (n_tests 3700): share_admitted 0.056, 0.056, 0.028, 0.056 on the first 85, 90, 95 and 100 percent of rows and 0.056 on a random 85 percent (tables/guard-fit-fraction.tex); $RUN6/paper-extras-upworthy-20261103-4bfb354b87e3.json guard_pairs.summary pooled_realized_p2_minus_p1 0.0010 (-0.0007, 0.0025); positive_control_sensitivity.rows n_tests 11100 text_effect_strength 1.0 recovered_fraction 0.180 (a fifth); $RUN6/paper-extras-upworthy-20261103-4bfb354b87e3.json guard_pairs_rows.pooled location_part 0.0036 (0.0025, 0.0046), scale_part -0.0026 (-0.0040, -0.0010); guard_pairs_rows.rows location_part positive in 7 of 8 pairs (docs/method.md call 233)
Under that rule the mean in-selection gain stays below its standard error at every size (ratios of 0.63 to 0.71, Table~\ref{tab:guard-gain-diagnostic}); what remains is a selection bias that shrinks with the cell: with a random 15 percent of tests held out, the in-selection gain of the best candidate exceeds its held-out gain by 0.0093, 0.0020, 0.0006 and 0.0000 nats at 300, 1,000, 3,700 and 11,100 tests (Table~\ref{tab:guard-gain-held-out}), and the decision at 3,700 tests does not move with the fit set (6, 6, 3 and 6 percent of cells on the first 85 to 100 percent of rows, 6 on a random 85 percent; Table~\ref{tab:guard-fit-fraction}). The guard reads as designed on these ledgers; what an admission on the real archive is worth is the out-of-time comparison of Table~\ref{tab:guard-pairs} (pooled plus 0.0010 nats, within its noise), and the positive control bounds the other side, a rung that recovers a fifth of a planted one-percent lift at the size of the real cell (Section~\ref{sec:synthetic-results}), so the control does not separate an archive with no precedent signal from one whose signal the rung mostly discards. Future work, named and not done: a held-out choice of P2's scale term, the candidate the decomposition of the admitted pairs points at (the location part is plus 0.0036 nats per row, 0.0025 to 0.0046, and positive in seven of the eight pairs; the scale part minus 0.0026, minus 0.0040 to minus 0.0010); a held-out split for the gain in small cells; and an empirical-Bayes variance for P1's period level.

\paragraph{The precedent rung at scale.}
% src: $RUN6/paper-extras-synthetic-20261103-bf4fd286e2b2.json top_m_tail_share.by_size share_p2_weight mean: contrasts 10000 0.5649, 100000 0.5185, 300000 0.6535
On a large ledger the instability check is a lower bound: a query lists its 1,000 highest-weight precedents, and the kernel weight beyond the list is 0.52 to 0.65 of the total at ten thousand to three hundred thousand synthetic contrasts (Appendix~\ref{app:walltime}).

\paragraph{Ethics, data norms and misuse.} Both archives are aggregate, published under CC BY 4.0 and credited by their published names; nothing in the system ingests user-level data, and only the exploratory split of the headline archive is used (Section~\ref{sec:setting}). The system never decides: it prioritizes and sizes, and the live test decides. Three misuses are foreseeable: a team that prioritizes by $\pwin$ writes a ledger selected on the forecaster's own output, so a ledger row should record whether a card influenced the test, with skill reported apart for those rows; P3 must be entered before the card is seen; and an abstained card's base rate is labelled as the quantity it declined to improve on.

% $RUN6 = paper/data/run8-randomization-window
\section{Reproducibility}
\label{sec:repro}
% src: data/CHECKSUMS.json; $RUN6/backtest-upworthy-20261106-2dc461f52ea4.json header.seed 0, header.bootstrap_resamples 500; $RUN6/README.md (backtest seed 0, power seed 20261102); paper/figures/figures.json (every plotted value with its table path, row and field); $RUN6/leakage-audit-upworthy.json, $RUN6/leakage-audit-asos.json, $RUN6/leakage-audit-synthetic.json rungs[].fit_and_scored 0, fit_after_origin 0, outside_window 0, usable_after_mismatch 0, s_star_mismatch 0 for every rung
\iftmlr The code in the anonymized supplementary material\else The released code\fi{} (Python package \texttt{mootcourt}, MIT license) regenerates every results table and every figure from the report files with one command each (\texttt{make tables}, \texttt{make paper-figures}), and the reports themselves with one command per report (Table~\ref{tab:commands}, Appendix~\ref{app:tables}). It lists every plotted value in \texttt{paper/figures/figures.json}. Neither archive is bundled: \texttt{make data} fetches the archive files and refuses any Upworthy file named holdout, confirmatory or undeployed. In a default clone it verifies every file against the committed checksums (Table~\ref{tab:checksums}). With a separate storage root it records the digests on the first download, and the reader compares them with that table. The repository does carry a small sample ledger for tests. Every main backtest and ablation is at seed 0, and Appendix~\ref{app:walltime} lists the identifiers and the other runs' seeds. The leakage audit recomputes that no scored row was visible to the model that scored it and no fit row was usable after its origin: zero violations on all three archives. The audit checks visibility against each test's assigned end, which on the headline archive is the assumed seven days of Section~\ref{sec:setting}, so its zeros hold under that assumption.
% src: .github/workflows/ci.yml (test job on ubuntu-latest, Python 3.11, 3.12 and 3.13; steps after the suite on each version) and scripts/ci_score_report.py: artifact score-report-<python> with the P0, P1 and P2 log score and Brier skill of a fixed-seed backtest of the bundled sample ledger (docs/method.md call 239); tests/test_quick_check.py against tests/fixtures/quick-check/backtest-sample.json (built on macOS; relative 1e-9, 2.0 points on the percentile fields; docs/method.md call 255); docs/troubleshooting.md, Ids differ between machines
On one runner image (ubuntu), continuous integration runs the backtest on the bundled sample ledger under three Python versions and uploads an artifact holding the pooled scores of P0, P1 and P2, so the job itself shows whether the Python versions agree. Within the same job, the quick check gives the cross-platform reading: it rebuilds the sample backtest on the runner and sets it against a fixture made on macOS, at a relative $10^{-9}$. It tolerates 2 points on the percentile fields, so any move smaller than that there passes unseen. Run on a different operating system, the same fit may come out with different last bits in its parameters, and that shifts the model id and the feature hash even though the scores match. So the comparison of cards and models uses the identifiers stored for them, instead of recomputing them (Appendix~\ref{app:walltime}).
% src: docs/method.md call 255 (under the strict comparison, ece_percentile and ece_percentile_independent moved by up to 1.1 points between macOS and ubuntu on the same seed and by 5.2 points in the worst case on the sample report; 26 tied draws of 0.2 points each, a bound and not an observed move; the mid-rank rule counts a draw within 1e-9 * max(1, |observed|) of the observed value as half, src/mootcourt/harness_stats.py null_percentile); scripts/quick_check.py PERCENTILE_TOLERANCE 2.0
We define the calibration percentile as the rank the observed calibration error takes among the draws simulated under its null. Any draw landing within $10^{-9}$ of the observed value is counted as one half. Under the strict comparison this rule replaced, draws like that could land on either side, which shifted the percentile by as much as 1.1 points from macOS to ubuntu with the seed unchanged. For the sample ledger, 5.2 points is the largest shift it could have given a single percentile.

%% tmlr:move-begin artifact
%% tmlr:note artifact How the protocol artifact is kept and how a result is added to it are in Appendix~\ref{sec:artifact}.
%% tmlr:title artifact The protocol artifact and its upkeep
\label{sec:artifact}
\paragraph{Protocol artifact.} Adding a result means a contributor reruns that same backtest against the same snapshot and puts out its report card. Two things go into the artifact, and one stays out.
\begin{itemize}
\item Inside it sit the ledger schema (Table~\ref{tab:schema}), the snapshot digests and admission rules of Section~\ref{sec:setting}, the quarterly origins and the names of the refused files.
% src: README.md, License and citations (MIT)
\item It holds the harness and the report card under the MIT license.
% src: README.md, License and citations (Upworthy Research Archive CC BY 4.0; ASOS Digital Experiments Dataset CC BY 4.0)
\item The raw archive files are absent from it; they are fetched under CC BY 4.0 and never redistributed.
\end{itemize}
The protocol is kept as follows.
\begin{itemize}
\item \iftmlr The authors will maintain it in a public repository after review.\else The author maintains it through the repository, by issues and pull requests.\fi
\item To submit a report card from a new archive, a pull request adds a run folder under \texttt{paper/data} together with its README, in the way the contributing page sets out.
% src: CHANGELOG.md (semantic versioning; before 1.0.0 a schema change is a minor version)
\item \iftmlr The code is under semantic versioning.\else The repository is under semantic versioning. Version 1.0.0 names the released snapshot digests, and before it a schema change counted as a minor version.\fi
\item The snapshot digest names the data a result was read on.
% src: Zenodo record https://zenodo.org/api/records/23988587 (doi 10.5281/zenodo.23988587, conceptdoi 10.5281/zenodo.23988586, read 5 November 2026); $RUN6/README.md (commit 5e4fedf, the code state of this paper's run; commit a7e9f44 for the two learning-curve reports; not repeated: scale/, run-log.json and run-log.md, the timed run, which stay run 7's); $RUN6/README.md (the fourth exception paragraph: five reports at commit 3971829; the fifth: eight files at the analysis code of commit 73d0d70)
\item \iftmlr An earlier version of the code predates\else Release 1.0.0 is archived at Zenodo under DOI\newline \mbox{10.5281/zenodo.23988587}. It predates\fi{} the date rule set out in Section~\ref{sec:setting}, under which tests that began inside the archive's flagged randomization window are excluded. This paper's results rest on \iftmlr one code state\else commit 5e4fedf\fi, apart from six exceptions: one comparison added at \iftmlr a later code state\else commit a7e9f44\fi, the ablation with a Student-t base rate, run at \iftmlr a later code state still\else commit f4566ca\fi, its repeat with exposure fixed, run at \iftmlr a code state that changes no file of the package\else commit e05247e, which changes no file under \texttt{src}\fi, the base rate with a quarter term and the reading by origin, run at \iftmlr a code state after that\else commit 3971829\fi, the runs with exposure fixed on visible rows, the interval with quarters as units and the no-effect reference, run at \iftmlr the latest code state\else commit 73d0d70\fi, and the wall times and scale measurements of Appendix~\ref{app:walltime}, taken from earlier code states that the appendix describes. The run folder \texttt{paper/data/run8-randomization-window} holds the reports.
% src: Zenodo record https://zenodo.org/api/records/23988587 (conceptdoi 10.5281/zenodo.23988586, read 5 November 2026)
\item \iftmlr The code, the reports and the plan are in the anonymized supplementary material.\else The concept DOI,\newline \mbox{10.5281/zenodo.23988586}, resolves to the latest release.\fi
% src: references.bib mootcourt2026plan (posted 2026-11-04)
\item \iftmlr The plan is the copy in that material.\else The plan is the copy posted publicly on 4 November 2026 \cite{mootcourt2026plan}.\fi
\end{itemize}
%% tmlr:move-end artifact

\section{Conclusion}
\label{sec:conclusion}
A team's concluded experiments predict how the next result will be distributed better than a base rate does, and on this evidence their skill on which test will be significant is at most small and depends on how that event is defined. A text classifier on the headlines ranks tests a little better than a random order matched on arm count. Precedents rank them a little better than the prior, an edge that may be traffic and not precedents. At realized precision neither interval on significance clears zero, and the positive control shows this backtest is too weak to rule out a precedent signal. Forecasts as issued are miscalibrated, mostly in a level they do not anticipate, and sizing a test from the forecast shows no advantage over sizing it from a base rate.

These claims have limits. Every number is exploratory: one split of one publisher's archive, read repeatedly, with the method changed between runs after results were seen. The archive of product experiments is too small to decide anything, so nothing here is shown for product experiments, and the system has never run in production. The confirmatory split is unopened and is the next test. What a team can take from this paper is the protocol: run the same backtest on its own ledger, and let that report, with its gates and its abstentions, say what its forecasts are worth.


\bibliographystyle{ACM-Reference-Format}
\bibliography{references}

% $REPORTS = /Volumes/Crucial X10/mootcourt-data/reports
\appendix
% Float limits for the appendix only: it is tables without text, so several floats per page must be allowed; the two-column
% floats come first so that they stack at page tops, and the final \clearpage flushes every pending float before the class
% balances the last columns (without it the class drops pending single-column floats).
\setcounter{topnumber}{8}\setcounter{bottomnumber}{8}\setcounter{dbltopnumber}{6}\setcounter{totalnumber}{16}
\renewcommand{\topfraction}{1}\renewcommand{\bottomfraction}{1}\renewcommand{\dbltopfraction}{1}\renewcommand{\textfraction}{0}\renewcommand{\floatpagefraction}{0.5}
\section{Additional Tables}
\label{app:tables}
Tables in this appendix are generated by the same script as those in the body and carry the same provenance comment. The two-column tables come first: the commands behind every table, the pooled deployment-shaped results, the learning curve and the weekly replay of the headline archive, then the checksums, the synthetic archive, and the deployment-shaped learning curve and replay; the single-column tables follow in the order of the sections that cite them.

\begin{table*}[t]\centering\small
\caption{One command per table. Commands run from the repository root with the data root set; \texttt{S} is \texttt{upworthy}, \texttt{asos} or \texttt{synthetic}, and the \texttt{make} targets wrap the \texttt{moot} commands shown in the repository's README.}
\label{tab:commands}
\begin{tabular}{ll}
\toprule
what & command \\
\midrule
data, checksums, ledger & \texttt{make data} \\
gates (Table~\ref{tab:gates}) & \texttt{moot simulate power} \\
main, calibration, decision, abstention, selection tables & \texttt{moot backtest -{}-source S -{}-origins quarterly -{}-seed 0} \\
challenger rows of Tables~\ref{tab:main-upworthy} and~\ref{tab:main-asos} & \texttt{make challengers SOURCE=S} \\
ablations (Tables~\ref{tab:ablations-upworthy} and~\ref{tab:ablations-asos}) & \texttt{make ablations SOURCE=S} \\
cautionary run (Table~\ref{tab:cautionary}) & \texttt{moot backtest -{}-source upworthy -{}-effect-scale relative} \\
 & \texttt{\ \ -{}-s-star realized} \\
guard comparison & \texttt{python scripts/compare\_guard.py} \\
learning curve (Tables~\ref{tab:learning-upworthy} and~\ref{tab:learning-asos}) & \texttt{make learning-curve SOURCE=S} \\
trigger ablation (Tables~\ref{tab:trigger-ablation} and~\ref{tab:trigger-ablation-asos}) & \texttt{moot backtest -{}-source S -{}-origins quarterly -{}-drop-trigger T -{}-seed 0} \\
skill by year (Tables~\ref{tab:drift} and~\ref{tab:drift-asos}) & the main backtest report, \texttt{tables.score\_by\_year} \\
weekly replay (Tables~\ref{tab:replay} and~\ref{tab:replay-asos}) & \texttt{make replay SOURCE=S} \\
leakage audit & \texttt{make audit SOURCE=S} \\
card latency & \texttt{python scripts/card\_latency.py -{}-source upworthy} \\
timed run (Table~\ref{tab:walltime}) & \texttt{scripts/timed\_run.sh} \\
scale runs & \texttt{moot benchmark -{}-scale 100000}, \texttt{-{}-scale 1000000} \\
LaTeX tables & \texttt{python paper/build\_tables.py -{}-reports \$ROOT/reports} \\
 & \texttt{\ \ -{}-out paper/tables} \\
\bottomrule
\end{tabular}
\end{table*}

\begin{table*}[t]\centering\scriptsize
\caption{Deployment-shaped archive, pooled over origins and metric cells. Logarithmic scores are positive because the standardized effects and their standard errors are small numbers.}
\label{tab:main-asos}
\resizebox{\textwidth}{!}{% generated by paper/build_tables.py from /Volumes/Crucial X10/mootcourt-data/reports/backtest-asos-20261103-6bb5a15c5ff0.json (source asos, snapshot 20261103-6bb5a15c5ff0, model 548c646b9cb1, 337 scored rows over 7 origins; gate decision target 0.9)
\begin{tabular}{llrrrrrrrr}
\toprule
kind & model & tests & log score & CRPS & Brier skill (90\%) & dir.\ acc.\ / ceiling & cov.\ 80 / 90 & abst. & decid. \\
\midrule
ladder & P0 & 68 & 4.383 & 0.00 & -0.443 (-0.836, -0.185) & 0.499 / 0.806 & 0.91 / 0.93 & 0.448 & no \\
ladder & P1 & 68 & 4.610 & 0.00 & -0.107 (-0.237, -0.002) & 0.522 / 0.806 & 0.78 / 0.86 & 0.448 & no \\
ladder & P2 & 68 & 4.610 & 0.00 & -0.107 (-0.237, -0.002) & 0.522 / 0.806 & 0.78 / 0.86 & 0.448 & no \\
baseline & always\_reference & 68 &  &  &  & 0.496 / 0.806 &  & 0.448 & no \\
baseline & always\_variant & 68 &  &  &  & 0.504 / 0.806 &  & 0.448 & no \\
challenger & GBM & 68 & 3.629 & 0.00 & -0.714 (-1.259, -0.377) & 0.493 / 0.806 & 0.99 / 1.00 & 0.448 & no \\
challenger & local\_eb & 68 & 4.428 & 0.00 & -0.383 (-0.736, -0.138) & 0.487 / 0.806 & 0.89 / 0.90 & 0.448 & no \\
\bottomrule
\end{tabular}
}
\end{table*}

\begin{table*}[t]\centering\scriptsize
\caption{Brier skill against the base rate by fraction of the visible ledger at the final origin (headline archive); 90 percent bootstrap intervals by held-out test of the seed-averaged skill. The direction baselines carry no probability and have no skill.}
\label{tab:learning-upworthy}
\resizebox{\textwidth}{!}{% generated by paper/build_tables.py from learning-curve-upworthy-20261106-2dc461f52ea4.json (held-out rows 1143, seeds 20, gates {'n_200': 175, 'n_400': 250, 'by_target': {'0.80': {'n_200': 125, 'n_400': 200}, '0.90': {'n_200': 175, 'n_400': 250}}, 'decision_target': 0.9})
\begin{tabular}{rrrrrrr}
\toprule
fraction of ledger & GBM skill & P0 skill & P1 skill & P2 skill & local\_eb skill & text\_lr skill \\
\midrule
0.01 & -0.058 (-0.094, -0.033) & -0.042 (-0.053, -0.034) & 0.009 (-0.004, 0.027) & 0.009 (-0.005, 0.027) & -0.011 (-0.020, 0.000) & -0.025 (-0.028, -0.021) \\
0.02 & -0.094 (-0.145, -0.058) & -0.027 (-0.034, -0.022) & 0.012 (-0.002, 0.030) & 0.011 (-0.002, 0.028) & -0.014 (-0.024, -0.003) & -0.014 (-0.019, -0.008) \\
0.05 & -0.088 (-0.143, -0.048) & -0.019 (-0.027, -0.014) & 0.018 (0.004, 0.037) & 0.018 (0.004, 0.037) & -0.005 (-0.019, 0.008) & -0.008 (-0.013, 0.000) \\
0.10 & -0.077 (-0.128, -0.039) & -0.008 (-0.013, -0.005) & 0.020 (0.008, 0.037) & 0.020 (0.007, 0.038) & 0.003 (-0.007, 0.017) & -0.006 (-0.015, 0.003) \\
0.25 & -0.054 (-0.096, -0.021) & -0.003 (-0.005, -0.002) & 0.024 (0.011, 0.043) & 0.025 (0.010, 0.046) & 0.012 (-0.003, 0.030) & 0.003 (-0.008, 0.014) \\
0.50 & -0.039 (-0.079, -0.010) & 0.000 (-0.002, 0.002) & 0.024 (0.011, 0.042) & 0.027 (0.012, 0.049) & 0.019 (0.004, 0.039) & 0.006 (-0.010, 0.021) \\
1.00 & -0.041 (-0.079, -0.015) & 0.000 (0.000, 0.000) & 0.024 (0.011, 0.042) & 0.029 (0.014, 0.051) & 0.013 (-0.011, 0.037) & 0.017 (-0.007, 0.037) \\
\bottomrule
\end{tabular}
}
\end{table*}

\IfFileExists{tables/replay-upworthy.tex}{%
\begin{table*}[t]\centering\scriptsize
\caption{Weekly replay of the headline archive, summed by calendar quarter of the join week: cards issued, outcomes joined, abstention share, typed outcomes (hit, magnitude miss, direction miss, confident but wrong, abstained but clear), rolling 90 percent coverage at the quarter's last week, and weeks in which a cell was marked drifted.}
\label{tab:replay}
\resizebox{\textwidth}{!}{% generated by paper/build_tables.py from replay-upworthy-20261106-2dc461f52ea4.json
\begin{tabular}{lrrrrrrrrrrr}
\toprule
quarter & weeks & cards & joined & abst. & hit & magn. & dir. & conf.\ wrong & abst.\ clear & cov.\ 90 (last wk) & flagged wks \\
\midrule
2013Q2 & 13 & 679 & 591 & 0.52 & 420 & 19 & 77 & 1 & 74 & 0.85 & 0 \\
2013Q3 & 14 & 0 & 88 &  & 53 & 6 & 22 & 0 & 7 & 0.62 & 0 \\
2013Q4 & 13 & 0 & 0 &  & 0 & 0 & 0 & 0 & 0 &  & 0 \\
2014Q1 & 13 & 1322 & 1066 & 0.15 & 664 & 62 & 263 & 27 & 50 & 0.79 & 1 \\
2014Q2 & 13 & 1943 & 1888 & 0.46 & 1372 & 56 & 290 & 1 & 169 & 0.92 & 1 \\
2014Q3 & 13 & 2837 & 2635 & 0.03 & 2057 & 94 & 417 & 46 & 21 & 0.95 & 0 \\
2014Q4 & 13 & 2510 & 2872 & 0.16 & 2135 & 128 & 317 & 182 & 110 & 0.88 & 0 \\
2015Q1 & 13 & 1173 & 1195 & 0.05 & 867 & 86 & 154 & 71 & 17 & 0.90 & 0 \\
2015Q2 & 6 & 233 & 362 & 0.08 & 277 & 22 & 40 & 14 & 9 & 0.90 & 4 \\
\midrule
all & 111 & 10697 & 10697 & 0.19 & 7845 & 473 & 1580 & 342 & 457 & 0.90 & 6 \\
\bottomrule
\end{tabular}
}
\end{table*}
}{}

\begin{table*}[t]\centering\scriptsize
\caption{Recorded checksums of the downloaded files (\texttt{data/CHECKSUMS.json}).}
\label{tab:checksums}
\begin{tabular}{lrl}
\toprule
file under \texttt{data/raw/} & bytes & sha256 \\
\midrule
\texttt{upworthy/upworthy-archive-exploratory-} & 14,260,949 & \texttt{8368313b060f4015a0c6fb34e6d788163cee29554144aba9390b14922eb9d8ce} \\
\texttt{\ \ packages-03.12.2020.csv} & & \\
\texttt{asos/asos\_digital\_experiments\_dataset.parquet} & 1,003,239 & \texttt{bdf88b27185d3f7e65912cbb421129a32ae347c2fb0d9442fe54d74266551524} \\
\texttt{asos/Datasheet - ASOS Digital Experiments Dataset.pdf} & 206,619 & \texttt{e1e42e29a30fa7ad9eeeb2ac7b61d2f0221a531d2b030a8abd9c9f6371a93563} \\
\bottomrule
\end{tabular}
\end{table*}

\begin{table*}[t]\centering\scriptsize
\caption{Synthetic archive with known truth, pooled over origins and cells.}
\label{tab:main-synthetic}
\resizebox{\textwidth}{!}{% generated by paper/build_tables.py from backtest-synthetic-20261103-bf4fd286e2b2.json (source synthetic, snapshot 20261103-bf4fd286e2b2, model 6377324df73a, 1814 scored rows over 8 origins; gate decision target 0.9)
\begin{tabular}{llrrrrrrrrrr}
\toprule
kind & model & tests & log score & minus P1 (90\%) & CRPS & Brier / base / floor & Brier skill (90\%) & dir.\ acc.\ / ceiling & cov.\ 80 / 90 & abst. & decid. \\
\midrule
challenger & GBM & 907 & -2.457 & -0.0623 (-0.0786, -0.0476) & 1.58 & 0.1470 / 0.1391 / 0.0683 & -0.0573 (-0.0860, -0.0309) & 0.639 / 0.809 & 0.81 / 0.90 & 0.182 & yes \\
challenger & GBM+cal & 907 &  &  &  & 0.1485 / 0.1391 / 0.0683 & -0.0678 (-0.0986, -0.0391) &  / 0.809 &  & 0.182 & yes \\
ladder & P0 & 907 & -2.405 & -0.0099 (-0.0183, -0.0009) & 1.53 & 0.1391 / 0.1391 / 0.0683 & 0.0000 (0.0000, 0.0000) & 0.648 / 0.809 & 0.78 / 0.88 & 0.182 & yes \\
ladder & P1 & 907 & -2.395 & 0.0000 (0.0000, 0.0000) & 1.53 & 0.1388 / 0.1391 / 0.0683 & 0.0023 (-0.0128, 0.0168) & 0.648 / 0.809 & 0.78 / 0.89 & 0.182 & yes \\
ladder & P2 & 907 & -2.395 & 0.0000 (0.0000, 0.0000) & 1.53 & 0.1388 / 0.1391 / 0.0683 & 0.0023 (-0.0128, 0.0168) & 0.648 / 0.809 & 0.78 / 0.89 & 0.182 & yes \\
baseline & always\_reference & 907 &  &  &  &  &  & 0.460 / 0.809 &  & 0.182 & yes \\
baseline & always\_variant & 907 &  &  &  &  &  & 0.540 / 0.809 &  & 0.182 & yes \\
challenger & local\_eb & 907 & -2.416 & -0.0216 (-0.0292, -0.0147) & 1.54 & 0.1404 / 0.1391 / 0.0683 & -0.0097 (-0.0274, 0.0083) & 0.641 / 0.809 & 0.79 / 0.89 & 0.182 & yes \\
challenger & local\_eb+cal & 907 &  &  &  & 0.1407 / 0.1391 / 0.0683 & -0.0115 (-0.0308, 0.0061) &  / 0.809 &  & 0.182 & yes \\
challenger & text\_lr & 907 &  &  &  & 0.1441 / 0.1391 / 0.0683 & -0.0359 (-0.0550, -0.0165) & 0.607 / 0.809 &  & 0.182 & yes \\
challenger & text\_lr+cal & 907 &  &  &  & 0.1430 / 0.1391 / 0.0683 & -0.0281 (-0.0463, -0.0119) &  / 0.809 &  & 0.182 & yes \\
\bottomrule
\end{tabular}
}
\end{table*}

\begin{table*}[t]\centering\scriptsize
\caption{Brier skill against the base rate by fraction of the visible ledger at the final origin (deployment-shaped archive; 56 held-out rows in 9 tests; the five-percent subsample could not be fit).}
\label{tab:learning-asos}
\resizebox{\textwidth}{!}{% generated by paper/build_tables.py from /Volumes/Crucial X10/mootcourt-data/reports/learning-curve-asos-20261103-6bb5a15c5ff0.json (held-out rows 56, seeds 20, gates {'n_200': 175, 'n_400': 250, 'by_target': {'0.80': {'n_200': 125, 'n_400': 200}, '0.90': {'n_200': 175, 'n_400': 250}}, 'decision_target': 0.9})
\begin{tabular}{rrrrrr}
\toprule
fraction of ledger & GBM skill & P0 skill & P1 skill & P2 skill & local\_eb skill \\
\midrule
0.05 & -0.202 (-1.454, 0.170) & 0.000 (0.000, 0.000) & -0.146 (-0.273, 0.045) & -0.146 (-0.273, 0.045) & -0.235 (-0.803, 0.019) \\
0.10 & -0.268 (-2.005, 0.192) & 0.000 (0.000, 0.000) & -0.121 (-0.281, -0.028) & -0.121 (-0.281, -0.028) & -0.198 (-0.673, -0.010) \\
0.25 & -0.295 (-2.155, 0.191) & 0.000 (0.000, 0.000) & -0.117 (-0.214, -0.014) & -0.117 (-0.212, -0.019) & -0.212 (-0.752, -0.011) \\
0.50 & -0.345 (-2.153, 0.153) & 0.000 (0.000, 0.000) & -0.142 (-0.271, -0.026) & -0.142 (-0.271, -0.026) & -0.249 (-0.911, 0.005) \\
1.00 & -0.358 (-2.233, 0.155) & 0.000 (0.000, 0.000) & -0.182 (-0.500, -0.037) & -0.182 (-0.500, -0.037) & -0.121 (-0.659, 0.065) \\
\bottomrule
\end{tabular}
}
\end{table*}

\IfFileExists{tables/replay-asos.tex}{%
\begin{table*}[t]\centering\scriptsize
\caption{Weekly replay (deployment-shaped archive, nominal calendar); columns as in Table~\ref{tab:replay}.}
\label{tab:replay-asos}
\resizebox{\textwidth}{!}{% generated by paper/build_tables.py from /Volumes/Crucial X10/mootcourt-data/reports/replay-asos-20261103-6bb5a15c5ff0.json
\begin{tabular}{lrrrrrrrrrrr}
\toprule
quarter & weeks & cards & joined & abst. & hit & magn. & dir. & conf.\ wrong & abst.\ clear & cov.\ 90 (last wk) & flagged wks \\
\midrule
2019Q2 & 13 & 45 & 20 & 1.00 & 13 & 0 & 0 & 0 & 7 &  & 0 \\
2019Q3 & 14 & 48 & 49 & 1.00 & 34 & 0 & 0 & 0 & 15 &  & 0 \\
2019Q4 & 13 & 48 & 40 & 0.75 & 38 & 0 & 0 & 0 & 2 & 1.00 & 0 \\
2020Q1 & 13 & 48 & 60 & 0.00 & 32 & 10 & 4 & 3 & 11 & 0.53 & 0 \\
2020Q2 & 13 & 36 & 28 & 0.00 & 9 & 5 & 9 & 5 & 0 & 0.42 & 0 \\
2020Q3 & 13 & 56 & 60 & 0.04 & 28 & 9 & 14 & 9 & 0 & 0.57 & 1 \\
2020Q4 & 13 & 56 & 40 & 0.29 & 20 & 7 & 7 & 1 & 5 & 0.62 & 13 \\
2021Q1 & 9 & 0 & 40 &  & 35 & 0 & 3 & 0 & 2 & 1.00 & 8 \\
\midrule
all & 101 & 337 & 337 & 0.44 & 209 & 31 & 37 & 18 & 42 & 0.66 & 22 \\
\bottomrule
\end{tabular}
}
\end{table*}
}{}

\begin{table}[t]
\centering\small
\caption{Ledger row.}
\label{tab:schema}
\resizebox{\columnwidth}{!}{\begin{tabular}{ll}
\toprule
field & meaning \\
\midrule
row\_id, test\_id, variant\_id & identifiers \\
cell, period & kind of change; quarter of the final read \\
started\_at, final\_read\_at, usable\_after & timeline \\
$y$, $s$, $n_{\mathrm{ref}}$, $n_{\mathrm{var}}$ & effect, standard error, exposures \\
status & admitted or an exclusion reason \\
features & pre-launch fields; post-launch block kept apart \\
\bottomrule
\end{tabular}}
\end{table}

\begin{table}[t]
\centering\small
\caption{Capabilities against tools (Y yes, P partial, N no, ND not documented publicly). Columns: pre-launch outcome distribution; calibrated win probability; abstention; precedent retrieval; scoring against live results; learned prior from the organization's own archive; experiment memory; open source.}
\label{tab:landscape}
\resizebox{\columnwidth}{!}{\begin{tabular}{lcccccccc}
\toprule
 & dist. & $\pwin$ & abst. & prec. & scored & prior & memory & OSS \\
\midrule
Commercial platforms (13) & N & N & N & P & N & P & P & N \\
GrowthBook & N & N & N & P & N & P & Y & Y \\
Other open-source libraries & N & N & N & N & N & N & N & Y \\
Internal platforms (public papers) & N & N & N & P & P & P & Y & N \\
Simulation and LLM agents & P & N & P & N & P & N & N & P \\
Archive priors (research) & N & N & N & N & N & Y & N & N \\
This work & Y & Y & Y & Y & Y & Y & Y & Y \\
\bottomrule
\end{tabular}}
\end{table}

\begin{table}[t]\centering\scriptsize
\caption{The cautionary run on the headline archive (relative lift, realized precision), by decile of realized standard error. Effects in percent relative lift.}
\label{tab:cautionary}
\resizebox{\columnwidth}{!}{% generated by paper/build_tables.py from paper-extras-upworthy-20261106-2dc461f52ea4.json
\begin{tabular}{rrrrrr}
\toprule
decile of realized $s$ & median $s$ & mean $y$ & event rate & P1 mean $\psig$ & P1 cov.\ 90 \\
\midrule
1 & 9.8 & -42.5 & 0.001 & 0.141 & 0.62 \\
2 & 12.7 & -29.2 & 0.001 & 0.104 & 0.89 \\
3 & 14.9 & -23.7 & 0.009 & 0.085 & 0.97 \\
4 & 17.0 & -17.2 & 0.009 & 0.071 & 0.99 \\
5 & 19.3 & -10.2 & 0.017 & 0.061 & 0.99 \\
6 & 21.9 & -4.8 & 0.022 & 0.052 & 0.98 \\
7 & 25.1 & 3.3 & 0.032 & 0.045 & 0.96 \\
8 & 29.1 & 14.7 & 0.065 & 0.038 & 0.92 \\
9 & 36.2 & 30.7 & 0.091 & 0.032 & 0.86 \\
10 & 52.0 & 84.4 & 0.200 & 0.026 & 0.66 \\
\bottomrule
\end{tabular}
}
\end{table}

\begin{table}[t]\centering\scriptsize
\caption{Calibration by decile of planned exposure per arm (headline archive).}
\label{tab:calibration-upworthy}
\resizebox{\columnwidth}{!}{% generated by paper/build_tables.py from backtest-upworthy-20261106-2dc461f52ea4.json
\begin{tabular}{rrrrrrr}
\toprule
decile & \multicolumn{3}{c}{P1: mean $\psig$ / event rate / cov.\ 90} & \multicolumn{3}{c}{P2: mean $\psig$ / event rate / cov.\ 90} \\
\midrule
1 & 0.074 & 0.064 & 0.86 & 0.072 & 0.064 & 0.86 \\
2 & 0.093 & 0.077 & 0.87 & 0.091 & 0.077 & 0.87 \\
3 & 0.092 & 0.074 & 0.87 & 0.088 & 0.074 & 0.87 \\
4 & 0.085 & 0.105 & 0.88 & 0.084 & 0.105 & 0.88 \\
5 & 0.085 & 0.086 & 0.89 & 0.084 & 0.086 & 0.89 \\
6 & 0.100 & 0.085 & 0.89 & 0.098 & 0.085 & 0.89 \\
7 & 0.106 & 0.101 & 0.91 & 0.104 & 0.101 & 0.90 \\
8 & 0.106 & 0.089 & 0.89 & 0.105 & 0.089 & 0.89 \\
9 & 0.119 & 0.084 & 0.94 & 0.119 & 0.084 & 0.94 \\
10 & 0.133 & 0.089 & 0.95 & 0.132 & 0.089 & 0.95 \\
\bottomrule
\end{tabular}
}
\end{table}

\begin{table}[t]\centering\small
\caption{Lift captured by prioritization policies at $K$ equal to a quarter of each period's tests (headline archive); regret is the oracle's capture minus the policy's, in percentage points of click-through rate summed over the period's tests.}
\label{tab:decision-upworthy}
\resizebox{\columnwidth}{!}{% generated by paper/build_tables.py from backtest-upworthy-20261106-2dc461f52ea4.json (tables.decision_value; each mean is over the 7 per-origin rows of the policy and excludes the pooled origin=ALL row)
\begin{tabular}{llrr}
\toprule
policy & model & captured fraction (mean over origins) & regret (mean) \\
\midrule
random & all & 0.251 & 60.5 \\
p\_win & GBM & 0.278 & 56.8 \\
p\_win & P0 & 0.251 & 60.6 \\
p\_win & P1 & 0.249 & 60.7 \\
p\_win & P2 & 0.289 & 57.0 \\
p\_win & local\_eb & 0.291 & 56.5 \\
p\_win & text\_lr & 0.324 & 52.6 \\
arm\_count & all & 0.311 & 52.3 \\
oracle & all & 0.686 & 0.000 \\
random\_matched & GBM & 0.258 & 59.2 \\
random\_matched & P0 & 0.250 & 60.3 \\
random\_matched & P1 & 0.251 & 60.2 \\
random\_matched & P2 & 0.258 & 59.5 \\
random\_matched & local\_eb & 0.267 & 59.1 \\
random\_matched & text\_lr & 0.265 & 59.0 \\
\bottomrule
\end{tabular}
}
\end{table}

\begin{table}[t]\centering\small
\caption{Abstention triggers (headline archive). A row can fire several triggers.}
\label{tab:abstention-upworthy}
\resizebox{\columnwidth}{!}{% generated by paper/build_tables.py from backtest-upworthy-20261103-4bfb354b87e3.json
\begin{tabular}{lrrrrrr}
\toprule
cell & n & rate & support & novelty & reliability & instability \\
\midrule
ALL & 13247 & 0.141 & 304 & 1121 & 781 & 0 \\
both & 471 & 0.221 & 68 & 36 & 68 & 0 \\
headline\_only & 7417 & 0.115 & 0 & 375 & 477 & 0 \\
image\_only & 5359 & 0.171 & 236 & 710 & 236 & 0 \\
\bottomrule
\end{tabular}
}
\end{table}

\IfFileExists{tables/trigger-ablation-upworthy.tex}{%
\begin{table}[t]\centering\scriptsize
\caption{One abstention trigger switched off at a time (headline archive); the precedent rung, scored on the cards that issue under each setting, with bootstrap-by-test intervals.}
\label{tab:trigger-ablation}
\resizebox{\columnwidth}{!}{% generated by paper/build_tables.py from /Volumes/Crucial X10/mootcourt-data/reports/backtest-upworthy-20261103-4bfb354b87e3.json; /Volumes/Crucial X10/mootcourt-data/reports/ablations/backtest-upworthy-20261103-4bfb354b87e3--drop-trigger-support.json; /Volumes/Crucial X10/mootcourt-data/reports/ablations/backtest-upworthy-20261103-4bfb354b87e3--drop-trigger-novelty.json; /Volumes/Crucial X10/mootcourt-data/reports/ablations/backtest-upworthy-20261103-4bfb354b87e3--drop-trigger-reliability.json; /Volumes/Crucial X10/mootcourt-data/reports/ablations/backtest-upworthy-20261103-4bfb354b87e3--drop-trigger-instability.json
\begin{tabular}{lrrrrr}
\toprule
setting & abst.\ rate & issued rows & log score, issued & Brier, issued & dir.\ acc., issued \\
\midrule
none (main run) & 0.132 & 11504 & -0.855 (-0.878, -0.833) & 0.082 (0.078, 0.086) & 0.582 (0.572, 0.592) \\
drop support & 0.132 & 11504 & -0.855 (-0.878, -0.833) & 0.082 (0.078, 0.086) & 0.582 (0.572, 0.592) \\
drop novelty & 0.059 & 12466 & -0.856 (-0.879, -0.834) & 0.082 (0.078, 0.086) & 0.580 (0.570, 0.590) \\
drop reliability & 0.096 & 11979 & -0.859 (-0.882, -0.837) & 0.083 (0.079, 0.087) & 0.577 (0.567, 0.586) \\
drop instability & 0.132 & 11504 & -0.855 (-0.878, -0.833) & 0.082 (0.078, 0.086) & 0.582 (0.572, 0.592) \\
\bottomrule
\end{tabular}
}
\end{table}
}{}

\IfFileExists{tables/trigger-ablation-asos.tex}{%
\begin{table}[t]\centering\scriptsize
\caption{One abstention trigger switched off at a time (deployment-shaped archive).}
\label{tab:trigger-ablation-asos}
\resizebox{\columnwidth}{!}{% generated by paper/build_tables.py from /Volumes/Crucial X10/mootcourt-data/reports/backtest-asos-20261103-6bb5a15c5ff0.json; /Volumes/Crucial X10/mootcourt-data/reports/ablations/backtest-asos-20261103-6bb5a15c5ff0--drop-trigger-support.json; /Volumes/Crucial X10/mootcourt-data/reports/ablations/backtest-asos-20261103-6bb5a15c5ff0--drop-trigger-novelty.json; /Volumes/Crucial X10/mootcourt-data/reports/ablations/backtest-asos-20261103-6bb5a15c5ff0--drop-trigger-reliability.json; /Volumes/Crucial X10/mootcourt-data/reports/ablations/backtest-asos-20261103-6bb5a15c5ff0--drop-trigger-instability.json
\begin{tabular}{lrrrrr}
\toprule
setting & abst.\ rate & issued rows & log score, issued & Brier, issued & dir.\ acc., issued \\
\midrule
none (main run) & 0.448 & 186 & 4.421 (3.781, 4.969) & 0.164 (0.110, 0.230) & 0.527 (0.438, 0.606) \\
drop support & 0.208 & 267 & 4.673 (4.212, 5.069) & 0.136 (0.097, 0.182) & 0.558 (0.468, 0.625) \\
drop novelty & 0.374 & 211 & 4.463 (3.855, 4.947) & 0.149 (0.097, 0.211) & 0.502 (0.417, 0.575) \\
drop reliability & 0.448 & 186 & 4.421 (3.781, 4.969) & 0.164 (0.110, 0.230) & 0.527 (0.438, 0.606) \\
drop instability & 0.448 & 186 & 4.421 (3.781, 4.969) & 0.164 (0.110, 0.230) & 0.527 (0.438, 0.606) \\
\bottomrule
\end{tabular}
}
\end{table}
}{}

\IfFileExists{tables/drift-upworthy.tex}{%
\begin{table}[t]\centering\scriptsize
\caption{Skill by calendar year of the test's start (headline archive), all cells pooled.}
\label{tab:drift}
\resizebox{\columnwidth}{!}{% generated by paper/build_tables.py from backtest-upworthy-20261106-2dc461f52ea4.json (tables.score_by_year)
\begin{tabular}{rlrrrrrrr}
\toprule
year & model & rows & tests & event rate & log score (90\%) & Brier skill (90\%) & dir.\ acc. & abst. \\
\midrule
2013 & P0 & 679 & 245 & 0.1208 & -1.174 (-1.365, -1.038) & 0.0000 (0.0000, 0.0000) & 0.511 & 1.000 \\
2013 & P1 & 679 & 245 & 0.1208 & -0.989 (-1.092, -0.895) & 0.0334 (0.0121, 0.0521) & 0.451 & 1.000 \\
2013 & P2 & 679 & 245 & 0.1208 & -0.993 (-1.099, -0.897) & 0.0373 (0.0120, 0.0603) & 0.485 & 1.000 \\
2014 & P0 & 8605 & 2303 & 0.0873 & -0.873 (-0.902, -0.848) & 0.0000 (0.0000, 0.0000) & 0.576 & 0.308 \\
2014 & P1 & 8605 & 2303 & 0.0873 & -0.745 (-0.773, -0.718) & -0.0036 (-0.0070, -0.0004) & 0.576 & 0.308 \\
2014 & P2 & 8605 & 2303 & 0.0873 & -0.744 (-0.773, -0.718) & -0.0035 (-0.0070, -0.0004) & 0.576 & 0.308 \\
2015 & P0 & 1413 & 401 & 0.0566 & -0.751 (-0.798, -0.705) & 0.0000 (0.0000, 0.0000) & 0.580 & 0.050 \\
2015 & P1 & 1413 & 401 & 0.0566 & -0.718 (-0.783, -0.654) & 0.0243 (0.0129, 0.0395) & 0.580 & 0.050 \\
2015 & P2 & 1413 & 401 & 0.0566 & -0.717 (-0.783, -0.651) & 0.0297 (0.0157, 0.0478) & 0.581 & 0.050 \\
\bottomrule
\end{tabular}
}
\end{table}
}{}

\begin{table}[t]\centering\scriptsize
\caption{Ablations on the headline archive, one setting changed per row, same origins and rows; pooled scores of the precedent rung. P1 scores minus 0.881 in every row except \emph{prior-normal} (minus 0.957) and the three $\kappa$ rows (minus 0.882, minus 0.880, minus 0.879). The \emph{effect-scale-relative} row is the cautionary run of Section~\ref{sec:cautionary} and is on a different effect scale.}
\label{tab:ablations-upworthy}
\resizebox{\columnwidth}{!}{% generated by paper/build_tables.py from ablations/backtest-upworthy-*.json (33 files; bootstrap resamples [500]; guard setting from each header p2_guard)
\begin{tabular}{lrrrr}
\toprule
setting & log score, guard on & Brier skill, guard on & log score, guard off & Brier skill, guard off \\
\midrule
admission-prelaunch & -0.831 & 0.002 (-0.001, 0.006) &  &  \\
bandwidth-mult-0.125 & -0.857 & 0.002 (-0.001, 0.005) & -0.858 & 0.002 (-0.001, 0.005) \\
bandwidth-mult-0.25 & -0.856 & 0.002 (-0.002, 0.005) & -0.857 & 0.001 (-0.003, 0.005) \\
bandwidth-mult-0.5 & -0.857 & 0.001 (-0.003, 0.004) & -0.858 & -0.000 (-0.004, 0.004) \\
bandwidth-mult-1.0 & -0.858 & 0.001 (-0.003, 0.004) & -0.859 & -0.000 (-0.004, 0.004) \\
contrasts-one & -0.870 & -0.002 (-0.007, 0.004) &  &  \\
effect-scale-relative+s-star-realized & -4.965 & -0.107 (-0.114, -0.101) &  &  \\
features-structured & -0.856 & 0.002 (-0.002, 0.005) & -0.857 & 0.001 (-0.003, 0.005) \\
features-text & -0.857 & 0.001 (-0.002, 0.005) & -0.860 & -0.001 (-0.005, 0.003) \\
kappa-0.0 & -0.859 & 0.000 (-0.003, 0.004) &  &  \\
kappa-5.0 & -0.858 & 0.001 (-0.003, 0.004) &  &  \\
kappa-80.0 & -0.855 & 0.002 (-0.002, 0.006) &  &  \\
lambda-0.0 & -0.857 & 0.001 (-0.002, 0.005) &  &  \\
lambda-0.2 & -0.858 & 0.002 (-0.002, 0.005) &  &  \\
lambda-0.5 & -0.861 & 0.000 (-0.003, 0.004) &  &  \\
lambda-1.0 & -0.869 & -0.002 (-0.007, 0.002) &  &  \\
p2-guard-1.0 & -0.857 & 0.001 (-0.003, 0.005) &  &  \\
p2-guard-off (main settings) &  &  & -0.857 & 0.001 (-0.003, 0.005) \\
p2-selection-loto & -0.861 & 0.002 (-0.002, 0.005) & -0.859 & 0.003 (-0.001, 0.006) \\
p2-selection-loto+recency-select & -0.861 & 0.002 (-0.002, 0.005) & -0.860 & 0.004 (0.001, 0.007) \\
pairwise-off & -0.858 & 0.001 (-0.002, 0.004) & -0.858 & 0.001 (-0.003, 0.005) \\
prior-normal & -0.957 & -0.014 (-0.019, -0.008) &  &  \\
recency-select & -0.856 & 0.002 (-0.002, 0.005) & -0.858 & 0.002 (-0.002, 0.005) \\
\bottomrule
\end{tabular}
}
\end{table}

\begin{table}[t]\centering\small
\caption{Gate sizes from the power simulation (synthetic, calibrated forecasts with the stated Brier skill against a base rate of 0.3).}
\label{tab:gates}
\resizebox{\columnwidth}{!}{% generated by paper/build_tables.py from /Volumes/Crucial X10/mootcourt-data/reports/gates.json (base rates ['0.09', '0.14', '0.30'], interval 0.9 by test, 3.5 rows per test, reps 1000, bootstrap 500, seed 20261102)
\begin{tabular}{rlrr}
\toprule
base rate & power target & tests for skill 0.20 & tests for skill 0.10 \\
\midrule
0.09 & 0.80 & 75 & 125 \\
0.09 & 0.90 & 100 & 175 \\
0.14 & 0.80 & 50 & 100 \\
0.14 & 0.90 & 75 & 125 \\
0.30 & 0.80 & 50 & 75 \\
0.30 & 0.90 & 50 & 100 \\
\bottomrule
\end{tabular}
}
\end{table}

% src: $REPORTS/run-log.json stages[].seconds, total_seconds 764.7
\begin{table}[t]\centering\scriptsize
\caption{Stages of the timed run (\texttt{reports/run-log.json}), seconds; all exit codes zero.}
\label{tab:walltime}
\begin{tabular}{lr}
\toprule
stage & seconds \\
\midrule
setup (environment from the uv cache) & 1.1 \\
data (checksum verification, no download) & 58.3 \\
power simulation (gates) & 47.6 \\
fit, headline archive & 29.8 \\
backtest, headline archive & 192.5 \\
fit, deployment-shaped archive & 1.8 \\
backtest, deployment-shaped archive & 7.4 \\
fit, synthetic & 3.0 \\
backtest, synthetic & 21.4 \\
gradient-boosted challenger & 374.6 \\
leakage audits (three archives) & 4.4 \\
data inventory & 3.8 \\
data verify & 9.9 \\
card latency & 9.1 \\
\midrule
total & 764.7 \\
\bottomrule
\end{tabular}
\end{table}

\begin{table}[t]\centering\scriptsize
\caption{Calibration by decile of planned exposure per arm (synthetic archive).}
\label{tab:calibration-synthetic}
\resizebox{\columnwidth}{!}{% generated by paper/build_tables.py from backtest-synthetic-20261103-bf4fd286e2b2.json
\begin{tabular}{rrrrrrr}
\toprule
decile & \multicolumn{3}{c}{P1: mean $\psig$ / event rate / cov.\ 90} & \multicolumn{3}{c}{P2: mean $\psig$ / event rate / cov.\ 90} \\
\midrule
1 & 0.181 & 0.220 & 0.90 & 0.181 & 0.220 & 0.90 \\
2 & 0.175 & 0.199 & 0.90 & 0.175 & 0.199 & 0.90 \\
3 & 0.181 & 0.210 & 0.91 & 0.181 & 0.210 & 0.91 \\
4 & 0.173 & 0.198 & 0.87 & 0.173 & 0.198 & 0.87 \\
5 & 0.200 & 0.149 & 0.91 & 0.200 & 0.149 & 0.91 \\
6 & 0.173 & 0.204 & 0.90 & 0.173 & 0.204 & 0.90 \\
7 & 0.194 & 0.214 & 0.92 & 0.194 & 0.214 & 0.92 \\
8 & 0.177 & 0.215 & 0.89 & 0.177 & 0.215 & 0.89 \\
9 & 0.188 & 0.160 & 0.87 & 0.188 & 0.160 & 0.87 \\
10 & 0.177 & 0.220 & 0.84 & 0.177 & 0.220 & 0.84 \\
\bottomrule
\end{tabular}
}
\end{table}

\begin{table}[t]\centering\small
\caption{Lift captured by prioritization policies at $K$ equal to a quarter of each period's tests (synthetic archive).}
\label{tab:decision-synthetic}
\resizebox{\columnwidth}{!}{% generated by paper/build_tables.py from /Volumes/Crucial X10/mootcourt-data/reports/backtest-synthetic-20261103-bf4fd286e2b2.json
\begin{tabular}{llrr}
\toprule
policy & model & captured fraction (mean over origins) & regret (mean) \\
\midrule
random & all & 0.249 & 172.1 \\
p\_win & P0 & 0.446 & 78.6 \\
p\_win & P1 & 0.446 & 78.6 \\
p\_win & P2 & 0.446 & 78.6 \\
oracle & all & 0.611 & 0.0 \\
\bottomrule
\end{tabular}
}
\end{table}

\begin{table}[t]\centering\small
\caption{Abstention triggers (synthetic archive).}
\label{tab:abstention-synthetic}
\resizebox{\columnwidth}{!}{% generated by paper/build_tables.py from backtest-synthetic-20261103-bf4fd286e2b2.json
\begin{tabular}{lrrrrrr}
\toprule
cell & n & rate & support & novelty & reliability & instability \\
\midrule
ALL & 1814 & 0.182 & 0 & 149 & 216 & 0 \\
both & 650 & 0.183 & 0 & 61 & 72 & 0 \\
headline\_only & 592 & 0.181 & 0 & 38 & 74 & 0 \\
image\_only & 572 & 0.182 & 0 & 50 & 70 & 0 \\
\bottomrule
\end{tabular}
}
\end{table}

\begin{table}[t]\centering\small
\caption{Precedent-weight selection by cell, averaged over origins (headline archive). \emph{loto} is leave-one-test-out.}
\label{tab:selection-upworthy}
\resizebox{\columnwidth}{!}{% generated by paper/build_tables.py from backtest-upworthy-20261103-4bfb354b87e3.json
\begin{tabular}{llrrrr}
\toprule
cell & selection & mean $\lambda$ & share $\lambda=0$ & mean bandwidth mult. & origins \\
\midrule
both & forward & 0.00 & 1.00 & 0.26 & 9 \\
headline\_only & forward & 0.47 & 0.22 & 0.32 & 9 \\
image\_only & forward & 0.04 & 0.89 & 0.51 & 9 \\
\bottomrule
\end{tabular}
}
\end{table}

\begin{table}[t]\centering\small
\caption{Precedent-weight selection by cell, averaged over origins (deployment-shaped archive).}
\label{tab:selection-asos}
\resizebox{\columnwidth}{!}{% generated by paper/build_tables.py from /Volumes/Crucial X10/mootcourt-data/reports/backtest-asos-20261103-6bb5a15c5ff0.json
\begin{tabular}{llrrrr}
\toprule
cell & selection & mean $\lambda$ & share $\lambda=0$ & mean bandwidth mult. & origins \\
\midrule
metric\_1 & loto & 0.00 & 1.00 & 0.75 & 7 \\
metric\_2 & loto & 0.00 & 1.00 & 0.66 & 7 \\
metric\_3 & loto & 0.00 & 1.00 & 0.66 & 7 \\
metric\_4 & loto & 0.00 & 1.00 & 0.55 & 7 \\
\bottomrule
\end{tabular}
}
\end{table}

\begin{table}[t]\centering\small
\caption{Ablations on the deployment-shaped archive, one setting changed per row; pooled scores of the precedent rung. P1 scores 4.610 in every row except \emph{prior-normal} (4.332) and the $\kappa$ rows (4.589, 4.602, 4.617).}
\label{tab:ablations-asos}
\resizebox{\columnwidth}{!}{% generated by paper/build_tables.py from ablations/backtest-asos-*.json (6 files; bootstrap resamples [500]; guard setting from each header p2_guard; P2 minus P1 is the paired logscore_vs_p1 of each run, nats per row, interval by test)
\begin{tabular}{lrrrrrr}
\toprule
setting & log score, guard on & Brier skill, guard on & P2 minus P1 log score, guard on (90\%) & log score, guard off & Brier skill, guard off & P2 minus P1 log score, guard off (90\%) \\
\midrule
kappa-0.0 & 4.693 & -0.119 (-0.251, -0.005) & 0.0000 (0.0000, 0.0000) &  &  &  \\
kappa-5.0 & 4.701 & -0.110 (-0.240, -0.003) & 0.0000 (0.0000, 0.0000) &  &  &  \\
kappa-80.0 & 4.717 & -0.107 (-0.232, -0.005) & 0.0007 (-0.0001, 0.0017) &  &  &  \\
p2-guard-1.0 & 4.712 & -0.109 (-0.239, -0.002) & 0.0029 (-0.0022, 0.0090) &  &  &  \\
p2-guard-off (main settings) &  &  &  & 4.717 & -0.109 (-0.245, -0.001) & 0.0083 (-0.0029, 0.0233) \\
prior-normal & 4.324 & -0.375 (-0.719, -0.134) & -0.0078 (-0.0273, 0.0087) &  &  &  \\
\bottomrule
\end{tabular}
}
\end{table}

\begin{table}[t]\centering\small
\caption{Abstention triggers (deployment-shaped archive).}
\label{tab:abstention-asos}
\resizebox{\columnwidth}{!}{% generated by paper/build_tables.py from backtest-asos-20261103-6bb5a15c5ff0.json
\begin{tabular}{lrrrrrr}
\toprule
cell & n & rate & support & novelty & reliability & instability \\
\midrule
ALL & 337 & 0.457 & 126 & 40 & 45 & 0 \\
metric\_1 & 85 & 0.365 & 24 & 7 & 12 & 0 \\
metric\_2 & 84 & 0.488 & 34 & 11 & 11 & 0 \\
metric\_3 & 84 & 0.488 & 34 & 11 & 11 & 0 \\
metric\_4 & 84 & 0.488 & 34 & 11 & 11 & 0 \\
\bottomrule
\end{tabular}
}
\end{table}

\begin{table}[t]\centering\scriptsize
\caption{Calibration by decile of planned exposure per arm (deployment-shaped archive; about 34 rows per decile).}
\label{tab:calibration-asos}
\resizebox{\columnwidth}{!}{% generated by paper/build_tables.py from backtest-asos-20261103-6bb5a15c5ff0.json
\begin{tabular}{rrrrrrr}
\toprule
decile & \multicolumn{3}{c}{P1: mean $\psig$ / event rate / cov.\ 90} & \multicolumn{3}{c}{P2: mean $\psig$ / event rate / cov.\ 90} \\
\midrule
1 & 0.037 & 0.000 & 0.94 & 0.037 & 0.000 & 0.94 \\
2 & 0.110 & 0.167 & 0.69 & 0.110 & 0.167 & 0.69 \\
3 & 0.082 & 0.156 & 0.84 & 0.082 & 0.156 & 0.84 \\
4 & 0.142 & 0.000 & 1.00 & 0.142 & 0.000 & 1.00 \\
5 & 0.180 & 0.250 & 0.66 & 0.180 & 0.250 & 0.66 \\
6 & 0.110 & 0.250 & 0.86 & 0.110 & 0.250 & 0.86 \\
7 & 0.209 & 0.250 & 0.97 & 0.209 & 0.250 & 0.97 \\
8 & 0.264 & 0.139 & 0.78 & 0.264 & 0.139 & 0.78 \\
9 & 0.329 & 0.031 & 0.88 & 0.329 & 0.031 & 0.88 \\
10 & 0.312 & 0.139 & 1.00 & 0.312 & 0.139 & 1.00 \\
\bottomrule
\end{tabular}
}
\end{table}

\begin{table}[t]\centering\small
\caption{Lift captured by prioritization policies at $K$ equal to a quarter of each period's tests (deployment-shaped archive; the primary metric only).}
\label{tab:decision-asos}
\resizebox{\columnwidth}{!}{% generated by paper/build_tables.py from /Volumes/Crucial X10/mootcourt-data/reports/backtest-asos-20261103-6bb5a15c5ff0.json (tables.decision_value; each mean is over the 7 per-origin rows of the policy and excludes the pooled origin=ALL row)
\begin{tabular}{llrr}
\toprule
policy & model & captured fraction (mean over origins) & regret (mean) \\
\midrule
random & all & 0.278 & 0.009 \\
p\_win & GBM & 0.290 & 0.008 \\
p\_win & P0 & 0.278 & 0.009 \\
p\_win & P1 & 0.278 & 0.009 \\
p\_win & P2 & 0.278 & 0.009 \\
p\_win & local\_eb & 0.278 & 0.009 \\
oracle & all & 0.857 & 0.000 \\
\bottomrule
\end{tabular}
}
\end{table}

\IfFileExists{tables/drift-asos.tex}{%
\begin{table}[t]\centering\scriptsize
\caption{Skill by nominal year (deployment-shaped archive; the calendar is nominal).}
\label{tab:drift-asos}
\resizebox{\columnwidth}{!}{% generated by paper/build_tables.py from /Volumes/Crucial X10/mootcourt-data/reports/backtest-asos-20261103-6bb5a15c5ff0.json (tables.score_by_year)
\begin{tabular}{rlrrrrrr}
\toprule
year & model & rows & tests & log score (90\%) & Brier skill (90\%) & dir.\ acc. & abst. \\
\midrule
2019 & P0 & 137 & 29 & 4.486 (4.281, 4.658) & -0.857 (-1.801, -0.247) & 0.380 & 0.920 \\
2019 & P1 & 137 & 29 & 4.901 (4.599, 5.161) & -0.108 (-0.404, 0.154) & 0.569 & 0.920 \\
2019 & P2 & 137 & 29 & 4.901 (4.599, 5.161) & -0.108 (-0.404, 0.154) & 0.569 & 0.920 \\
2020 & P0 & 200 & 39 & 4.313 (3.932, 4.670) & -0.201 (-0.552, 0.018) & 0.580 & 0.125 \\
2020 & P1 & 200 & 39 & 4.410 (3.904, 4.898) & -0.107 (-0.177, -0.034) & 0.490 & 0.125 \\
2020 & P2 & 200 & 39 & 4.410 (3.904, 4.898) & -0.107 (-0.177, -0.034) & 0.490 & 0.125 \\
\bottomrule
\end{tabular}
}
\end{table}
}{}

\section{Synthetic Truth}
\label{app:synthetic}
% src: src/mootcourt/sources/synthetic.py _DEFAULT_CELLS headline_only (2.0, 3.0), image_only (-1.0, 1.5), both (0.5, 2.0); text_effect_strength 0.0; $REPORTS/backtest-synthetic-20261103-bf4fd286e2b2.json header.admitted_rows 2000, header.scored_rows 1814, header.origins_scored 8, tables.aggregate[cell=ALL] n_tests 907
The synthetic archive has 2,000 contrasts in 1,000 two-arm tests from three cells with known means and spreads (2.0 and 3.0, minus 1.0 and 1.5, 0.5 and 2.0 percent lift) and no text signal; 1,814 contrasts from 907 tests are scored over eight origins.
% src: $REPORTS/backtest-synthetic-20261103-bf4fd286e2b2.json tables.split_half cell=headline_only mu_odd 2.02 mu_even 1.86 tau_odd 2.97 tau_even 3.10; cell=image_only mu -0.86 and -0.97, tau 1.25 and 1.34; cell=both mu 0.43 and 0.46, tau 1.74 and 1.81
The pooled prior recovers the generator: fit on odd and even tests separately, the cell means come back as 2.02 and 1.86, minus 0.86 and minus 0.97, and 0.43 and 0.46, and the spreads as 2.97 and 3.10, 1.25 and 1.34, and 1.74 and 1.81, each within a few tenths of its true value.
% src: $REPORTS/backtest-synthetic-20261103-bf4fd286e2b2.json tables.aggregate[cell=ALL,rung=P1] coverage_80 0.7834, coverage_90 0.8908, ece 0.0265, ece_percentile 97.8, brier_skill 0.0023 (-0.0128, 0.0168), decidable_by_tests "decidable at skill 0.20 and 0.10"; tables.p2_selection_by_cell share_lambda_zero 1.00 in all cells
Its intervals cover 78 and 89 percent, its expected calibration error of 0.026 sits at the 97.8th percentile of the null, its Brier skill is 0.002 (minus 0.013 to 0.017) and decidable at both targets, and P2 equals P1 at every origin because the guard correctly finds no text signal to admit.
% src: $REPORTS/backtest-synthetic-20261103-bf4fd286e2b2.json tables.decision_value origin=ALL captured_fraction random 0.2490, p_win (P0, P1, P2) 0.4453, oracle 0.6104
Here, unlike on the headline archive, prioritizing by $\pwin$ has decision value, 0.445 of the positive lift against 0.249 for random order and 0.610 for the oracle, because the three cells differ in their means and the cell is known before launch: the measure responds to between-cell differences the ledger can see, not to within-cell differences it cannot (tables in Appendix~\ref{app:tables}).

\section{Tooling Landscape}
\label{app:landscape}
A fetch-verified survey of thirteen commercial experimentation platforms, seventeen open-source projects, sixteen internal platforms described in public papers and the 2023 to 2026 wave of simulation and language-model tools (released as \texttt{docs/landscape.md}) found no product that issues a pre-launch outcome distribution or win probability for a proposed experiment, abstains, scores its forecasts against live results, or backtests over an archive; what exists are the parts (power calculators fed by historical variance, hand-typed priors, experiment knowledge bases, one platform-wide learned prior, two archive-trained ranking scores, and simulation agents validated only by directional backtests on a few dozen tests). Table~\ref{tab:landscape} (Appendix~\ref{app:tables}) compresses the survey's gap table.

\section{Weekly Replay Detail}
\label{app:replay}
% src: $REPORTS/replay-upworthy-20261103-4bfb354b87e3.json drift_onsets (both 2013-05-06, headline_only 2013-07-29, both and image_only 2013-08-12, both 2013-11-11, both 2014-02-24, both 2015-01-19, both 2015-04-13), flagged_weeks (49), weeks[].drifted_cells, weeks[].coverage_90 (minimum 0.4955 at 2013-09-30)
In the weekly replay of the headline archive (Table~\ref{tab:replay}) the cell that changed both headline and image, the smallest, is marked drifted in May and June 2013, in late 2013 and early 2014 (24 February to 7 April 2014), and in the first two quarters of 2015 (one week in January and the weeks from 13 April to the end of the archive); the headline-only cell is marked from 29 July 2013 and the image-only cell from 12 August 2013; between 12 August and 18 November 2013 all three cells are marked in most weeks; the image-only cell stays marked until 10 February 2014; and the rolling 90 percent coverage over the trailing eight weeks falls from 0.84 in July 2013 to 0.50 at the end of September 2013 before recovering to 0.95 by December. From April 2014 to the end of the archive no cell other than the smallest is ever marked.

\section{Wall Time, Latency and Scale}
\label{app:walltime}
% src: $REPORTS/run-log.json total_seconds 764.7, main_pipeline_seconds 362.9, stages[challenger-gbm].seconds 374.6, download_seconds 4.6, data_stage_kind "checksum verification", downloaded_files 0
Table~\ref{tab:walltime} (Appendix~\ref{app:tables}) gives the stages of the timed run: 765 seconds in all, of which the main pipeline (power simulation, three fits, three backtests) took 363 and the gradient-boosted challenger 375; checksum verification took 4.6 seconds and nothing was downloaded. Ablations and learning curves are not in the timed run.
% src: $REPORTS/card-latency.json impressions 20000, in_process_ms p50 16.8 p95 17.4 max 19.8 n 30, statuses issue 19 abstain 11, service_ms p50 19.0 max 46.9, cli_fresh_process_ms p50 1124.1, bundle_load_ms 72.6
A card for a headline test at 20,000 impressions per arm takes 16.8 milliseconds at the median (19.8 at most) in process with the bundle preloaded (30 cards, 19 issued, 11 abstained), 19.0 (46.9) through the loopback service, and 1.12 seconds in a fresh process, 73 milliseconds of which load the bundle.
% src: /Volumes/Crucial X10/mootcourt-data/archive-pre-f/reports/scale-100000.json stages.fit.seconds 93.45, wall_seconds 331.17, card_latency_ms p50 23.652, peak_rss_gb 1.48; scale-1000000.json stages.fit.seconds 1209.18, wall_seconds 3082.21, card_latency_ms p50 35.146, peak_rss_gb 6.54
On synthetic ledgers of one hundred thousand and one million contrasts the fit took 93 and 1,209 seconds, the median card 23.7 and 35.1 milliseconds, and peak memory was 1.48 and 6.54 GiB.
\clearpage


\end{document}
